% Sensibilisation sur la nécessité d'adaptation de l'Homme avec son environnement — SVT Tle D — Cours signé OnBuch+
% Programme officiel MINESEC. Compiler avec : tectonic N16-evolution-adaptation-homme.tex
\newcommand{\DOCMATIERE}{SVT}
\newcommand{\DOCNIVEAU}{Tle D}
\newcommand{\DOCMODULE}{Module III : L'Éducation à l'Environnement et au Développement Durable — Évolution de l'Homme}
\newcommand{\DOCLECON}{N16}
\newcommand{\DOCTITRE}{Sensibilisation sur la nécessité d'adaptation de l'Homme avec son environnement}
% OnBuch+ — préambule « cours signés » ENRICHI (compilé avec tectonic / XeLaTeX)
% Système visuel « Pop » d'OnBuch+ : fond crème (jamais blanc), bordures encre
% épaisses, ombres « dures » décalées, titres Archivo Black en capitales.
% Version MATHÉMATIQUES : graphes (pgfplots/tikz), boîtes pédagogiques
% (définition, propriété, méthode, attention, le savais-tu, exercice résolu,
% exercice, TP, bilan), page de garde et signature OnBuch+ ancrée en bas.
%
% Macros attendues dans main.tex AVANT \input{preamble} :
%   \DOCMATIERE  \DOCNIVEAU  \DOCMODULE  \DOCLECON (numéro)  \DOCTITRE
\documentclass[11pt]{article}

\usepackage{fontspec}
\usepackage[french]{babel}
\usepackage[a4paper,margin=2cm,top=3.4cm,bottom=2.8cm,headheight=1.4cm,footskip=1.3cm]{geometry}
\usepackage{amsmath,amssymb,amsfonts}
\usepackage{mathtools} % \xmapsto, \coloneqq... souvent utilisés par les modèles
\usepackage{cancel} % simplifications barrées (bilans de forces)
% \abs{z} (module, valeur absolue) et \norm{u} (norme), utilisés par les modèles
\providecommand{\abs}{}\let\abs\relax\DeclarePairedDelimiter{\abs}{\lvert}{\rvert}
\providecommand{\norm}{}\let\norm\relax\DeclarePairedDelimiter{\norm}{\lVert}{\rVert}
\usepackage[table]{xcolor} % [table] : fournit aussi \rowcolors (alternance de lignes)
\usepackage{pagecolor}
\usepackage{tikz}
\usetikzlibrary{arrows.meta,positioning,calc,shapes.geometric,shapes.misc,shapes.arrows,decorations.pathmorphing,decorations.pathreplacing,decorations.markings,patterns,shadows,mindmap,trees,fit,backgrounds,matrix,intersections,angles,quotes}
\usepackage{pgfplots}
\usepackage[version=4]{mhchem} % équations nucléaires \ce{^{235}_{92}U}
\usepackage[european,siunitx]{circuitikz} % schémas électriques
\pgfplotsset{compat=1.18}
\usepgfplotslibrary{fillbetween,groupplots} % aires entre courbes (calcul intégral)
\usepackage{enumitem}
\usepackage{fancyhdr}
% [most] charge toutes les bibliothèques tcolorbox (listings, minted, raster,
% documentation...) et gonfle inutilement la mémoire TeX sur un document long
% (~300 boîtes) : on ne charge que ce qui est réellement utilisé.
\usepackage[skins,breakable]{tcolorbox}
\usepackage{graphicx}
\usepackage{colortbl}
\usepackage{booktabs}
\usepackage{tabularx}
\usepackage{diagbox} % case d'en-tête diagonale des tableaux à double entrée
% Colonnes extensibles centrée / gauche / droite (C, L, R), souvent utilisées par les modèles
\newcolumntype{C}{>{\centering\arraybackslash}X}
\newcolumntype{L}{>{\raggedright\arraybackslash}X}
\newcolumntype{R}{>{\raggedleft\arraybackslash}X}
\usepackage{array}
\usepackage{multicol}
\usepackage{siunitx}
% parse-numbers=false : accepte aussi \SI{2,5 \times 10^{-3}}{...}
\sisetup{locale=FR,output-decimal-marker={,},group-minimum-digits=4,parse-numbers=false}
\DeclareSIUnit\ohm{\ensuremath{\symup{\Omega}}} % U+2126 absent de la police
\DeclareSIUnit\division{div} % réglages d'oscilloscope (ms/div, V/div)
\DeclareSIUnit\tr{tr} % tours (vitesse de rotation, ex. \SI{1500}{\tr\per\minute})
\DeclareSIUnit\molar{M} % concentration molaire (ex. \SI{3}{\molar}, \SI{1}{\micro\molar})
\DeclareSIUnit\an{an} % année (le modèle confond parfois avec \year, un compteur TeX) — ex. \SI{5}{\kilo\meter\per\an}
\DeclareSIUnit\FCFA{FCFA} % franc CFA (ex. \SI{500}{\FCFA\per\kilo\gram})
\DeclareSIUnit\degreeBrix{\textdegree Brix} % teneur en sucre d'un sirop/jus (réfractométrie)
\DeclareSIUnit\month{mois} % le modèle utilise parfois \month, un compteur TeX, comme unité de durée
\usepackage{qrcode}
% \degree : commande fréquemment utilisée par les modèles pour un angle ou une
% température en dehors de \SI{}{} (non fournie par les paquets chargés).
\providecommand{\degree}{^{\circ}}
% \qed (fin de démonstration), utilisé par les modèles sans amsthm
\providecommand{\qed}{\hfill$\square$}
% \barre : double barre des valeurs interdites dans un tableau de variations (tkz-tab)
\providecommand{\barre}{\ensuremath{\|}}
% environnement proof (amsthm non chargé) utilisé par les modèles
\ifdefined\proof\else\newenvironment{proof}[1][Démonstration]{\par\noindent\textit{#1.}\ }{\qed\par}\fi
\usepackage{lastpage}
\usepackage{hyperref}
\hypersetup{hidelinks}

% ============================================================
% Palette « Pop » OnBuch+ — mêmes hex que la classe P (plus_ui.dart)
% ============================================================
\definecolor{poporange}{HTML}{F59321}
\definecolor{popdark}{HTML}{E07A0C}
\definecolor{popcream}{HTML}{FFF6E8}
\definecolor{popink}{HTML}{1C1714}
\definecolor{poppurple}{HTML}{7A5AE0}
\definecolor{popgreen}{HTML}{2E9E6A}
\definecolor{poppink}{HTML}{E0457B}
\definecolor{popblue}{HTML}{2F7DE1}
\definecolor{popgold}{HTML}{C98A12}
\definecolor{popmuted}{HTML}{8A7F76}
\definecolor{popline}{HTML}{EADFD2}
\definecolor{popgray}{HTML}{8A7F76}
\colorlet{popgrey}{popgray}
% Alias tolérants (le modèle nomme parfois les couleurs différemment)
\colorlet{popwhite}{white}
\colorlet{popblack}{popink}
\colorlet{popred}{poppink}
\colorlet{popyellow}{popgold}
% Teintes claires (fonds de tableaux/encadrés) — le modèle nomme parfois
% les couleurs avec un suffixe L (ex. popblueL) sans les avoir définies.
\colorlet{poporangeL}{poporange!20}
\colorlet{popdarkL}{popdark!20}
\colorlet{poppurpleL}{poppurple!20}
\colorlet{popgreenL}{popgreen!20}
\colorlet{poppinkL}{poppink!20}
\colorlet{popblueL}{popblue!20}
\colorlet{popgoldL}{popgold!20}
\colorlet{popredL}{popred!20}
\colorlet{popgrayL}{popgray!20}
% Teintes claires pour fonds de tableaux / zones
\colorlet{popblueL}{popblue!10!white}
\colorlet{poporangeL}{poporange!14!white}
\colorlet{popgreenL}{popgreen!12!white}
\colorlet{poppurpleL}{poppurple!12!white}
\colorlet{poppinkL}{poppink!12!white}
\colorlet{popgoldL}{popgold!14!white}

\pagecolor{popcream}

% ============================================================
% Typographie
% ============================================================
\defaultfontfeatures{Ligatures=TeX}
\setmainfont{PlusJakartaSans-Medium.ttf}[Path=../../../fonts/,BoldFont=PlusJakartaSans-Bold.ttf]
% Maths en sans-serif (Fira Math) avec chiffres et lettres droites de Plus
% Jakarta, pour que les formules se fondent dans le texte.
\usepackage[math-style=ISO,bold-style=ISO]{unicode-math}
\setmathfont{FiraMath-Regular.otf}[Path=../../../fonts/]
\setmathfont{PlusJakartaSans-Medium.ttf}[Path=../../../fonts/,range={up/num,up/{Latin,latin},\mathup}]
\usepackage{icomma} % virgule décimale sans espace en mode math
% Caractères mathématiques Unicode absents des polices de texte (Plus Jakarta,
% Archivo Black) : rendus par leur équivalent mathématique, en texte comme en
% formule — sinon ils s'affichent en carré vide (ex. « Arithmétique dans ℤ »).
\usepackage{newunicodechar}
\newunicodechar{ℝ}{\ensuremath{\mathbb{R}}}
\newunicodechar{ℤ}{\ensuremath{\mathbb{Z}}}
\newunicodechar{ℂ}{\ensuremath{\mathbb{C}}}
\newunicodechar{ℕ}{\ensuremath{\mathbb{N}}}
\newunicodechar{ℚ}{\ensuremath{\mathbb{Q}}}
\newunicodechar{𝔻}{\ensuremath{\mathbb{D}}}
\newunicodechar{α}{\ensuremath{\alpha}}
\newunicodechar{β}{\ensuremath{\beta}}
\newunicodechar{γ}{\ensuremath{\gamma}}
\newunicodechar{δ}{\ensuremath{\delta}}
\newunicodechar{ε}{\ensuremath{\varepsilon}}
\newunicodechar{θ}{\ensuremath{\theta}}
\newunicodechar{λ}{\ensuremath{\lambda}}
\newunicodechar{μ}{\ensuremath{\mu}}
\newunicodechar{σ}{\ensuremath{\sigma}}
\newunicodechar{τ}{\ensuremath{\tau}}
\newunicodechar{φ}{\ensuremath{\varphi}}
\newunicodechar{ψ}{\ensuremath{\psi}}
\newunicodechar{ω}{\ensuremath{\omega}}
\newunicodechar{Σ}{\ensuremath{\Sigma}}
% majuscules produites par \MakeUppercase dans les titres (α-aminés -> Α)
\newunicodechar{Α}{\ensuremath{\alpha}}
\newunicodechar{Β}{\ensuremath{\beta}}
\newunicodechar{Γ}{\ensuremath{\Gamma}}
\newunicodechar{Δ}{\ensuremath{\Delta}}
\newunicodechar{Ε}{\ensuremath{\varepsilon}}
\newunicodechar{Θ}{\ensuremath{\Theta}}
\newunicodechar{Λ}{\ensuremath{\Lambda}}
\newunicodechar{Μ}{\ensuremath{\mu}}
\newunicodechar{Τ}{\ensuremath{\tau}}
\newunicodechar{Φ}{\ensuremath{\Phi}}
\newunicodechar{Ψ}{\ensuremath{\Psi}}
\newunicodechar{Ω}{\ensuremath{\Omega}}
\newunicodechar{↦}{\ensuremath{\mapsto}}
\newunicodechar{∈}{\ensuremath{\in}}
\newunicodechar{∉}{\ensuremath{\notin}}
\newunicodechar{∧}{\ensuremath{\wedge}}
\newunicodechar{⇔}{\ensuremath{\Leftrightarrow}}
\newunicodechar{⟺}{\ensuremath{\Longleftrightarrow}}
\newunicodechar{⇒}{\ensuremath{\Rightarrow}}
\newunicodechar{⟹}{\ensuremath{\Longrightarrow}}
\newunicodechar{∘}{\ensuremath{\circ}}
\newunicodechar{∪}{\ensuremath{\cup}}
\newunicodechar{∩}{\ensuremath{\cap}}
\newunicodechar{≡}{\ensuremath{\equiv}}
\newunicodechar{⊂}{\ensuremath{\subset}}
\newunicodechar{⊥}{\ensuremath{\perp}}
\newunicodechar{∃}{\ensuremath{\exists}}
\newunicodechar{∀}{\ensuremath{\forall}}
\newunicodechar{∛}{\ensuremath{\sqrt[3]{\,}}}
\newunicodechar{∜}{\ensuremath{\sqrt[4]{\,}}}
\newunicodechar{⃗}{\ensuremath{{}^{\to}}}
\newunicodechar{ⁿ}{\ifmmode^{n}\else\textsuperscript{\ensuremath{n}}\fi}
\newunicodechar{ˣ}{\ifmmode^{x}\else\textsuperscript{\ensuremath{x}}\fi}
\newunicodechar{ᵇ}{\ifmmode^{b}\else\textsuperscript{\ensuremath{b}}\fi}
\newunicodechar{ʳ}{\ifmmode^{r}\else\textsuperscript{\ensuremath{r}}\fi}
\newunicodechar{ᵘ}{\ifmmode^{u}\else\textsuperscript{\ensuremath{u}}\fi}
\newunicodechar{ʸ}{\ifmmode^{y}\else\textsuperscript{\ensuremath{y}}\fi}
\newunicodechar{ᵃ}{\ifmmode^{a}\else\textsuperscript{\ensuremath{a}}\fi}
\newunicodechar{ᵉ}{\ifmmode^{e}\else\textsuperscript{\ensuremath{e}}\fi}
\newunicodechar{ᵏ}{\ifmmode^{k}\else\textsuperscript{\ensuremath{k}}\fi}
\newunicodechar{ᵖ}{\ifmmode^{p}\else\textsuperscript{\ensuremath{p}}\fi}
\newunicodechar{ᴺ}{\ifmmode^{N}\else\textsuperscript{\ensuremath{N}}\fi}
\newunicodechar{ᶜ}{\ifmmode^{c}\else\textsuperscript{\ensuremath{c}}\fi}
\newunicodechar{ʰ}{\ifmmode^{h}\else\textsuperscript{\ensuremath{h}}\fi}
\newunicodechar{ᵠ}{\ifmmode^{q}\else\textsuperscript{\ensuremath{q}}\fi}
\newunicodechar{ₙ}{\ifmmode_{n}\else\textsubscript{\ensuremath{n}}\fi}
\newunicodechar{ₐ}{\ifmmode_{a}\else\textsubscript{\ensuremath{a}}\fi}
\newunicodechar{ᵢ}{\ifmmode_{i}\else\textsubscript{\ensuremath{i}}\fi}
\newunicodechar{ᵣ}{\ifmmode_{r}\else\textsubscript{\ensuremath{r}}\fi}
\newunicodechar{ₖ}{\ifmmode_{k}\else\textsubscript{\ensuremath{k}}\fi}
\newunicodechar{ᵦ}{\ifmmode_{\beta}\else\textsubscript{\ensuremath{\beta}}\fi}
\newunicodechar{ₓ}{\ifmmode_{x}\else\textsubscript{\ensuremath{x}}\fi}
\newfontfamily\popdisplay{ArchivoBlack-Regular.ttf}[Path=../../../fonts/]
\newfontfamily\poplogo{SpaceGrotesk-Bold.ttf}[Path=../../../fonts/]
\newfontfamily\popsemi{PlusJakartaSans-SemiBold.ttf}[Path=../../../fonts/]
\newfontfamily\popxbold{PlusJakartaSans-ExtraBold.ttf}[Path=../../../fonts/]
\newfontfamily\popxboldls{PlusJakartaSans-ExtraBold.ttf}[Path=../../../fonts/,LetterSpace=3.0]

\setlist[enumerate]{leftmargin=1.7em,itemsep=5pt,topsep=4pt}
\setlist[itemize]{leftmargin=1.7em,itemsep=4pt,topsep=4pt,label={\color{poporange}\textbullet}}
\setlength{\parindent}{0pt}
\setlength{\parskip}{5pt}
\renewcommand{\arraystretch}{1.25}

% pgfplots : style Pop par défaut
\pgfplotsset{
  every axis/.append style={axis line style={popink,line width=1pt},tick style={popink},
    grid style={popline,line width=0.5pt},label style={font=\small},tick label style={font=\footnotesize}},
  popcurve/.style={poporange,line width=1.8pt,no markers,smooth},
}
% Même style disponible dans un \draw TikZ simple (hors pgfplots)
\tikzset{popcurve/.style={poporange,line width=1.8pt}}
\tikzset{popgrid/.style={popline,very thin}}
\colorlet{popcurve}{poporange} % le nom du style est parfois utilisé comme couleur

% ============================================================
% Pill / squiggle
% ============================================================
\newcommand{\poppill}[3]{%
  \tikz[baseline=-0.55ex]\node[draw=popink,line width=1.2pt,rounded corners=2.6pt,%
    fill=#2,inner xsep=6.5pt,inner ysep=3pt]{%
    \popxboldls\fontsize{7}{7}\selectfont\color{#3}\MakeUppercase{#1}};%
}
\newcommand{\popsquiggle}[1]{%
  \tikz[baseline=-0.4ex]{
    \draw[#1,line width=2.6pt,line cap=round]
      plot [smooth] coordinates {(0,0.09) (0.34,0.01) (0.68,0.21) (1.02,0.01) (1.36,0.10)};
  }%
}
% Mot-clé surligné (marqueur orange sous le texte)
\newcommand{\cle}[1]{{\popxbold\color{popdark}#1}}

% ============================================================
% En-tête / pied
% ============================================================
\pagestyle{fancy}
\fancyhf{}
\renewcommand{\headrulewidth}{0pt}
\renewcommand{\footrulewidth}{0pt}
\lhead{\raisebox{-0.15ex}{\poplogo\fontsize{13}{13}\selectfont\color{popink}OnBuch\color{poporange}+}}
\rhead{\poppill{\DOCMATIERE\ \textbullet\ \DOCNIVEAU\ \textbullet\ Leçon \DOCLECON}{white}{popink}}
\lfoot{\popxbold\fontsize{7.6}{7.6}\selectfont\color{popmuted}\MakeUppercase{Cours signé OnBuch+}\ \textbullet\ {\popsemi\DOCTITRE}}
\rfoot{\tikz[baseline=-0.6ex]\node[rounded corners=8.5pt,fill=popink,minimum height=17pt,minimum width=17pt,inner xsep=5pt,inner sep=0]{\popxbold\fontsize{8}{8}\selectfont\color{popcream}\thepage\,/\,\pageref*{LastPage}};}

% ============================================================
% Titres
% ============================================================
\newcommand{\chaptitre}[2]{%
  \begin{tcolorbox}[enhanced,colback=poporange,colframe=popink,boxrule=2.6pt,arc=11pt,
      drop shadow=popink,left=15pt,right=15pt,top=12pt,bottom=11pt,before skip=3pt,after skip=18pt]
    \poppill{#2}{popink}{popcream}\par\vspace{9pt}
    {\raggedright\hyphenpenalty=10000\exhyphenpenalty=10000\popdisplay\fontsize{22}{24}\selectfont\color{popcream}\MakeUppercase{#1}\par}
  \end{tcolorbox}
}
% Section numérotée : pastille orange avec le numéro + titre Archivo
\newcounter{popsec}
\newcommand{\coursec}[1]{%
  \stepcounter{popsec}\setcounter{popsub}{0}%
  \par\vspace{16pt}\noindent
  \begin{minipage}{\linewidth}
  \tikz[baseline=-0.7ex]\node[draw=popink,line width=1.6pt,fill=poporange,rounded corners=4pt,minimum size=22pt,inner sep=0]{\popdisplay\fontsize{12}{12}\selectfont\color{popcream}\thepopsec};%
  \hspace{8pt}{\popdisplay\fontsize{15}{17}\selectfont\color{popink}\MakeUppercase{#1}}\par
  \vspace{4pt}\noindent{\color{poporange}\rule{36pt}{3.6pt}}
  \end{minipage}\par\nopagebreak\vspace{8pt}
}
\newcounter{popsub}
\newcommand{\courssub}[1]{%
  \stepcounter{popsub}%
  \par\vspace{9pt}\noindent{\popxbold\fontsize{11.5}{13}\selectfont\color{popink}{\color{poporange}\thepopsec.\thepopsub}\hspace{6pt}#1}\par\nopagebreak\vspace{4pt}
}

% ============================================================
% Boîtes pédagogiques — même recette : fond blanc, bordure épaisse colorée,
% ombre dure encre, badge plein. Titre optionnel [#1].
% ============================================================
\tcbset{popbase/.style={breakable,enhanced,colback=white,boxrule=2.3pt,arc=9pt,
  shadow={3pt}{-3pt}{0pt}{popink},left=14pt,right=14pt,top=11pt,bottom=11pt,before skip=11pt,after skip=15pt,
  fontupper=\popsemi\fontsize{10.6}{14.5}\selectfont\color{popink},
  fontlower=\popsemi\fontsize{10.6}{14.5}\selectfont\color{popink}}}
% Coche dessinée (le glyphe ✓ n'existe pas dans les polices du document)
\newcommand{\popcheck}[1]{\tikz[baseline=-0.5ex]\draw[#1,line width=1.6pt,line cap=round,line join=round](-0.13,0.0)--(-0.04,-0.09)--(0.13,0.1);}
\providecommand{\coche}{\popcheck{popgreen}}
\providecommand{\resultat}[1]{\par\smallskip\noindent\fcolorbox{popgreen}{popgreenL}{\parbox{\dimexpr\linewidth-2\fboxsep-2\fboxrule\relax}{\textbf{Résultat.} #1}}\par\smallskip}
\newcommand{\popboxhead}[3]{\poppill{#1}{#2}{white}\ifx\relax#3\relax\else\hspace{6pt}{\popxbold\fontsize{10.6}{12}\selectfont\color{popink}#3}\fi\par\vspace{7pt}}

\newtcolorbox{aretenir}[1][]{popbase,colframe=popgreen,before upper={\popboxhead{À retenir}{popgreen}{#1}}}
\newtcolorbox{exemplebox}[1][]{popbase,colframe=poporange,before upper={\popboxhead{Exemple}{poporange}{#1}}}
\newtcolorbox{definition}[1][]{popbase,colframe=popblue,before upper={\popboxhead{Définition}{popblue}{#1}}}
\newtcolorbox{propriete}[1][]{popbase,colframe=poppurple,before upper={\popboxhead{Propriété}{poppurple}{#1}}}
\newtcolorbox{methode}[1][]{popbase,colframe=popgold,before upper={\popboxhead{Méthode}{popgold}{#1}}}
\newtcolorbox{attention}[1][]{popbase,colframe=poppink,before upper={\popboxhead{Attention}{poppink}{#1}}}
\newtcolorbox{experience}[1][]{popbase,colframe=popblue,colback=popblueL,before upper={\popboxhead{Expérience}{popblue}{#1}}}
% « Le savais-tu ? » : carte violette pleine (contexte, culture, Cameroun)
\newtcolorbox{savaistu}[1][]{popbase,colback=poppurple,colframe=popink,
  fontupper=\popsemi\fontsize{10.4}{14.2}\selectfont\color{white},
  before upper={\poppill{Le savais-tu\,?}{white}{poppurple}\ifx\relax#1\relax\else\hspace{6pt}{\popxbold\color{white}#1}\fi\par\vspace{7pt}}}
% Exercice résolu : énoncé en haut, \tcblower puis la solution sur fond crème
\newcounter{popexo}\newcounter{popexr}
\newtcolorbox{exoresolu}[1][]{popbase,colframe=poporange,colbacklower=popcream,
  segmentation style={popink,line width=1.2pt,dashed},
  before upper={\stepcounter{popexr}\popboxhead{Exercice résolu \thepopexr}{poporange}{#1}},
  before lower={\poppill{Solution}{popink}{popcream}\par\vspace{6pt}}}
% Exercice (non résolu en place) — la correction est dans la partie Corrigés
\newtcolorbox{exercice}[1][]{popbase,colframe=popink,
  before upper={\stepcounter{popexo}\popboxhead{Exercice \thepopexo}{popink}{#1}}}
% Corrigé
\newtcolorbox{corrige}[1][]{popbase,colframe=popgreen,colback=popgreenL,
  before upper={\popboxhead{Corrigé}{popgreen}{#1}}}
% Objectifs / prérequis / bilan
\newtcolorbox{objectifs}{popbase,colframe=popink,colback=poporangeL,
  before upper={\popboxhead{Objectifs de la leçon}{popink}{}}}
\newtcolorbox{prerequis}{popbase,colframe=popink,colback=popblueL,
  before upper={\popboxhead{Prérequis}{popblue}{}}}
\newtcolorbox{bilan}{popbase,colframe=popink,colback=white,boxrule=2.8pt,
  before upper={\popboxhead{Fiche bilan}{poporange}{L'essentiel à connaître pour l'examen}}}
% Figure encadrée avec légende
\newtcolorbox{popfigure}[1][]{enhanced,breakable=false,colback=white,colframe=popline,boxrule=1.4pt,arc=8pt,
  left=10pt,right=10pt,top=10pt,bottom=8pt,before skip=10pt,after skip=12pt,halign=center,
  lower separated=false,halign lower=center,
  fontlower=\popsemi\fontsize{9}{11.5}\selectfont\color{popmuted},
  #1}
\newcommand{\legende}[1]{\par\vspace{6pt}{\popsemi\fontsize{9}{11.5}\selectfont\color{popmuted}#1}}

% ============================================================
% Signature OnBuch+
% ============================================================
\newcommand{\poplogoinline}[1]{{\poplogo\fontsize{#1}{#1}\selectfont\color{popink}OnBuch\color{poporange}+}}
\newcommand{\signatureonbuch}{%
  \begin{tcolorbox}[enhanced,colback=popink,colframe=popink,boxrule=2.2pt,arc=10pt,
      drop shadow=poporange,left=14pt,right=14pt,top=10pt,bottom=10pt,before skip=0pt,after skip=0pt]
    \tikz[baseline=-0.7ex]\node[circle,fill=popgreen,draw=popcream,line width=1.3pt,minimum size=22pt,inner sep=0]{\popcheck{white}};%
    \hspace{9pt}\begin{minipage}[c]{0.62\linewidth}
      {\popxbold\fontsize{12}{13}\selectfont\color{popcream}Signé\ \textbullet\ {\poplogo OnBuch\color{poporange}+}}\par\vspace{2pt}
      {\raggedright\popsemi\fontsize{8}{10.5}\selectfont\color{popline}\MakeUppercase{Équipe pédagogique OnBuch+}\\\MakeUppercase{Conforme au programme officiel MINESEC}\par}
    \end{minipage}\hfill
    {\poplogo\fontsize{22}{22}\selectfont\color{popcream}OnBuch\color{poporange}+}
  \end{tcolorbox}%
}

% ============================================================
% Page de garde
% #1 titre · #2 matière · #3 niveau · #4 module · #5 n° leçon · #6 durée ·
% #7 description / objectifs (texte ou itemize)
% ============================================================
\newcommand{\pagegarde}[7]{%
  \thispagestyle{empty}
  \newgeometry{a4paper,margin=1.7cm,top=1.6cm,bottom=1.4cm}%
  \begin{tikzpicture}[remember picture,overlay]
    \fill[poporange!18] ([xshift=-3.2cm,yshift=-3.6cm]current page.north east) circle (4.6cm);
    \fill[poppurple!14] ([xshift=2.4cm,yshift=7.6cm]current page.south west) circle (3.8cm);
  \end{tikzpicture}%
  {\poplogo\fontsize{30}{30}\selectfont\color{popink}OnBuch\color{poporange}+}\hfill
  \raisebox{0.6ex}{\poppill{Cours signé \textbullet\ Premium}{popink}{popcream}}\par
  \vspace{1pt}\popsquiggle{poppurple}\par
  \vspace{8pt}
  \begin{tcolorbox}[enhanced,colback=poporange,colframe=popink,boxrule=2.8pt,arc=13pt,
      drop shadow=popink,left=18pt,right=18pt,top=12pt,bottom=13pt,after skip=10pt]
    \poppill{#2 \textbullet\ #3}{popink}{popcream}\hspace{5pt}\poppill{#4}{white}{popink}\par\vspace{8pt}
    {\popdisplay\fontsize{14}{14}\selectfont\color{popink}LEÇON #5}\par\vspace{4pt}
    {\raggedright\hyphenpenalty=10000\exhyphenpenalty=10000\popdisplay\fontsize{25}{27}\selectfont\color{popcream}\MakeUppercase{#1}\par}
    \vspace{8pt}
    {\popxbold\fontsize{9.5}{9.5}\selectfont\color{popink}\MakeUppercase{Durée conseillée : #6}}
  \end{tcolorbox}
  \begin{tcolorbox}[enhanced,colback=white,colframe=popink,boxrule=2.2pt,arc=10pt,
      drop shadow=popink,left=16pt,right=16pt,top=10pt,bottom=10pt,after skip=8pt]
    \poppill{Dans cette leçon}{poppurple}{white}\par\vspace{5pt}
    \popsemi\fontsize{9.3}{12.6}\selectfont\color{popink}#7
  \end{tcolorbox}
  \noindent\begin{minipage}[t]{0.487\linewidth}
    \begin{tcolorbox}[enhanced,colback=popgreen,colframe=popink,boxrule=2.2pt,arc=10pt,
        drop shadow=popink,left=11pt,right=11pt,top=10pt,bottom=10pt,height=3.1cm,valign=center]
      \tcbox[colback=white,colframe=popink,boxrule=1.3pt,arc=5pt,boxsep=0pt,nobeforeafter,
        left=3pt,right=3pt,top=3pt,bottom=3pt,valign=center]{\qrcode[height=1.9cm]{https://play.google.com/store/apps/details?id=com.luvvix.onbuch&hl=fr}}%
      \hspace{8pt}\begin{minipage}[c]{0.5\linewidth}
        {\popdisplay\fontsize{11}{12.5}\selectfont\color{white}\MakeUppercase{Télécharge OnBuch}}\par\vspace{4pt}
        {\raggedright\popsemi\fontsize{8.4}{11}\selectfont\color{white}Gratuit \textbullet\ Google Play\\Tuteur IA, annales, concours, résultats\par}
      \end{minipage}
    \end{tcolorbox}
  \end{minipage}\hfill
  \begin{minipage}[t]{0.487\linewidth}
    \begin{tcolorbox}[enhanced,colback=poppurple,colframe=popink,boxrule=2.2pt,arc=10pt,
        drop shadow=popink,left=11pt,right=11pt,top=10pt,bottom=10pt,height=3.1cm,valign=center]
      \tikz[baseline=-0.7ex]\node[circle,fill=white,draw=popink,line width=1.3pt,minimum size=1.6cm,inner sep=0]{\popdisplay\fontsize{24}{24}\selectfont\color{poppurple}+};%
      \hspace{8pt}\begin{minipage}[c]{0.6\linewidth}
        {\popdisplay\fontsize{11}{12.5}\selectfont\color{white}\MakeUppercase{Passe à OnBuch+}}\par\vspace{4pt}
        {\raggedright\popsemi\fontsize{8.4}{11}\selectfont\color{white}Tous les cours signés, Tuteur IA illimité, fiches PDF — onglet OnBuch+ de l'app\par}
      \end{minipage}
    \end{tcolorbox}
  \end{minipage}
  \vspace{8pt}
  \popcontenu
  \vfill
  \signatureonbuch
  \restoregeometry
  \newpage
}

% Rangée « Ce cours contient » (page de garde)
\newcommand{\poptile}[3]{% #1 couleur #2 gros texte #3 légende
  \begin{tcolorbox}[enhanced,colback=#1,colframe=popink,boxrule=2pt,arc=9pt,shadow={2.5pt}{-2.5pt}{0pt}{popink},
      left=6pt,right=6pt,top=9pt,bottom=9pt,halign=center,valign=center,height=1.75cm,width=\linewidth,nobeforeafter]
    {\popdisplay\fontsize{12.5}{13}\selectfont\color{white}\MakeUppercase{#2}}\par\vspace{3pt}
    {\centering\popsemi\fontsize{7.8}{9.5}\selectfont\color{white}#3\par}
  \end{tcolorbox}}
\newcommand{\popcontenu}{%
  \par\noindent{\popxboldls\fontsize{7.5}{7.5}\selectfont\color{popmuted}CE COURS SIGNÉ CONTIENT}\par\vspace{7pt}
  \noindent\begin{minipage}[t]{0.235\linewidth}\poptile{popblue}{Cours}{complet, illustré, pas à pas}\end{minipage}\hfill
  \begin{minipage}[t]{0.235\linewidth}\poptile{poporange}{Méthodes}{et exercices résolus}\end{minipage}\hfill
  \begin{minipage}[t]{0.235\linewidth}\poptile{poppink}{Exercices}{type examen + corrigés}\end{minipage}\hfill
  \begin{minipage}[t]{0.235\linewidth}\poptile{popgreen}{Fiche bilan}{l'essentiel à retenir}\end{minipage}\par}

% Sommaire
\newtcolorbox{sommaire}{enhanced,colback=white,colframe=popink,boxrule=2.2pt,arc=10pt,
  drop shadow=popink,left=18pt,right=18pt,top=13pt,bottom=13pt,before skip=6pt,after skip=20pt,
  before upper={\poppill{Sommaire}{poppurple}{white}\par\vspace{9pt}\popsemi\fontsize{10.6}{14.5}\selectfont\color{popink}}}

% Signature de fin de cours — poussée tout en bas de la dernière page
\newcommand{\findecours}{%
  \par\vfill\nopagebreak
  \begin{center}\popsquiggle{poporange}\end{center}\vspace{4pt}
  \signatureonbuch
}
\AtBeginDocument{\def\llbracket{\mathopen{[\mkern-3.5mu[}}\def\rrbracket{\mathclose{]\mkern-3.5mu]}}} % intervalles d'entiers (glyphe ⟦ absent de la police)
% Glyphes absents de Fira Math : redéfinis à partir de glyphes disponibles
\AtBeginDocument{%
  \def\setminus{\mathbin{\backslash}}\let\smallsetminus\setminus
  \def\vdots{\mathord{\vbox{\baselineskip3.2pt\lineskiplimit0pt\kern4pt\hbox{.}\hbox{.}\hbox{.}}}}%
  \def\ddots{\mathinner{\mkern1mu\raise6.5pt\hbox{.}\mkern2mu\raise3.6pt\hbox{.}\mkern2mu\raise0.7pt\hbox{.}\mkern1mu}}%
  \def\triangle{\mathord{\scalebox{0.9}{$\symup{\Delta}$}}}%
  \def\nabla{\mathord{\raisebox{\depth}{\scalebox{1}[-1]{$\symup{\Delta}$}}}}%
  \def\checkmark{\popcheck{popdark}}%
  \def\star{\mathbin{\ast}}\def\bigstar{\mathord{\ast}}%
  \def\diamond{\mathbin{\scalebox{0.6}{$\Diamond$}}}%
  \def\bigcup{\mathop{\mathchoice{\raisebox{-0.3em}{\scalebox{1.7}{$\cup$}}}{\scalebox{1.25}{$\cup$}}{\cup}{\cup}}\displaylimits}%
  \def\bigcap{\mathop{\mathchoice{\raisebox{-0.3em}{\scalebox{1.7}{$\cap$}}}{\scalebox{1.25}{$\cap$}}{\cap}{\cap}}\displaylimits}%
  \def\lesseqqgtr{\mathrel{\lessgtr}}\def\bigoplus{\mathop{\oplus}\displaylimits}%
}
\providecommand{\ssi}{si et seulement si} % abréviation fréquente
\AtBeginDocument{\providecommand{\wideparen}{\overparen}} % arc de cercle

%% FIGFIX-BEGIN v4 — correctifs globaux des figures (docs/onbuch-generation/tools/figfix)
\usepackage{adjustbox}\usepackage{etoolbox}
% toute figure TikZ/pgfplots plus large que la ligne est réduite à la largeur disponible
\BeforeBeginEnvironment{tikzpicture}{\begin{adjustbox}{max width=\linewidth}}
\AfterEndEnvironment{tikzpicture}{\end{adjustbox}}
% légendes pgfplots lisibles par-dessus les courbes
\pgfplotsset{every axis/.append style={legend style={fill=white,fill opacity=0.92,text opacity=1,draw=black!15,inner sep=3pt}}}
% étiquettes d'arêtes (syntaxe edge["..."]) avec fond blanc
\tikzset{every edge quotes/.append style={fill=white,inner sep=1.2pt}}
%% FIGFIX-END
\begin{document}

\pagegarde{Sensibilisation sur la nécessité d'adaptation de l'Homme avec son environnement}{SVT}{Tle D}{Module III : L'Éducation à l'Environnement et au Développement Durable — Évolution de l'Homme}{N16}{6 h}{Cette leçon établit, à partir de faits paléontologiques, anatomiques, embryologiques et moléculaires, les preuves de l'évolution humaine et les critères de l'hominisation. Elle reconstitue la lignée humaine, construit son arbre phylogénétique avec les grands singes anthropomorphes, et débouche sur la nécessité pour l'Homme moderne d'adapter ses comportements à son environnement dans une logique de développement durable.
\vspace{4pt}
\begin{multicols}{2}\small
\begin{itemize}
\item À la fin de cette leçon, je sais citer et interpréter les preuves paléontologiques, anatomiques, embryologiques et moléculaires de l'évolution de l'Homme
\item À la fin de cette leçon, je sais définir et illustrer les critères de l'hominisation (bipédie, cerveau, face, denture, station debout)
\item À la fin de cette leçon, je sais reconstituer chronologiquement les grandes étapes de la lignée humaine (Australopithèques, Homo habilis, Homo erectus, Néandertalien, Homo sapiens)
\item À la fin de cette leçon, je sais construire et lire un arbre phylogénétique montrant les liens de parenté entre l'Homme et les grands singes anthropomorphes
\item À la fin de cette leçon, je sais exploiter des comparaisons anatomiques (bassins, crânes, membres) et moléculaires (séquences d'ADN, protéines) pour établir des degrés de parenté
\item À la fin de cette leçon, je sais distinguer ce qui relève d'un fait expérimental, d'une interprétation et d'un savoir admis en paléoanthropologie
\end{itemize}\end{multicols}}

\begin{sommaire}
\begin{multicols}{2}
\begin{enumerate}
\item I. Les preuves de l'évolution de l'Homme
\item II. Les critères de l'hominisation
\item III. Les grandes étapes de la lignée humaine
\item IV. Parenté entre l'Homme et les grands singes anthropomorphes : l'arbre phylogénétique
\item V. Capacité d'adaptation de l'Homme à son environnement et développement durable
\item Activité d'intégration
\item Méthodes et exercices résolus
\item Exercices
\item Corrigés des exercices
\item Fiche bilan
\end{enumerate}
\end{multicols}
\end{sommaire}

\begin{objectifs}
\begin{itemize}
\item À la fin de cette leçon, je sais citer et interpréter les preuves paléontologiques, anatomiques, embryologiques et moléculaires de l'évolution de l'Homme
\item À la fin de cette leçon, je sais définir et illustrer les critères de l'hominisation (bipédie, cerveau, face, denture, station debout)
\item À la fin de cette leçon, je sais reconstituer chronologiquement les grandes étapes de la lignée humaine (Australopithèques, Homo habilis, Homo erectus, Néandertalien, Homo sapiens)
\item À la fin de cette leçon, je sais construire et lire un arbre phylogénétique montrant les liens de parenté entre l'Homme et les grands singes anthropomorphes
\item À la fin de cette leçon, je sais exploiter des comparaisons anatomiques (bassins, crânes, membres) et moléculaires (séquences d'ADN, protéines) pour établir des degrés de parenté
\item À la fin de cette leçon, je sais distinguer ce qui relève d'un fait expérimental, d'une interprétation et d'un savoir admis en paléoanthropologie
\item À la fin de cette leçon, je sais établir un lien entre l'évolution humaine et la capacité d'adaptation de l'Homme à son environnement
\item À la fin de cette leçon, je sais expliquer la nécessité du développement durable comme forme moderne d'adaptation de l'Homme à son milieu
\end{itemize}
\end{objectifs}

\begin{prerequis}
\begin{itemize}
\item Notion d'évolution des espèces et de sélection naturelle (classe de Première)
\item Notion de parenté entre espèces et de classification du vivant
\item Structure de l'ADN et comparaison des séquences (génétique de Première)
\item Notion de fossile et de datation des roches (Seconde / début d'année de Terminale)
\item Notions de base sur les écosystèmes et les interactions Homme-milieu (début du Module III)
\end{itemize}
\end{prerequis}

\newpage

\coursec{I. Les preuves de l'évolution de l'Homme}

\begin{savaistu}[Situation d'accroche : les ossements de Shum Laka]
À \cle{Shum Laka}, un abri sous roche célèbre de la région du Nord-Ouest du Cameroun (Grassfields, non loin de Bamenda), des fouilles archéologiques ont livré des ossements humains vieux de plusieurs millénaires. L'analyse de leur \cle{ADN ancien}, publiée en 2020, a révélé des informations inédites sur les populations qui vivaient en Afrique centrale il y a \num{8000} à \num{3000} ans, montrant que ces groupes étaient très différents des populations actuelles de la région. Face à de telles découvertes, un élève de Terminale se pose naturellement trois questions :
\begin{itemize}
\item Comment les scientifiques peuvent-ils affirmer que l'Homme actuel partage des ancêtres communs avec les grands singes ?
\item Quels indices prouvent que notre espèce a évolué au cours du temps ?
\item Alors que le climat du Cameroun change et que les ressources se raréfient, quelles capacités d'adaptation l'Homme possède-t-il pour survivre durablement dans son environnement ?
\end{itemize}
\textbf{Problématique :} sur quels faits observables repose l'affirmation selon laquelle l'Homme est le produit d'une longue évolution biologique ?
\end{savaistu}

Avant de répondre, rappelons un point essentiel : \cle{l'évolution} des espèces, y compris de l'espèce humaine, n'est pas une simple « théorie » au sens vague du terme. C'est un \textbf{fait établi} par la \textbf{convergence de plusieurs catégories de preuves indépendantes} : les fossiles, l'anatomie comparée, l'embryologie et la biologie moléculaire. Chaque catégorie, prise isolément, suggère une histoire évolutive ; leur accord mutuel la rend quasi certaine. Étudions-les une à une, en adoptant la démarche du chercheur : observation d'un document, hypothèse, interprétation, conclusion.

\courssub{Les preuves paléontologiques : les fossiles, archives de l'histoire humaine}

\textbf{Observation.} Dans les couches sédimentaires d'Afrique de l'Est et du Tchad, les paléontologues exhument régulièrement des restes d'êtres qui ressemblent à l'Homme sans lui être identiques : crânes plus petits, faces projetées en avant, bassins différents. Comment interpréter ces vestiges ?

\begin{definition}[Fossile]
Un \cle{fossile} est un reste (os, dent, coquille) ou une empreinte (trace de pas, moulage) d'un organisme ancien, conservé dans les roches sédimentaires grâce à des processus de minéralisation. L'ensemble des fossiles connus constitue le \cle{registre fossile}.
\end{definition}

\begin{experience}[Deux découvertes fondatrices]
\textbf{Document 1 — Lucy.} En 1974, à Hadar (Afar, Éthiopie), l'équipe de Donald Johanson découvre un squelette féminin complet à environ 40 \% : \textbf{Lucy}, \emph{Australopithecus afarensis}, datée d'environ \num{3,2} millions d'années (Ma). Son bassin large et court et l'angle de son fémur montrent qu'elle marchait debout (\cle{bipédie}), mais son volume crânien n'atteint que 400 à \SI{500}{cm^3} (contre environ \SI{1400}{cm^3} chez l'Homme actuel).\par
\textbf{Document 2 — Toumaï.} En 2001, dans le désert du Djourab (Tchad), l'équipe de Michel Brunet découvre un crâne fossile de \emph{Sahelanthropus tchadensis}, surnommé \textbf{Toumaï}, daté d'environ \num{7} Ma. C'est l'un des plus anciens représentants connus de la lignée humaine.\par
\textbf{Datation.} L'âge de ces fossiles est déterminé par les \cle{méthodes radiométriques}, notamment la méthode \cle{potassium-argon} (K/Ar) : le potassium 40 radioactif contenu dans les cendres volcaniques encaissant les fossiles se désintègre en argon 40 avec une période de \num{1,25} milliard d'années ; le rapport Ar/K mesuré dans la roche donne son âge absolu.\par
\textbf{Interprétation.} En classant les fossiles du plus ancien au plus récent, on observe une \textbf{succession ordonnée de formes} : le volume crânien augmente progressivement, la face devient de moins en moins \cle{prognathe} (projetée en avant), la bipédie se perfectionne.\par
\textbf{Conclusion.} Les espèces humaines ont changé au cours du temps : c'est une preuve directe de l'évolution.
\end{experience}

\begin{popfigure}
\begin{center}
\begin{tikzpicture}[font=\small]
\draw[popline,thick] (0,0) -- (12,0);
\foreach \x/\date/\nom/\desc in {1/{7 Ma}/{Toumaï}/{\emph{S. tchadensis}},
4/{3,2 Ma}/{Lucy}/{\emph{A. afarensis}},
7/{1,8 Ma}/{\emph{H. erectus}}/{bipède, outils},
10/{0,3 Ma}/{\emph{H. sapiens}}/{Homme actuel}}{
\fill[poporange] (\x,0) circle (2.2pt);
\draw[popline] (\x,0) -- (\x,0.5);
\node[above,align=center,text=popdark] at (\x,0.55) {\textbf{\nom}\\ \desc};
\node[below,text=popblue] at (\x,-0.15) {\date};
}
\node[right,text=popmuted] at (12.2,0) {temps};
\draw[->,popmuted] (12,0) -- (12.9,0);
\end{tikzpicture}
\end{center}
\legende{Frise chronologique schématique de quelques fossiles majeurs de la lignée humaine. Les dates (Ma = millions d'années) sont obtenues par datation radiométrique des couches volcaniques encaissantes.}
\end{popfigure}

\begin{aretenir}
Les fossiles montrent que des espèces aujourd'hui disparues, anatomiquement intermédiaires entre les grands singes et l'Homme actuel, ont existé et se sont succédé dans le temps. Le registre fossile humain, quoique lacunaire, documente un \textbf{changement progressif des caractères} sur environ 7 millions d'années.
\end{aretenir}

\courssub{Les preuves anatomiques comparées : homologies et organes vestigiaux}

\textbf{Observation.} Disseque (en imagination ou en salle de TP) le membre antérieur d'un Homme, d'un chimpanzé et d'une chauve-souris. Les fonctions diffèrent totalement : saisir, grimper, voler. Pourtant, le plan d'organisation osseuse est le même : un os au bras (humérus), deux à l'avant-bras (radius et cubitus), des carpes, des métacarpes et des phalanges.

\begin{definition}[Homologie]
On dit que deux organes sont \cle{homologues} lorsqu'ils présentent une ressemblance de structure (même plan d'organisation, même origine embryonnaire) malgré des fonctions parfois différentes. L'homologie est un indice d'\textbf{ascendance commune} : les organes homologues ont hérité leur plan d'un ancêtre commun, puis se sont diversifiés par adaptation à des modes de vie différents.
\end{definition}

\begin{popfigure}
\begin{center}
\begin{tikzpicture}[font=\footnotesize,scale=0.95]
% Homme
\begin{scope}[shift={(0,0)}]
\draw[fill=popblueL,draw=popblue] (0,2.2) rectangle (0.35,3.4);
\draw[fill=popgreenL,draw=popgreen] (-0.05,1.1) rectangle (0.16,2.15);
\draw[fill=popgreenL,draw=popgreen] (0.22,1.1) rectangle (0.42,2.15);
\draw[fill=poporangeL,draw=poporange] (0.0,0.55) rectangle (0.38,1.05);
\foreach \i in {0,...,4}{\draw[fill=poppinkL,draw=poppink] (-0.08+0.11*\i,0.05) rectangle (0.0+0.11*\i,0.5);}
\node[align=center] at (0.2,-0.5) {\textbf{Homme}\\ saisir};
\end{scope}
% Chimpanzé
\begin{scope}[shift={(3.5,0)}]
\draw[fill=popblueL,draw=popblue] (0,2.2) rectangle (0.35,3.6);
\draw[fill=popgreenL,draw=popgreen] (-0.05,0.9) rectangle (0.16,2.15);
\draw[fill=popgreenL,draw=popgreen] (0.22,0.9) rectangle (0.42,2.15);
\draw[fill=poporangeL,draw=poporange] (0.0,0.35) rectangle (0.38,0.85);
\foreach \i in {0,...,4}{\draw[fill=poppinkL,draw=poppink] (-0.08+0.11*\i,-0.35) rectangle (0.0+0.11*\i,0.3);}
\node[align=center] at (0.2,-0.9) {\textbf{Chimpanzé}\\ grimper};
\end{scope}
% Chauve-souris
\begin{scope}[shift={(7,0)}]
\draw[fill=popblueL,draw=popblue] (0,2.2) rectangle (0.35,3.4);
\draw[fill=popgreenL,draw=popgreen] (-0.05,1.1) rectangle (0.16,2.15);
\draw[fill=popgreenL,draw=popgreen] (0.22,1.1) rectangle (0.42,2.15);
\draw[fill=poporangeL,draw=poporange] (0.0,0.6) rectangle (0.38,1.05);
\foreach \i/\len in {0/1.3,1/1.7,2/1.9,3/1.7,4/1.2}{
\draw[fill=poppinkL,draw=poppink] (0.05+0.35*\i,0.55) rectangle (0.13+0.35*\i,0.55-\len*0.5);}
\draw[poppurple,dashed] (0.05,0.55) -- (1.85,-0.4) -- (0.13,0.55);
\node[align=center] at (0.9,-1.4) {\textbf{Chauve-souris}\\ voler};
\end{scope}
% Légende couleurs
\node[fill=popblueL,draw=popblue,minimum width=0.3cm,minimum height=0.2cm] at (10.2,3.2) {};
\node[right] at (10.4,3.2) {humérus};
\node[fill=popgreenL,draw=popgreen,minimum width=0.3cm,minimum height=0.2cm] at (10.2,2.7) {};
\node[right] at (10.4,2.7) {radius et cubitus};
\node[fill=poporangeL,draw=poporange,minimum width=0.3cm,minimum height=0.2cm] at (10.2,2.2) {};
\node[right] at (10.4,2.2) {carpes};
\node[fill=poppinkL,draw=poppink,minimum width=0.3cm,minimum height=0.2cm] at (10.2,1.7) {};
\node[right] at (10.4,1.7) {métacarpes et phalanges};
\end{tikzpicture}
\end{center}
\legende{Homologie des membres antérieurs de l'Homme, du chimpanzé et de la chauve-souris : même plan d'organisation osseuse (mêmes couleurs = os homologues), fonctions différentes. Ce plan commun s'explique par l'héritage d'un ancêtre commun.}
\end{popfigure}

\textbf{Interprétation.} Si chaque espèce avait été créée indépendamment pour sa fonction, rien n'imposerait que l'aile d'une chauve-souris contienne les mêmes os qu'une main humaine. L'explication la plus simple et la plus économique est la \textbf{filiation} : un plan ancestral commun, modifié par l'évolution.

\textbf{Organes vestigiaux.} L'anatomie humaine conserve aussi des structures réduites et sans fonction majeure, héritées d'ancêtres chez qui elles étaient développées : l'\cle{appendice} (reste d'un caecum développé chez les herbivores), le \cle{coccyx} (vestige d'une queue), les \cle{dents de sagesse} (souvent inutiles, parfois absentes, héritées d'ancêtres à mâchoires plus longues). Ces « reliques » anatomiques sont des témoins de notre histoire évolutive.

\courssub{Les preuves embryologiques : le développement révèle la parenté}

\textbf{Observation.} Compare des embryons précoces d'Homme, de chimpanzé, de poulet et de poisson : ils se ressemblent étonnamment. L'embryon humain de 4 à 5 semaines possède des \cle{fentes branchiales} (sillons pharyngés) et une \cle{queue embryonnaire} bien visible, qui régressera ensuite (elle ne laissera que le coccyx).

\textbf{Interprétation.} Ces structures, inutiles à l'Homme adulte, rappellent des états ancestraux : les fentes branchiales sont fonctionnelles chez les poissons, la queue chez la plupart des vertébrés. Le développement embryonnaire conserve ainsi des traces du passé évolutif. Plus les embryons de deux espèces restent semblables longtemps, plus leur parenté est proche : les embryons d'Homme et de chimpanzé sont quasi indiscernables aux stades précoces.

\begin{attention}[Ne pas caricaturer]
Il ne faut pas dire que « l'embryon humain passe par un stade poisson puis un stade singe ». L'embryon ne rejoue pas l'évolution de façon littérale ; il partage simplement avec d'autres vertébrés des \textbf{étapes de développement communes héritées d'ancêtres communs}, ce qui constitue un indice de parenté.
\end{attention}

\courssub{Les preuves moléculaires : l'ADN, archive de la parenté}

\textbf{Observation.} Depuis les années 1980, on sait lire et comparer les séquences d'\cle{ADN} et de protéines (cytochrome c, hémoglobine) de différentes espèces. Résultat frappant : l'ADN de l'Homme et celui du chimpanzé sont identiques à plus de \textbf{98 \%}.

\begin{center}
\begin{tabularx}{\linewidth}{|l|X|X|}\hline
\rowcolor{popblueL} \textbf{Espèce comparée à l'Homme} & \textbf{Identité d'ADN approximative} & \textbf{Ancêtre commun estimé} \\ \hline
Chimpanzé & $\approx 98{,}8\ \%$ & $\approx$ 6 à 7 Ma \\ \hline
Gorille & $\approx 98\ \%$ & $\approx$ 8 à 10 Ma \\ \hline
Orang-outan & $\approx 97\ \%$ & $\approx$ 12 à 16 Ma \\ \hline
\end{tabularx}
\end{center}

\textbf{Interprétation.} Les mutations s'accumulent dans l'ADN au fil des générations, à un rythme à peu près régulier (on parle d'\emph{horloge moléculaire}). Par conséquent, \textbf{plus deux espèces présentent de différences dans leurs séquences, plus leur ancêtre commun est ancien}. La proximité moléculaire Homme--chimpanzé indique un ancêtre commun récent à l'échelle géologique.

\begin{methode}[Établir un degré de parenté à partir de séquences d'ADN]
\begin{enumerate}
\item Aligner les séquences homologues (même gène) des espèces comparées.
\item Compter les différences (substitutions de nucléotides) entre chaque paire d'espèces.
\item Calculer le pourcentage d'identité : $\text{identité} = \dfrac{\text{sites identiques}}{\text{sites totaux}} \times 100$.
\item Conclure : moins il y a de différences, plus la parenté est proche et plus l'ancêtre commun est récent. On peut alors ordonner les espèces et proposer un arbre de parenté.
\end{enumerate}
\end{methode}

\begin{exoresolu}[Comparaison de séquences]
On compare un fragment de 100 nucléotides d'un même gène chez trois espèces. Homme--chimpanzé : 2 différences ; Homme--gorille : 3 différences ; chimpanzé--gorille : 3 différences. Quel couple d'espèces est le plus proche parent ? Justifier.
\tcblower
Le couple Homme--chimpanzé présente le moins de différences (2 sur 100, soit 98 \% d'identité). D'après le principe de l'horloge moléculaire, moins il y a de différences, plus l'ancêtre commun est récent. L'Homme est donc plus proche du chimpanzé que du gorille ; l'ancêtre commun Homme--chimpanzé est postérieur à l'ancêtre commun des trois espèces.
\end{exoresolu}

\courssub{Faits et interprétations : la rigueur du raisonnement scientifique}

\begin{attention}[Erreur fréquente à bannir]
\textbf{L'Homme ne descend pas du singe actuel.} Ni du chimpanzé, ni du gorille. Homme et chimpanzé sont des « cousins » évolutifs : ils partagent un \textbf{ancêtre commun disparu} (il y a 6 à 7 Ma), dont chaque lignée a évolué de son côté. Dire « descendre du singe » est une caricature qui fausse tout le raisonnement phylogénétique.
\end{attention}

\begin{aretenir}[Faits versus interprétations]
Il faut distinguer clairement :
\begin{itemize}
\item les \textbf{faits} : les fossiles existent, les séquences d'ADN sont mesurables, les homologies anatomiques sont observables — ce sont des données objectives ;
\item les \textbf{interprétations} : la filiation entre les formes, la position exacte de chaque fossile dans l'arbre — ce sont des reconstructions scientifiques \textbf{robustes} (car fondées sur la convergence de preuves indépendantes) mais \textbf{révisables} : chaque nouveau fossile (comme Toumaï en 2001) peut préciser ou modifier le scénario. C'est précisément cette autocorrection qui fait la force de la science.
\end{itemize}
\end{aretenir}

\begin{savaistu}[Shum Laka, une fierté scientifique camerounaise]
Le site de Shum Laka (Nord-Ouest du Cameroun) a livré des ADN anciens de plusieurs individus datés d'environ \num{8000} à \num{3000} ans. Publiées en 2020, ces analyses ont montré que les anciens habitants de l'Ouest camerounais étaient génétiquement distincts des locuteurs bantous actuels, et ont éclairé l'histoire des migrations humaines en Afrique centrale. Preuve que le Cameroun est un acteur majeur de la recherche sur l'évolution humaine !
\end{savaistu}

\textbf{Bilan de la section.} Quatre catégories de preuves indépendantes — paléontologiques, anatomiques, embryologiques et moléculaires — convergent : l'espèce humaine est le résultat d'une évolution, et elle partage une parenté étroite avec les grands singes. Reste à comprendre \emph{comment} cette évolution s'est concrétisée : c'est l'objet des critères de l'hominisation, que nous étudierons dans la section suivante.

\coursec{Les critères de l'hominisation}

La section précédente nous a convaincus, fossiles et arguments anatomiques à l'appui, que l'Homme est le produit d'une longue évolution. Mais une question surgit immédiatement : face à un os fossile découvert dans les sédiments de l'Afrique de l'Est, comment le paléoanthropologue décide-t-il qu'il appartient à la \cle{lignée humaine} plutôt qu'à celle des grands singes ? Il ne suffit pas de dire « cela ressemble à un humain » : il faut des \cle{critères} précis, mesurables, objectifs. Ces critères, réunis sous le nom de \cle{critères de l'hominisation}, sont l'objet de cette section. Tu verras qu'ils ne sont pas apparus en même temps, et que le cerveau — contrairement à une idée très répandue — n'est pas venu en premier.

\courssub{Définition de l'hominisation}

\begin{definition}[Hominisation]
L'\cle{hominisation} est l'ensemble des transformations \cle{morphologiques} (bipédie, augmentation du volume crânien, réduction de la face, libération de la main) et \cle{comportementales} ou culturelles (fabrication d'outils, maîtrise du feu, langage, sépultures, art) qui, au cours de plusieurs millions d'années, ont conduit de l'ancêtre commun aux grands singes vers la lignée humaine actuelle, \emph{Homo sapiens}.
\end{definition}

\begin{attention}[L'hominisation est un processus en mosaïque]
Les critères de l'hominisation ne sont \textbf{pas apparus simultanément}. La bipédie est acquise très tôt (il y a environ 7 à 4 millions d'années), alors que l'augmentation importante du volume crânien est beaucoup plus tardive (à partir d'environ 2,5 millions d'années). On dit que l'hominisation s'est faite « en \cle{mosaïque} » : chaque caractère évolue à son propre rythme. Ne rédige jamais au Bac que « les Australopithèques avaient un gros cerveau parce qu'ils étaient bipèdes » : c'est faux.
\end{attention}

\courssub{Démarche d'investigation : comparer pour dégager les critères}

\begin{experience}[TP — Comparaison de crânes et de bassins]
\textbf{Matériel :} moulages (ou photographies cotées) de crânes et de bassins de chimpanzé (\emph{Pan troglodytes}), d'Australopithèque (\emph{Australopithecus afarensis}, spécimen « Lucy ») et d'Homme moderne (\emph{Homo sapiens}) ; pied à coulisse ; ficelle et rapporteur.

\textbf{Protocole :}
\begin{enumerate}
\item Observer chaque crâne de profil et repérer la position du \cle{trou occipital} (foramen magnum) : vers l'arrière, intermédiaire ou sous le crâne (centré).
\item Mesurer ou estimer le volume de la boîte crânienne de chaque spécimen.
\item Observer le prognathisme (projection des mâchoires vers l'avant) et la forme de l'arcade dentaire.
\item Comparer les bassins : haut et étroit, ou court et large.
\end{enumerate}

\textbf{Observations attendues :}

\begin{center}
\begin{tabularx}{\linewidth}{|l|X|X|X|}
\hline
\rowcolor{popblueL} \textbf{Critère} & \textbf{Chimpanzé} & \textbf{Australopithèque} & \textbf{Homo sapiens} \\
\hline
Trou occipital & Vers l'arrière & Position intermédiaire & Centré sous le crâne \\
\hline
Volume crânien & $\approx$ \SI{400}{cm^3} & $\approx$ \SI{450}{cm^3} & $\approx$ \SI{1350}{cm^3} \\
\hline
Prognathisme & Très marqué & Moyen & Absent (face verticale, menton) \\
\hline
Arcade dentaire & En U & Intermédiaire & En parabole \\
\hline
Bassin & Haut et étroit & Court et élargi & Court et large \\
\hline
\end{tabularx}
\end{center}

\textbf{Interprétation :} l'Australopithèque présente des caractères « humains » (position du trou occipital, bassin) et des caractères « simiesques » (faible volume crânien, prognathisme). C'est un \cle{intermédiaire évolutif}. \textbf{Conclusion :} l'hominisation se définit par un faisceau de critères apparus progressivement.
\end{experience}

\begin{popfigure}
\begin{tikzpicture}[scale=0.9, every node/.style={font=\small}]
% Chimpanzé
\begin{scope}[shift={(0,0)}]
\draw[thick, popdark, fill=popcream] (0.4,0.6) ellipse (1.1 and 0.85);
\draw[thick, popdark, fill=popcream] (1.7,0.25) ellipse (0.85 and 0.45);
\fill[popink] (0.95,-0.15) circle (0.09);
\draw[->, popink, thick] (0.95,-0.15) -- (0.95,-0.75) node[below, font=\footnotesize]{trou occipital (arrière)};
\node[font=\bfseries] at (0.9,1.9) {Chimpanzé};
\node[font=\footnotesize, popmuted] at (0.9,1.55) {$\approx$ \SI{400}{cm^3}, prognathe};
\end{scope}
% Australopithèque
\begin{scope}[shift={(5.2,0)}]
\draw[thick, popdark, fill=poporangeL] (0.3,0.7) ellipse (1.05 and 0.95);
\draw[thick, popdark, fill=poporangeL] (1.45,0.2) ellipse (0.7 and 0.4);
\fill[popink] (0.35,-0.2) circle (0.09);
\draw[->, popink, thick] (0.35,-0.2) -- (0.35,-0.75) node[below, font=\footnotesize]{position intermédiaire};
\node[font=\bfseries] at (0.75,2.0) {Australopithèque};
\node[font=\footnotesize, popmuted] at (0.75,1.65) {$\approx$ \SI{450}{cm^3}, bipède};
\end{scope}
% Homo sapiens
\begin{scope}[shift={(10.4,0)}]
\draw[thick, popdark, fill=popgreenL] (0,0.85) ellipse (1.15 and 1.15);
\draw[thick, popdark, fill=popgreenL] (0.95,0.05) ellipse (0.5 and 0.35);
\fill[popink] (0,-0.28) circle (0.09);
\draw[->, popink, thick] (0,-0.28) -- (0,-0.85) node[below, font=\footnotesize]{trou occipital centré};
\node[font=\bfseries] at (0.3,2.25) {\emph{Homo sapiens}};
\node[font=\footnotesize, popmuted] at (0.3,1.9) {$\approx$ \SI{1350}{cm^3}, face verticale, menton};
\end{scope}
\end{tikzpicture}
\legende{Crânes vus de profil (schématisés) : chez le chimpanzé, la boîte crânienne est petite, la face projetée en avant (prognathisme) et le trou occipital dirigé vers l'arrière ; chez \emph{Homo sapiens}, la boîte crânienne est volumineuse, la face verticale sous le crâne et le trou occipital centré, signe d'un port vertical de la tête.}
\end{popfigure}

\courssub{Critère 1 : la bipédie, premier critère de l'hominisation}

\begin{definition}[Bipédie]
La \cle{bipédie} est un mode de locomotion dans lequel l'animal se déplace debout, sur ses deux seuls membres postérieurs. La bipédie permanente et exclusive est le \cle{premier critère de l'hominisation} : c'est le caractère le plus anciennement acquis dans la lignée humaine.
\end{definition}

La bipédie n'est pas un simple « détail de déplacement » : elle a entraîné une véritable réorganisation de tout le squelette. Cinq indices anatomiques permettent de la reconnaître chez un fossile.

\begin{enumerate}
\item \textbf{La position du trou occipital (foramen magnum).} Le trou occipital est l'orifice par lequel la moelle épinière rejoint l'encéphale. Chez le singe quadrupède, la colonne vertébrale est horizontale et le trou occipital est situé \emph{vers l'arrière} du crâne. Chez l'Homme bipède, la tête est portée verticalement, en équilibre au sommet de la colonne : le trou occipital est \emph{centré sous le crâne}.
\item \textbf{Les courbures de la colonne vertébrale.} La colonne humaine n'est pas droite : elle présente quatre courbures qui amortissent les chocs de la marche — deux \cle{lordoses} (cervicale et lombaire, courbures vers l'avant) et deux \cle{cyphoses} (dorsale et sacrée, courbures vers l'arrière). La colonne du singe quadrupède forme une seule courbe en arc.
\item \textbf{Le bassin court et large.} Chez l'Homme, le bassin est élargi en « coupole » : il soutient les viscères en position verticale et offre une large surface d'insertion aux muscles fessiers qui stabilisent le tronc pendant la marche. Le bassin du chimpanzé est haut et étroit.
\item \textbf{Le fémur oblique.} Chez l'Homme, les fémurs convergent vers les genoux (angle bicondylien), ce qui place les pieds sous le centre de gravité du corps : la marche est stable et économe.
\item \textbf{Le pied voûté.} Le pied humain forme une voûte plantaire amortissante, et le \cle{gros orteil} (hallux) est aligné avec les autres orteils, \emph{non opposable} : il sert de point d'appui propulsif. Chez le singe, le gros orteil est opposable, car le pied reste une main de préhension arboricole.
\end{enumerate}

\begin{methode}[Déterminer si un fossile était bipède à partir d'un crâne]
\begin{enumerate}
\item Observer le crâne de profil et repérer le trou occipital (foramen magnum).
\item Situer sa position : vers l'arrière du crâne, intermédiaire, ou centrée sous le crâne.
\item Conclure : un trou occipital \textbf{centré} signifie que la tête s'équilibrait verticalement au sommet d'une colonne dressée : l'individu était \textbf{bipède}. Un trou occipital rejeté vers l'arrière indique une colonne horizontale : l'animal était quadrupède.
\item Confirmer si possible par d'autres indices (bassin large, fémur oblique, gros orteil non opposable).
\end{enumerate}
\end{methode}

\begin{attention}[Piège classique du Bac]
Erreur fréquente : affirmer que « le grand volume crânien est le premier critère de l'hominisation ». C'est \textbf{faux}. Le fossile « Lucy » (\emph{Australopithecus afarensis}, environ 3,2 Ma) présente un bassin large, un fémur oblique et un trou occipital en position avancée : elle était \textbf{bipède}… avec un cerveau d'environ \SI{450}{cm^3} (moyenne de l'espèce, Lucy \emph{stricto sensu} $\approx$ \SI{400}{cm^3}), à peine plus gros que celui d'un chimpanzé. La bipédie \emph{précède} de plusieurs millions d'années l'augmentation du volume crânien.
\end{attention}

\courssub{Critère 2 : l'augmentation du volume crânien}

Le deuxième critère est l'accroissement progressif de la boîte crânienne et donc de l'encéphale, en particulier du \cle{cortex cérébral} (siège des fonctions supérieures : langage, abstraction, fabrication d'outils). Les valeurs moyennes parlent d'elles-mêmes.

\begin{exemplebox}[Volumes crâniens moyens dans la lignée humaine]
\begin{center}
\begin{tabularx}{\linewidth}{|l|X|}
\hline
\rowcolor{popblueL} \textbf{Espèce} & \textbf{Volume crânien moyen} \\
\hline
Chimpanzé (\emph{Pan troglodytes}) & $\approx$ \SI{400}{cm^3} \\
\hline
Australopithèque (\emph{A. afarensis}) & $\approx$ \SI{450}{cm^3} \\
\hline
\emph{Homo habilis} & $\approx$ \SI{650}{cm^3} \\
\hline
\emph{Homo erectus} & $\approx$ \SI{900}{cm^3} \\
\hline
\emph{Homo sapiens} & $\approx$ \SI{1350}{cm^3} \\
\hline
\end{tabularx}
\end{center}
Entre l'Australopithèque et l'Homme moderne, le volume crânien a été multiplié par trois, en moins de 3 millions d'années.
\end{exemplebox}

\begin{popfigure}
\begin{tikzpicture}
\begin{axis}[
width=\linewidth, height=7cm,
xlabel={Temps (millions d'années avant le présent)},
ylabel={Volume crânien (cm$^3$)},
xlabel style={font=\small}, ylabel style={font=\small},
x dir=reverse,
xmin=7.5, xmax=-0.3,
ymin=0, ymax=1500,
grid=both, grid style={popline!40},
legend style={font=\footnotesize, at={(0.35,0.97)}, anchor=north},
tick label style={font=\footnotesize}
]
\addplot[popcurve, mark=*, mark size=1.5pt] coordinates {
(7,350) (4,400) (3.2,450) (2.4,650) (1.8,900) (0.3,1400) (0,1350)
};
\legend{Volume crânien moyen}
\addplot[only marks, mark=none] coordinates {(3.2,450)} node[above right, font=\footnotesize, popdark]{Australopithèque};
\addplot[only marks, mark=none] coordinates {(2.4,650)} node[above, font=\footnotesize, popdark]{\emph{H. habilis}};
\addplot[only marks, mark=none] coordinates {(1.8,900)} node[above, font=\footnotesize, popdark]{\emph{H. erectus}};
\addplot[only marks, mark=none] coordinates {(0,1350)} node[above left, font=\footnotesize, popdark]{\emph{H. sapiens}};
\end{axis}
\end{tikzpicture}
\legende{Évolution du volume crânien au cours du temps. La courbe reste presque plate pendant plus de 4 millions d'années (la bipédie est déjà acquise), puis s'élève fortement à partir du genre \emph{Homo} (environ 2,5 Ma) : preuve que l'augmentation du cerveau est tardive par rapport à la bipédie.}
\end{popfigure}

\courssub{Critère 3 : la réduction de la face}

\begin{definition}[Prognathisme]
Le \cle{prognathisme} est la projection en avant de la face, c'est-à-dire l'avancée des mâchoires (os maxillaires et mandibule) par rapport au plan vertical du front. Très marqué chez les grands singes, il s'est progressivement \cle{réduit} au cours de l'hominisation, jusqu'à donner chez \emph{Homo sapiens} une face verticale, placée sous la boîte crânienne.
\end{definition}

La réduction de la face s'accompagne de plusieurs transformations liées au régime alimentaire (nourriture plus tendre, puis cuite) et à la diminution des contraintes de mastication :
\begin{itemize}
\item la \textbf{réduction des canines}, qui perdent leur fonction de défense et d'intimidation ;
\item la transformation de l'\textbf{arcade dentaire} : en \textbf{U} (rangées presque parallèles) chez le singe, elle devient une \textbf{parabole} régulière chez l'Homme ;
\item l'apparition d'un \textbf{menton} (éminence mentonnière), caractère propre à \emph{Homo sapiens}.
\end{itemize}

\begin{popfigure}
\begin{tikzpicture}[scale=0.95, every node/.style={font=\small}]
% Arcade en U (singe)
\begin{scope}[shift={(0,0)}]
\draw[very thick, poporange, rounded corners=2pt] (-1.1,1.1) -- (-1.1,-0.6) arc[start angle=180, end angle=360, radius=1.1] -- (1.1,1.1);
\foreach \x in {-0.9,-0.45,0,0.45,0.9} {\fill[popdark] (\x,-1.02) circle (0.07);}
\foreach \x in {-1.1,1.1} {\foreach \y in {0.9,0.45,0} {\fill[popdark] (\x,\y) circle (0.07);}}
\node[font=\bfseries] at (0,1.6) {Singe : arcade en U};
\node[font=\footnotesize, popmuted] at (0,-1.55) {rangées parallèles, grandes canines};
\end{scope}
% Arcade en parabole (Homme)
\begin{scope}[shift={(5.5,0)}]
\draw[very thick, popgreen, domain=-1.15:1.15, samples=40] plot (\x, {1.1 - 2.1*(1 - \x*\x/1.32)});
\foreach \x in {-1.05,-0.75,-0.45,-0.15,0.15,0.45,0.75,1.05} {\fill[popdark] (\x, {1.1 - 2.1*(1 - \x*\x/1.32)}) circle (0.07);}
\node[font=\bfseries] at (0,1.6) {Homme : arcade en parabole};
\node[font=\footnotesize, popmuted] at (0,-1.55) {courbe régulière, canines réduites};
\end{scope}
% Bassins
\begin{scope}[shift={(11,0)}]
\draw[thick, popdark, fill=popblueL] (-0.35,1.2) ellipse (0.28 and 1.15);
\draw[thick, popdark, fill=popblueL] (0.35,1.2) ellipse (0.28 and 1.15);
\node[font=\footnotesize] at (0,-0.35) {Singe : haut, étroit};
\begin{scope}[shift={(2.6,0)}]
\draw[thick, popdark, fill=popgreenL] (-0.75,0.75) ellipse (0.85 and 0.55);
\draw[thick, popdark, fill=popgreenL] (0.75,0.75) ellipse (0.85 and 0.55);
\node[font=\footnotesize] at (0,-0.35) {Homme : court, large};
\end{scope}
\node[font=\bfseries] at (1.3,1.9) {Bassins comparés};
\end{scope}
\end{tikzpicture}
\legende{À gauche : arcades dentaires vues d'en bas — en U avec rangées parallèles chez le singe, en parabole chez l'Homme. À droite : bassins schématisés — haut et étroit chez le quadrupède, court et large en coupole chez le bipède (soutien des viscères, insertion des muscles fessiers).}
\end{popfigure}

\courssub{Critère 4 : la libération de la main}

Conséquence directe de la bipédie : les membres antérieurs ne servant plus à la locomotion, la main est « libérée » et peut se spécialiser. La main humaine possède un \cle{pouce opposable} long et mobile, qui peut venir au contact de chacun des autres doigts. Cette \cle{préhension fine} (pince pouce-index) permet de saisir, manipuler, tailler : c'est le support anatomique de la fabrication d'outils. On parle de couple fonctionnel « main--cerveau » : le développement du cortex a permis des gestes plus habiles, et la pratique de ces gestes a en retour favorisé le développement cérébral.

\courssub{Critère 5 : les critères culturels}

L'hominisation n'est pas seulement anatomique : elle est aussi \cle{comportementale}. Les traces matérielles retrouvées par les archéologues attestent de comportements de plus en plus complexes :
\begin{itemize}
\item la \textbf{fabrication d'outils} : galets taillés (environ 2,5 Ma, associés à \emph{Homo habilis}), puis \textbf{bifaces} soigneusement façonnés (associés à \emph{Homo erectus}) ;
\item la \textbf{maîtrise du feu} (environ 400 000 ans, peut-être davantage) : cuisson des aliments, protection, vie sociale autour du foyer ;
\item les \textbf{sépultures} : inhumation volontaire des morts (Néandertaliens, puis \emph{Homo sapiens}), indice d'une pensée symbolique ;
\item l'\textbf{art pariétal} : peintures et gravures des grottes ornées, expression d'une pensée abstraite ;
\item le \textbf{langage articulé}, rendu possible par l'évolution du cerveau, du pharynx et du larynx : outil de communication et de transmission culturelle.
\end{itemize}

\begin{aretenir}[Les cinq critères de l'hominisation]
\begin{enumerate}
\item \textbf{Bipédie} : trou occipital centré, colonne à courbures (lordoses et cyphoses), bassin court et large, fémur oblique, pied voûté à gros orteil non opposable. \emph{Premier critère apparu.}
\item \textbf{Augmentation du volume crânien} : de $\approx$ \SI{400}{cm^3} (Australopithèque) à $\approx$ \SI{1350}{cm^3} (\emph{Homo sapiens}), avec développement du cortex.
\item \textbf{Réduction de la face} : disparition du prognathisme, canines réduites, arcade en parabole, apparition du menton.
\item \textbf{Libération de la main} : pouce opposable, préhension fine, fabrication d'outils.
\item \textbf{Critères culturels} : outils, feu, sépultures, art, langage.
\end{enumerate}
Ces critères sont apparus à des époques différentes : l'hominisation est un processus \textbf{en mosaïque}.
\end{aretenir}

\begin{exoresolu}[Identifier un fossile]
Au cours d'une fouille dans la vallée de l'Omo (Éthiopie), on découvre un crâne fossile dont le volume est estimé à \SI{500}{cm^3}. Le trou occipital occupe une position avancée, presque centrée sous le crâne, et l'arcade dentaire est intermédiaire entre le U et la parabole. Ce fossile appartenait-il à un bipède ? Peut-on le classer dans le genre \emph{Homo} ? Justifier.
\tcblower
\textbf{1. La bipédie.} Le trou occipital est en position avancée, presque centrée : la tête s'équilibrait verticalement au sommet d'une colonne vertébrale dressée. Le fossile appartenait donc à un individu \textbf{bipède}.

\textbf{2. L'appartenance au genre \emph{Homo}.} Le volume crânien de \SI{500}{cm^3} est très inférieur à celui des représentants du genre \emph{Homo} (au moins $\approx$ \SI{650}{cm^3} chez \emph{Homo habilis}), et l'arcade dentaire reste intermédiaire. Ce fossile ne peut donc \textbf{pas} être classé dans le genre \emph{Homo} : il s'agit plutôt d'un \textbf{Australopithèque}, bipède à cerveau encore petit — ce qui illustre le caractère \emph{en mosaïque} de l'hominisation.
\end{exoresolu}

\begin{exercice}[À toi de jouer]
On découvre un bassin fossile court et très élargi, associé à un fémur dont l'axe converge fortement vers le genou. Le crâne retrouvé à proximité présente un fort prognathisme et un volume de \SI{430}{cm^3}. (1) L'individu était-il bipède ? Justifier par deux arguments. (2) Peut-on dire qu'il s'agit d'un Homme moderne ? Expliquer pourquoi la réponse n'est pas contradictoire avec la bipédie.
\end{exercice}

\begin{corrige}[Exercice]
(1) Oui : le bassin court et large soutient les viscères en position verticale et le fémur oblique place les pieds sous le centre de gravité : ce sont deux indices certains de bipédie. (2) Non : le fort prognathisme et le faible volume crânien (\SI{430}{cm^3}) sont des caractères archaïques incompatibles avec \emph{Homo sapiens}. Il n'y a pas contradiction : l'hominisation s'est faite en mosaïque — la bipédie est apparue plusieurs millions d'années avant l'augmentation du volume crânien et la réduction de la face. Ce fossile correspond au profil d'un Australopithèque.
\end{corrige}

\begin{savaistu}[Lucy, la grand-mère de l'humanité]
Découverte en 1974 à Hadar (Éthiopie) par l'équipe de Donald Johanson, « Lucy » (\emph{Australopithecus afarensis}, 3,2 Ma) est l'un des fossiles les plus célèbres du monde : son squelette était complet à près de 40 \%. Son bassin et son fémur prouvent qu'elle marchait debout, alors que son cerveau ne dépassait pas \SI{450}{cm^3}. C'est elle qui a définitivement démontré que la bipédie avait précédé le « gros cerveau » — renversant les idées reçues de l'époque. Les empreintes de pas fossilisées de Laetoli (Tanzanie, 3,6 Ma), laissées dans des cendres volcaniques par des Australopithèques marchant à deux pieds, confirment spectaculairement cette conclusion.
\end{savaistu}

En résumé, l'hominisation se lit dans les os : un crâne posé verticalement, un bassin en coupole, une face qui s'efface sous une boîte crânienne grandissante, une main devenue outil. Mais ces critères ne disent rien, encore, de l'\emph{ordre} ni des \emph{acteurs} de cette histoire : c'est l'objet de la section suivante, où nous reconstituerons les grandes étapes de la lignée humaine, de \emph{Sahelanthropus} à \emph{Homo sapiens}.

\coursec{III. Les grandes étapes de la lignée humaine}

Dans la section précédente, nous avons identifié les \cle{critères de l'hominisation} : la bipédie permanente, l'augmentation du volume crânien, la réduction du prognathisme et de la denture, la fabrication d'outils et l'apparition de comportements culturels. Ces critères ne sont pas apparus simultanément : ils se sont mis en place progressivement, au fil de millions d'années, chez des espèces successives — et souvent contemporaines les unes des autres. L'objectif de cette section est de \cle{reconstituer la chronologie de la lignée humaine}, de décrire les principales espèces qui la composent et de montrer que cette histoire ressemble à un \emph{buisson} plutôt qu'à une échelle.

\courssub{Qu'est-ce qu'un homininé ?}

\begin{definition}[Homininé]
Un \cle{homininé} est un représentant de la \cle{lignée humaine}, c'est-à-dire l'ensemble des espèces apparues \emph{après} la séparation évolutive d'avec la lignée du chimpanzé (il y a environ 7 à 8 millions d'années), et plus étroitement apparentées à l'Homme actuel qu'au chimpanzé. Tous les homininés fossiles sont éteints, à l'exception d'\emph{Homo sapiens}.
\end{definition}

\begin{attention}[Ne pas confondre]
Hominidés (famille incluant les grands singes : gorilles, chimpanzés, orangs-outans \emph{et} l'Homme), homininés (lignée humaine \emph{stricto sensu}) et hominines (sous-ensemble encore plus restreint) ne sont pas synonymes. Au Baccalauréat, emploie le terme \textbf{homininé} pour désigner les fossiles de la lignée humaine (Toumaï, Lucy, Néandertal...).
\end{attention}

\courssub{Démarche d'investigation : reconstituer une chronologie à partir de fiches fossiles}

\begin{methode}[Reconstituer la chronologie de la lignée humaine]
Face à une collection de fiches fossiles (document fréquent au Bac), procède en quatre étapes :
\begin{enumerate}
\item \textbf{Repérer les dates} de chaque fossile (datation radiométrique des couches encaissantes, par exemple par la méthode potassium-argon pour les terrains volcaniques d'Afrique de l'Est).
\item \textbf{Classer les fossiles} du plus ancien au plus récent sur un axe chronologique gradué en millions d'années (Ma).
\item \textbf{Identifier les caractères morphologiques} de chaque espèce : position du trou occipital (bipédie), volume crânien, prognathisme, outils associés, et faire correspondre ces caractères aux critères de l'hominisation.
\item \textbf{Croiser dates et caractères} : tracer les barres d'existence de chaque espèce et repérer les \emph{périodes de coexistence}, qui interdisent toute interprétation linéaire.
\end{enumerate}
\end{methode}

Appliquons cette démarche aux six groupes majeurs de la lignée humaine.

\courssub{Les premiers homininés : \emph{Sahelanthropus} et les Australopithèques}

\textbf{\emph{Sahelanthropus tchadensis}, le plus ancien homininé connu.} En 2001, une équipe franco-tchadienne dirigée par Michel Brunet découvre dans le désert du Djourab (Tchad) un crâne fossile surnommé \cle{Toumaï} (« espoir de vie » en langue gorane), daté d'environ \SI{7}{Ma}. Malgré un très petit volume crânien (environ \SI{350}{cm^3}, comparable à celui d'un chimpanzé), Toumaï présente un caractère décisif : son \cle{trou occipital} (orifice par lequel la moelle épinière rejoint le cerveau) est \emph{avancé} vers le centre du crâne, signe que la tête se portait au-dessus de la colonne vertébrale. \emph{Sahelanthropus} était donc probablement déjà, au moins partiellement, \textbf{bipède}. Cette découverte, faite à \SI{2500}{km} à l'ouest de la vallée du Rift, a montré que les premiers homininés n'étaient pas confinés à l'Afrique de l'Est.

\textbf{Les Australopithèques (environ 4 à 2 Ma).} Ce groupe, bien documenté en Afrique de l'Est (Hadar en Éthiopie, Olduvai en Tanzanie) et en Afrique australe, rassemble des homininés \textbf{bipèdes confirmés} mais conservant des aptitudes arboricoles (bras longs, doigts courbés). L'espèce type est \emph{Australopithecus afarensis}, représentée par le célèbre squelette de \cle{Lucy} (Hadar, 1974, environ \SI{3,2}{Ma}), conservé à 40 \%. Les Australopithèques présentent :
\begin{itemize}
\item un petit cerveau (\SI{400}{cm^3} à \SI{550}{cm^3}) ;
\item une face \cle{prognathe} (projetée en avant) et de fortes mâchoires ;
\item une bipédie attestée aussi par les empreintes de pas fossilisées de Laetoli (Tanzanie, \SI{3,6}{Ma}).
\end{itemize}
Ils ne fabriquaient pas d'outils de pierre taillée (ou seulement des outils très rudimentaires chez les formes tardives).

\courssub{Le genre \emph{Homo} : outils, feu et conquête du monde}

\textbf{\emph{Homo habilis} (environ 2,4 à 1,6 Ma).} Découvert dans les gorges d'Olduvai (Tanzanie) par Louis et Mary Leakey, c'est le \cle{premier fabricant d'outils} : les \emph{galets taillés} (industrie oldowayenne), obtenus en frappant un galet pour en détacher des éclats tranchants. Son volume crânien atteint environ \SI{650}{cm^3} et sa face est moins prognathe. Le nom \emph{habilis} (« habile ») souligne cette innovation technique majeure : la culture matérielle devient un critère d'hominisation.

\textbf{\emph{Homo erectus} (environ 1,9 Ma à 100 000 ans).} Avec un volume crânien de \SI{900}{} à \SI{1100}{cm^3}, un corps de proportions modernes et une bipédie pleinement efficace, \emph{Homo erectus} est le \cle{premier homininé à sortir d'Afrique} : ses fossiles se retrouvent de Géorgie (Dmanissi, \SI{1,8}{Ma}) à la Chine et à Java. Il fabrique des \cle{bifaces} (outils symétriques taillés sur deux faces, industrie acheuléenne) et \textbf{maîtrise le feu} (foyers structurés attestés vers 400 000 ans, par exemple à Zhoukoudien en Chine ou Beeches Pit en Angleterre ; des indices plus anciens existent mais sont discutés), ce qui transforme l'alimentation, la protection et la vie sociale.

\textbf{\emph{Homo neanderthalensis} (environ 300 000 à 30 000 ans, Europe et Proche-Orient).} Adapté aux climats froids (corps trapu, grande capacité crânienne de \SI{1200}{} à \SI{1600}{cm^3}), Néandertal élabore un outillage fin (moustérien), pratique la chasse collective et \textbf{inhume ses morts} : les sépultures constituent le premier indice solide d'une pensée symbolique. L'analyse de l'\cle{ADN fossile} a établi qu'il s'est partiellement hybridé avec \emph{Homo sapiens} à la sortie d'Afrique de ce dernier.

\textbf{\emph{Homo sapiens} (environ 300 000 ans à aujourd'hui).} Les plus anciens fossiles de notre espèce proviennent de \cle{Jebel Irhoud} (Maroc, environ 300 000 ans), ce qui repousse de 100 000 ans l'âge antérieurement admis. \emph{Homo sapiens} se caractérise par un front vertical, un menton saillant, un volume crânien moyen de \SI{1350}{cm^3}, et surtout une \textbf{culture symbolique} explosive : art pariétal, parures, sépultures ritualisées, langage articulé complexe. À partir d'environ 60 000 ans, il quitte définitivement l'Afrique et peuple toute la planète.

\begin{propriete}[Origine africaine de l'Humanité (admise)]
L'Humanité est \cle{originaire d'Afrique}. Cette propriété repose sur deux faits convergents : (1) les plus anciens fossiles d'homininés (Toumaï, Australopithèques, premiers \emph{Homo}) sont tous africains ; (2) la \emph{diversité génétique} des populations humaines actuelles est maximale en Afrique, signe que c'est là que l'espèce a évolué le plus longtemps. (Démonstration non exigible ; arguments à savoir mobiliser.)
\end{propriete}

\courssub{Une histoire en buisson, pas en échelle}

\begin{attention}[L'erreur classique à bannir]
Il est \textbf{faux} de représenter l'évolution humaine comme une échelle linéaire « du singe à l'Homme » (Toumaï $\rightarrow$ Australopithèque $\rightarrow$ \emph{habilis} $\rightarrow$ \emph{erectus} $\rightarrow$ Néandertal $\rightarrow$ \emph{sapiens}). En réalité, \textbf{plusieurs espèces ont coexisté} pendant de longues périodes : \emph{Homo habilis} a côtoyé les derniers Australopithèques ; \emph{Homo erectus} a vécu en même temps que les premiers \emph{Homo sapiens} ; Néandertal et \emph{sapiens} ont partagé l'Europe pendant plus de 10 000 ans. La lignée humaine est un \cle{buisson évolutif} dont presque tous les rameaux se sont éteints ; \emph{Homo sapiens} est le seul rameau survivant. De plus, le chimpanzé actuel n'est pas notre ancêtre : nous partageons avec lui un \emph{ancêtre commun} disparu.
\end{attention}

\begin{popfigure}
\begin{center}
\begin{tikzpicture}[x=1.55cm,y=0.52cm]
% Axe du temps
\draw[->,thick,popdark] (-7.4,0) -- (0.4,0) node[below left=2pt and -30pt,font=\small] {};
\foreach \x/\lab in {-7/{-7 Ma},-6/{-6},-5/{-5},-4/{-4},-3/{-3},-2/{-2},-1/{-1},0/{0}}
  \draw[popdark] (\x,0.12) -- (\x,-0.12) node[below,font=\scriptsize] {\lab};
% Barres d'existence
\draw[line width=5pt,popgold]   (-7.0,1) -- (-6.6,1);
\node[anchor=west,font=\scriptsize,popdark] at (-7.35,1.65) {\emph{Sahelanthropus tchadensis} (Toumaï)};
\draw[line width=5pt,poporange] (-4.2,2) -- (-2.0,2);
\node[anchor=west,font=\scriptsize,popdark] at (-6.6,2.65) {Australopithèques (dont \emph{A. afarensis}, Lucy)};
\draw[line width=5pt,popgreen]  (-2.4,3) -- (-1.6,3);
\node[anchor=west,font=\scriptsize,popdark] at (-5.4,3.65) {\emph{Homo habilis}};
\draw[line width=5pt,popblue]   (-1.9,4) -- (-0.1,4);
\node[anchor=west,font=\scriptsize,popdark] at (-4.6,4.65) {\emph{Homo erectus}};
\draw[line width=5pt,poppurple] (-0.3,5) -- (-0.03,5);
\node[anchor=west,font=\scriptsize,popdark] at (-3.4,5.65) {\emph{Homo neanderthalensis}};
\draw[line width=5pt,poppink]   (-0.3,6) -- (0,6);
\node[anchor=west,font=\scriptsize,popdark] at (-2.2,6.65) {\emph{Homo sapiens}};
% Zones de coexistence
\draw[<->,popmuted,thick] (-2.4,-0.9) -- (-1.6,-0.9);
\node[font=\scriptsize,popmuted] at (-2.0,-1.5) {coexistence \emph{habilis} / Australopithèques};
\draw[<->,popmuted,thick] (-0.3,-2.3) -- (-0.03,-2.3);
\node[font=\scriptsize,popmuted,anchor=west] at (0.1,-2.3) {coexistence Néandertal / \emph{sapiens}};
\end{tikzpicture}
\end{center}
\legende{Frise chronologique de la lignée humaine (échelle en millions d'années, non linéaire pour les périodes récentes). Chaque barre colorée représente la durée d'existence d'un groupe d'homininés ; les doubles flèches grises soulignent deux périodes de coexistence, preuve du caractère buissonnant de l'évolution humaine.}
\end{popfigure}

\begin{popfigure}
\begin{center}
\begin{tikzpicture}[scale=1.05]
% Contour simplifié de l'Afrique
\draw[thick,popdark,fill=popcream]
  (-2.6,2.6) .. controls (-3.0,2.0) and (-2.9,1.4) .. (-2.4,0.9)
  .. controls (-2.0,0.5) and (-2.1,-0.2) .. (-1.7,-0.9)
  .. controls (-1.3,-1.7) and (-0.7,-2.3) .. (0.0,-2.9)
  .. controls (0.5,-3.3) and (1.0,-3.2) .. (1.3,-2.6)
  .. controls (1.7,-1.9) and (1.6,-1.1) .. (1.9,-0.4)
  .. controls (2.3,0.4) and (2.9,0.7) .. (3.1,1.3)
  .. controls (3.3,1.9) and (2.6,2.3) .. (1.9,2.5)
  .. controls (1.0,2.8) and (0.0,2.7) .. (-0.9,2.75)
  .. controls (-1.7,2.8) and (-2.2,2.9) .. (-2.6,2.6) -- cycle;
% Madagascar
\draw[thick,popdark,fill=popcream] (2.35,-1.9) ellipse (0.14 and 0.45);
% Sites
\fill[poppink] (-0.6,2.55) circle (0.07) node[above,font=\scriptsize,popdark] {Jebel Irhoud (Maroc) : \emph{H. sapiens}, 300 000 ans};
\fill[popgold] (-1.9,1.5) circle (0.07) node[left,font=\scriptsize,popdark] {Djourab (Tchad) : Toumaï, 7 Ma};
\fill[poporange] (1.55,0.9) circle (0.07) node[right,font=\scriptsize,popdark] {Hadar (Éthiopie) : Lucy, 3,2 Ma};
\fill[popgreen] (1.5,0.1) circle (0.07) node[right,font=\scriptsize,popdark] {Olduvai / Laetoli (Tanzanie)};
\fill[popblue] (0.9,-2.6) circle (0.07) node[below,font=\scriptsize,popdark] {Afrique australe : Australopithèques};
\end{tikzpicture}
\end{center}
\legende{Carte simplifiée de l'Afrique localisant les principaux sites fossiles de la lignée humaine. Tous les homininés anciens (plus de 2 Ma) sont africains : c'est l'argument paléontologique majeur de l'origine africaine de l'Humanité.}
\end{popfigure}

\begin{savaistu}[Un héritage néandertalien en nous]
Le séquençage de l'ADN extrait d'os fossiles de Néandertal (prix Nobel de médecine 2022 attribué à Svante Pääbo) a montré qu'environ \textbf{1 à 4 \%} du génome des humains actuels non africains provient de l'Homme de Néandertal, hérité de métissages survenus au Proche-Orient il y a environ 50 000 à 60 000 ans. Certains de ces gènes hérités influencent encore notre immunité. L'« extinction » de Néandertal n'a donc pas été totale : une partie de lui vit en nous.
\end{savaistu}

\begin{exoresolu}[Reconstituer une chronologie]
\textbf{Énoncé.} On dispose de quatre fiches fossiles : (A) crâne de volume \SI{650}{cm^3} associé à des galets taillés, Olduvai, daté de \SI{1,8}{Ma} ; (B) squelette féminin d'Australopithèque, Hadar, \SI{3,2}{Ma} ; (C) crâne de \SI{1400}{cm^3} découvert dans une sépulture en Europe, daté de \SI{60000}{ans} ; (D) crâne au trou occipital avancé, Djourab, \SI{7}{Ma}. 1) Attribue chaque fiche à une espèce. 2) Range-les par ordre chronologique. 3) Montre que deux de ces espèces ont pu coexister avec d'autres homininés.
\tcblower
\textbf{Solution.} 1) (A) : \emph{Homo habilis} (galets taillés, \SI{650}{cm^3}) ; (B) : \emph{Australopithecus afarensis} (Lucy) ; (C) : \emph{Homo neanderthalensis} (sépulture, Europe, grand crâne) ; (D) : \emph{Sahelanthropus tchadensis} (Toumaï). 2) Ordre chronologique : D (\SI{7}{Ma}) $\rightarrow$ B (\SI{3,2}{Ma}) $\rightarrow$ A (\SI{1,8}{Ma}) $\rightarrow$ C (60 000 ans). 3) \emph{Homo habilis} (A, 2,4--1,6 Ma) a coexisté avec les derniers Australopithèques (jusqu'à 2 Ma) ; \emph{Homo neanderthalensis} (C) a coexisté en Europe avec \emph{Homo sapiens} entre environ 45 000 et 30 000 ans. Conclusion : la chronologie ne forme pas une succession simple d'espèces mais un buisson évolutif avec des rameaux contemporains, dont la plupart se sont éteints.
\end{exoresolu}

\begin{aretenir}[L'essentiel de la section]
\begin{itemize}
\item La lignée humaine compte six groupes majeurs : \emph{Sahelanthropus} (7 Ma, bipédie probable), Australopithèques (4--2 Ma, bipèdes à petit cerveau), \emph{Homo habilis} (outils, \SI{650}{cm^3}), \emph{Homo erectus} (sortie d'Afrique, feu, bifaces), \emph{Homo neanderthalensis} (sépultures, hybridation avec \emph{sapiens}) et \emph{Homo sapiens} (300 000 ans, art et expansion mondiale).
\item Les critères de l'hominisation sont apparus \emph{progressivement} : la bipédie précède de plusieurs millions d'années l'augmentation du volume crânien.
\item L'évolution humaine est un \textbf{buisson} : plusieurs espèces ont coexisté, et toutes sauf une se sont éteintes.
\item Tous les fossiles anciens sont africains : l'Humanité est née en Afrique.
\end{itemize}
\end{aretenir}

Cette chronologie en buisson pose naturellement une question : comment ces espèces sont-elles reliées entre elles et aux grands singes actuels ? C'est l'objet de la section suivante, consacrée à la construction de l'arbre phylogénétique de la lignée humaine.

\coursec{IV. Parenté entre l'Homme et les grands singes anthropomorphes : l'arbre phylogénétique}

Après avoir reconstitué les grandes étapes de la lignée humaine à travers les fossiles (section III), nous disposons maintenant des outils pour situer notre espèce parmi les primates actuels. Comment classer l'Homme par rapport au chimpanzé, au gorille ou à l'orang-outan ? Les fossiles nous donnent l'histoire, mais la biologie moléculaire et la cytogénétique nous révèlent la parenté profonde. C'est ce que nous allons établir rigoureusement.

\courssub{Les grands singes anthropomorphes et les indices de parenté}

Les \cle{grands singes anthropomorphes} (famille des \textbf{Hominidés}) regroupent quatre genres actuels : \emph{Pan} (chimpanzé \emph{Pan troglodytes} et bonobo \emph{Pan paniscus}), \emph{Gorilla} (gorille de l'Ouest \emph{Gorilla gorilla} et de l'Est \emph{Gorilla beringei}), \emph{Pongo} (orang-outan de Bornéo \emph{Pongo pygmaeus} et de Sumatra \emph{Pongo abelii}) et \emph{Homo} (l'Homme, \emph{Homo sapiens}). Le gibbon (famille des \textbf{Hylobatidés}), bien que singe sans queue, en est exclu car plus éloigné (petit singe ou \emph{petit hominoïde}). L'ensemble forme la super-famille des \textbf{Hominoïdes} (ou Hominoidea).

\begin{definition}[Caractères partagés avec les grands singes (Synapomorphies des Hominidés)]
L'Homme et les grands singes anthropomorphes partagent des \cle{caractères dérivés communs} (synapomorphies) absents chez les singes à queue (Cercopithécidés) :
\begin{itemize}
    \item Absence de queue externe (caudectomie embryonnaire, vertèbres coccygiennes réduites).
    \item Morphologie du bassin élargi et raccourci (ilions évasés), colonne vertébrale lombaire raccourcie, adaptation à la station \textbf{orthograde} (buste droit) et à la brachiation (sauf homme).
    \item Main préhensile avec pouce opposable (réduit chez l'orang-outan), ongles plats à la place de griffes.
    \item Encéphale volumineux par rapport au corps (indice d'encéphalisation élevé), cortex cérébral très plissé (gyrification).
    \item Denture : arcade dentaire en parabole (sauf gorille/orang-outan plus en U), absence de diastème (sauf mâles gorille/orang-outan), canines réduites et non émoulues (sauf dimorphisme sexuel fort chez gorille/orang-outan), molaires inférieures à 5 cuspides (patron \textbf{Y-5}).
    \item Embryologie : développement placentaire \textbf{hémochorial} identique, stades embryonnaires précoces quasi superposables.
\end{itemize}
\end{definition}

\begin{experience}[Observation comparative : caryotypes de l'Homme et du Chimpanzé]
\textbf{Matériel :} Photographies de caryotypes en bandes G (résolution 400-550 bandes) d'un homme (46,XY) et d'un chimpanzé (48,XY) au même grossissement.
\textbf{Protocole :} Compter les chromosomes. Identifier les chromosomes homologues par leur taille, position du centromère et profil de bandes. Rechercher les correspondances bande par bande.
\textbf{Observations :}
\begin{itemize}
    \item L'Homme possède \textbf{46 chromosomes} (23 paires). Le chimpanzé (comme le gorille et l'orang-outan) en possède \textbf{48} (24 paires).
    \item Les bandes des chromosomes 1, 3, 4, 5, 6, 7, 8, 9, 10, 11, 12, 13, 14, 15, 16, 17, 18, 19, 20, 21, 22, X et Y sont \emph{quasi identiques} entre les deux espèces (homologie de bandage conservée).
    \item \textbf{Le chromosome 2 humain} (métacentrique) correspond \emph{exactement} à la fusion bout-à-bout (fusion \textbf{télomère-télomère}) de \textbf{deux chromosomes acrocentriques distincts} chez le chimpanzé (nommés 2A et 2B, orthologues aux bras 2p et 2q humains).
\end{itemize}
\textbf{Interprétation :} L'ancêtre commun Homme/Chimpanzé possédait 48 chromosomes. Une mutation de fusion chromosomique (translocation télo-télique) s'est produite dans la lignée menant à l'Homme \emph{après} la divergence avec le chimpanzé (il y a \textasciitilde 6-7 Ma), réduisant le nombre à 46. C'est une \textbf{preuve cytogénétique majeure} de notre parenté étroite.
\legende{Comparaison schématique des caryotypes : le chromosome 2 humain résulte de la fusion de deux chromosomes ancestraux (2A et 2B) restés séparés chez le chimpanzé. Les bandes G sont conservées.}
\end{experience}

\begin{popfigure}
\begin{tikzpicture}[scale=0.9, every node/.style={font=\small}]
    % Human Chr 2 (metacentric)
    \draw[popblue, very thick] (0,0) -- (0,5);
    \draw[popblue, very thick] (0.35,0) -- (0.35,5);
    % Centromere human (metacentric ~ middle)
    \fill[popblue!70!popdark] (0,2.4) rectangle (0.35,2.7);
    % Bands human (simplified pattern)
    \foreach \y/\h in {0.2/0.3, 0.6/0.15, 1.0/0.3, 1.4/0.15, 1.8/0.3, 2.8/0.3, 3.2/0.15, 3.6/0.3, 4.0/0.15, 4.4/0.3}
        \fill[popblue!30] (0,\y) rectangle (0.35,\y+\h);
    \node[align=center] at (0.175, -0.6) {Chr 2 Humain\\(métacentrique, 46 chr)};
    \node[popblue, rotate=90, anchor=south] at (-0.6, 2.55) {Centromère fonctionnel};
    \node[popred, rotate=90, anchor=south] at (-0.6, 4.7) {Télomères};
    \node[popred, rotate=90, anchor=south] at (-0.6, 0.3) {Télomères};
    % Residual telomeres at fusion point (2q13)
    \node[popred, font=\bfseries] at (0.175, 2.55) {$\ast$};
    \node[font=\tiny, align=center, text=popred] at (0.175, 2.0) {Séquences\\télomériques\\résiduelles\\(TTAGGG)n};
    % Inactive centromere (2q21)
    \node[poppurple, font=\bfseries] at (0.175, 3.6) {$\bullet$};
    \node[font=\tiny, align=center, text=poppurple] at (0.175, 3.2) {Centromère\\inactivé (vestige)};

    % Ancestral 2A (Chimp chr 2A ~ human 2p)
    \draw[poporange, very thick] (2.5,0) -- (2.5,2.5);
    \draw[poporange, very thick] (2.85,0) -- (2.85,2.5);
    \fill[poporange!70!popdark] (2.5,0.1) rectangle (2.85,0.4); % Acrocentric centromere near top
    \foreach \y/\h in {0.6/0.3, 1.0/0.15, 1.4/0.3, 1.8/0.15, 2.2/0.3}
        \fill[poporange!30] (2.5,\y) rectangle (2.85,\y+\h);
    \node[align=center] at (2.675, -0.6) {Chr 2A (Chimpanzé)\\(acrocentrique $\sim$ 2p)};
    \node[poporange, rotate=90, anchor=south] at (2.0, 0.25) {Centromère};
    \node[popred, rotate=90, anchor=south] at (2.0, 2.4) {Télomère};

    % Ancestral 2B (Chimp chr 2B ~ human 2q)
    \draw[popgreen, very thick] (4,0) -- (4,2.5);
    \draw[popgreen, very thick] (4.35,0) -- (4.35,2.5);
    \fill[popgreen!70!popdark] (4,0.1) rectangle (4.35,0.4); % Acrocentric centromere near top
    \foreach \y/\h in {0.6/0.3, 1.0/0.15, 1.4/0.3, 1.8/0.15, 2.2/0.3}
        \fill[popgreen!30] (4,\y) rectangle (4.35,\y+\h);
    \node[align=center] at (4.175, -0.6) {Chr 2B (Chimpanzé)\\(acrocentrique $\sim$ 2q)};
    \node[popgreen, rotate=90, anchor=south] at (3.5, 0.25) {Centromère};
    \node[popred, rotate=90, anchor=south] at (3.5, 2.4) {Télomère};

    % Arrow fusion
    \draw[->, thick, popdark] (0.6, 2.55) -- (1.9, 2.55) node[midway, above] {Fusion télomère-télomère};
    \draw[->, thick, popdark] (0.6, 2.3) -- (1.9, 1.8);
\end{tikzpicture}
\legende{Schéma de la fusion chromosomique à l'origine du chromosome 2 humain. Les bandes de coloration (bandes G) sont conservées entre les deux chromosomes ancestraux (2A, 2B) et le chromosome 2 fusionné. La présence de séquences télomériques résiduelles (TTAGGG)n au point de fusion (2q13, $\ast$ rouge) et d'un centromère inactivé vestigial (2q21, $\bullet$ violet) confirme l'hypothèse.}
\end{popfigure}

\courssub{La preuve moléculaire décisive : comparaison des séquences d'ADN et des protéines}

Au-delà de l'anatomie et des chromosomes, la comparaison des \cle{séquences nucléotidiques} (ADN) et des \cle{séquences amino-acidiques} (protéines) fournit une mesure quantitative de la parenté.

\begin{propriete}[Distance génétique et parenté]
Le nombre de différences de nucléotides (ou d'acides aminés) entre deux espèces pour un gène (ou une protéine) orthologue est proportionnel au temps écoulé depuis leur divergence (principe de l'horloge moléculaire). \textbf{Moins il y a de différences, plus la parenté est proche.}
\end{propriete}

\begin{exemplebox}[Exemple chiffré : Protéine $\beta$-hémoglobine (146 acides aminés)]
\begin{center}
\begin{tabularx}{\linewidth}{|l|X|X|}\hline
\rowcolor{popblueL}
\textbf{Comparaison} & \textbf{Nombre de différences (AA)} & \textbf{Pourcentage de divergence} \\ \hline
Homme -- Chimpanzé & 1 & $\approx 0,7\%$ \\ \hline
Homme -- Gorille & 2 & $\approx 1,4\%$ \\ \hline
Homme -- Orang-outan & 8 & $\approx 5,5\%$ \\ \hline
Homme -- Gibbon & 15 & $\approx 10,3\%$ \\ \hline
Homme -- Macaque (Cercopithèque) & 18 & $\approx 12,3\%$ \\ \hline
\end{tabularx}
\end{center}
\legende{Données classiques (Zuckerkandl \& Pauling, 1965 ; confirmées par génomique moderne). La divergence Homme/Chimpanzé est la plus faible.}
\end{exemplebox}

\begin{aretenir}[Résultat génomique global]
Le séquençage complet des génomes confirme : l'Homme et le chimpanzé (et le bonobo) partagent \textbf{$\approx 98,8\%$ d'identité nucléotidique} (soit $\approx 1,2\%$ de différences sur l'alignement global), contre $\approx 98,4\%$ avec le gorille ($\approx 1,6\%$ de différences). \textbf{Le chimpanzé (et le bonobo) est le plus proche parent vivant de l'Homme.}
\end{aretenir}

\courssub{Principe de construction et lecture de l'arbre phylogénétique}

\begin{definition}[Arbre phylogénétique (Chronogramme)]
Un \cle{arbre phylogénétique} est un graphe acyclique représentant les relations de parenté entre des taxons (espèces, genres). 
\begin{itemize}
    \item Les \textbf{feuilles} (extrémités) représentent les espèces \textbf{actuelles} (ou fossiles terminales). Elles sont toutes alignées au temps présent (0 Ma).
    \item Les \textbf{nœuds} (points de branchement) représentent les \textbf{derniers ancêtres communs} (DAC) : des espèces \textbf{disparues}.
    \item Les \textbf{branches} représentent les lignées évolutives.
    \item Dans un \textbf{chronogramme}, l'axe horizontal (ou vertical) est proportionnel au \textbf{temps géologique} (Ma). La distance entre deux feuilles = somme des longueurs de branches = divergence cumulée.
\end{itemize}
\end{definition}

\begin{definition}[Ancêtre commun et Clade]
Un \cle{ancêtre commun} d'un groupe d'espèces est une espèce disparue dont descendent toutes les espèces de ce groupe. Le \emph{dernier ancêtre commun} (DAC) de deux espèces est le nœud le plus récent qui les relie. Un \textbf{clade} (ou groupe monophylétique) regroupe un ancêtre commun et \textbf{tous} ses descendants.
\end{definition}

\begin{attention}[Piège majeur : Lecture de l'arbre]
\textbf{Ne jamais lire un arbre phylogénétique de gauche à droite (ou bas en haut) comme une "progression" ou une "échelle de perfection".}
\begin{itemize}
    \item L'ordre des branches aux nœuds (gauche/droite) est \textbf{arbitraire} (on peut pivoter les branches sans changer l'information phylogénétique).
    \item Seule la \textbf{topologie} (qui est regroupé avec qui) compte : les espèces sœurs partagent un nœud exclusif.
    \item L'Homme n'est pas "au sommet" ni "le but" de l'évolution ; il est une feuille parmi d'autres, \textbf{cousin} du chimpanzé et du bonobo (groupe frère), \textbf{neveu} du gorille, etc.
    \item Toutes les espèces actuelles ont la même \textbf{profondeur évolutive} (même temps écoulé depuis la racine).
\end{itemize}
\end{attention}

\begin{definition}[Horloge moléculaire (savoir admis)]
L'\cle{horloge moléculaire} est une méthode estimant la date de divergence entre deux lignées à partir du nombre de mutations neutres accumulées dans leur ADN, en supposant un \textbf{taux de mutation constant} dans le temps pour une molécule donnée. 
\[
T_{\text{divergence}} = \frac{\text{Nombre de différences observées (par site)}}{2 \times \text{Taux de mutation (substitutions/site/an)}}
\]
(Le facteur 2 tient compte de l'accumulation indépendante dans les deux lignées depuis la séparation). Le taux est calibré sur des datations fossiles fiables (points de calibration).
\end{definition}

\begin{popfigure}
\begin{tikzpicture}[scale=0.75, every node/.style={font=\small}]
    % TIME AXIS (Horizontal: Past -> Present)
    % Root at x=0 (approx 15-16 Ma for Hominidae crown), Tips at x=12 (0 Ma)
    % 1 Ma = 0.8 units approx. 15 Ma = 12 units.
    \draw[->, popmuted, thick] (-0.5, -4.5) -- (12.5, -4.5) node[right, popmuted] {Temps (Ma avant présent)};
    \foreach \x/\t in {0/15, 2/13.5, 4/10.5, 6/7.5, 8/4.5, 10/1.5, 12/0}
        \draw[popmuted] (\x, -4.6) -- (\x, -4.4) node[below, popmuted] {\t};
    
    % NODES (x = time, y = topology)
    % Root: Hominidae crown (Hominidae + Hylobatidae split ~18-20 Ma, but we show Hominidae crown ~14-15 Ma)
    \coordinate (root) at (0, 0); % ~15 Ma
    % Node 1: Orangutan split (Ponginae vs Homininae) ~14 Ma
    \coordinate (n_orang) at (1.5, 3.5); 
    % Node 2: African Apes ancestor (Homininae) ~14 Ma
    \coordinate (n_afr) at (1.5, -3.5);
    % Node 3: Gorilla split (Gorillini vs Hominini) ~9-10 Ma
    \coordinate (n_gor) at (5.5, -3.5);
    % Node 4: Human-Chimp ancestor (Hominini) ~9-10 Ma
    \coordinate (n_hc) at (5.5, 1.0);
    % Node 5: Human split (Hominina) ~6.5-7 Ma
    \coordinate (n_hum) at (8.5, 2.5);
    % Node 6: Pan split (Panina: Chimp vs Bonobo) ~1.5-2 Ma
    \coordinate (n_pan) at (10.5, -0.5);
    
    % LEAVES (All at x=12, Present)
    \coordinate (l_orang) at (12, 3.5);
    \coordinate (l_gor) at (12, -3.5);
    \coordinate (l_hum) at (12, 2.5);
    \coordinate (l_bon) at (12, 0.5);
    \coordinate (l_chimp) at (12, -1.5);
    
    % BRANCHES
    \draw[popdark, thick] (root) -- (n_orang);
    \draw[popdark, thick] (root) -- (n_afr);
    \draw[popdark, thick] (n_afr) -- (n_gor);
    \draw[popdark, thick] (n_afr) -- (n_hc);
    \draw[popdark, thick] (n_hc) -- (n_hum);
    \draw[popdark, thick] (n_hc) -- (n_pan);
    \draw[popdark, thick] (n_pan) -- (l_bon);
    \draw[popdark, thick] (n_pan) -- (l_chimp);
    \draw[popdark, thick] (n_hum) -- (l_hum);
    \draw[popdark, thick] (n_gor) -- (l_gor);
    \draw[popdark, thick] (n_orang) -- (l_orang);
    
    % NODE MARKERS (circles)
    \foreach \c in {root, n_orang, n_afr, n_gor, n_hc, n_hum, n_pan}
        \fill[popdark] (\c) circle (2.5pt);
    
    % DIVERGENCE DATES at nodes
    \node[popblue, anchor=west, font=\small\bfseries] at (1.6, 3.5) {$\sim$14--15 Ma};
    \node[popblue, anchor=west, font=\small\bfseries] at (5.6, -3.5) {$\sim$9--10 Ma};
    \node[popblue, anchor=west, font=\small\bfseries] at (5.6, 1.0) {$\sim$9--10 Ma};
    \node[popblue, anchor=west, font=\small\bfseries] at (8.6, 2.5) {$\sim$6--7 Ma};
    \node[popblue, anchor=west, font=\small\bfseries] at (10.6, -0.5) {$\sim$1,5--2 Ma};
    
    % LEAF LABELS
    \node[anchor=west, poporange, font=\bfseries] at (l_orang) {Orang-outan (\emph{Pongo})};
    \node[anchor=west, popgreen, font=\bfseries] at (l_gor) {Gorille (\emph{Gorilla})};
    \node[anchor=west, popblue, font=\bfseries] at (l_hum) {Homme (\emph{Homo sapiens})};
    \node[anchor=west, poppurple, font=\bfseries] at (l_bon) {Bonobo (\emph{Pan paniscus})};
    \node[anchor=west, poppurple, font=\bfseries] at (l_chimp) {Chimpanzé (\emph{Pan troglodytes})};
    
    % CLADE LABELS (vertical brackets or text)
    \node[rotate=90, anchor=south, popmuted, font=\small] at (-0.8, 0) {Hominidés (Grands singes)};
    \node[rotate=90, anchor=south, popmuted, font=\small] at (3.5, -1.0) {Homininés (Africains / Homininés s.l.)};
    \node[rotate=90, anchor=south, popmuted, font=\small] at (7.0, 1.75) {Homininés (s.s. / Lignée humaine)};
    \node[rotate=90, anchor=south, popmuted, font=\small] at (9.5, -0.5) {Panins};
    
    % Gibbon outgroup (optional, dashed line to root)
    \draw[popmuted, dashed] (root) -- (-1, 0) node[left, popmuted] {Gibbon (Hylobatidés) $\sim$18-20 Ma};
\end{tikzpicture}
\legende{Chronogramme phylogénétique des Hominoïdes (Hominidés + Hylobatidés). L'axe horizontal représente le temps (Ma). Les nœuds (cercles noirs) sont les derniers ancêtres communs datés par horloge moléculaire calibrée. Topologie : \emph{((Homme, (Chimpanzé, Bonobo)), Gorille), Orang-outan), Gibbon}. L'Homme et Pan (Chimpanzé/Bonobo) sont groupes frères (clade Hominini).}
\end{popfigure}

\courssub{Construction guidée d'un arbre phylogénétique à partir d'une matrice de distances}

C'est une compétence clé du programme : savoir passer d'un tableau de différences (matrice de distances) à l'arbre de parenté. On utilise la méthode UPGMA (\emph{Unweighted Pair Group Method with Arithmetic Mean}), qui suppose une horloge moléculaire (arbre ultramétrique).

\begin{methode}[Construire un arbre phylogénétique par regroupement successif (UPGMA simplifiée)]
\begin{enumerate}
    \item \textbf{Identifier la plus petite distance} dans le tableau (hors diagonale). Les deux espèces correspondantes sont \textbf{sœurs} : on les relie à un nœud commun.
    \item \textbf{Calculer la hauteur du nœud} : $h = d_{\min} / 2$ (distance depuis les feuilles jusqu'au nœud).
    \item \textbf{Calculer les distances du nouveau groupe} (cluster) aux autres espèces : faire la \textbf{moyenne arithmétique} des distances de chaque membre du groupe aux espèces extérieures.
    \item \textbf{Remplacer} les deux espèces par le groupe dans le tableau (réduire la matrice).
    \item \textbf{Répéter} les étapes 1 à 4 jusqu'à ce qu'il ne reste qu'un seul groupe (la racine).
    \item \textbf{Tracer l'arbre} : positionner les nœuds à leurs hauteurs respectives ; la longueur de branche = différence de hauteur entre nœud parent et nœud enfant.
\end{enumerate}
\end{methode}

\begin{exoresolu}[Construction de l'arbre des primates à partir de distances génomiques]
\textbf{Énoncé :} On donne le pourcentage de différences nucléotidiques (génome complet aligné) entre 5 primates.
\begin{center}
\begin{tabular}{|c|ccccc|}\hline
 & Homme & Chimpanzé & Gorille & Orang-outan & Gibbon \\ \hline
Homme & 0 & \textbf{1,2\%} & 1,6\% & 3,1\% & 5,0\% \\
Chimpanzé & 1,2\% & 0 & 1,7\% & 3,2\% & 5,1\% \\
Gorille & 1,6\% & 1,7\% & 0 & 3,0\% & 4,9\% \\
Orang-outan & 3,1\% & 3,2\% & 3,0\% & 0 & 4,5\% \\
Gibbon & 5,0\% & 5,1\% & 4,9\% & 4,5\% & 0 \\ \hline
\end{tabular}
\end{center}
Construire l'arbre phylogénétique correspondant (topologie et hauteurs des nœuds).

\tcblower
\textbf{Solution détaillée :}

\textbf{Étape 1 :} Plus petite distance = \textbf{1,2\% (Homme -- Chimpanzé)}. 
$\rightarrow$ Regroupement \textbf{[Homme, Chimpanzé]} (Nœud N1). Hauteur N1 = $1,2 / 2 = \mathbf{0,60\%}$.

\textbf{Étape 2 :} Calcul distances du groupe [Homme, Chimpanzé] aux autres (moyennes) :
\begin{itemize}
    \item vers Gorille : $(1,6 + 1,7) / 2 = \mathbf{1,65\%}$
    \item vers Orang-outan : $(3,1 + 3,2) / 2 = \mathbf{3,15\%}$
    \item vers Gibbon : $(5,0 + 5,1) / 2 = \mathbf{5,05\%}$
\end{itemize}
Nouvelle matrice réduite :
\begin{center}
\begin{tabular}{|c|cccc|}\hline
 & [Homme,Chimp] & Gorille & Orang-outan & Gibbon \\ \hline
[Homme,Chimp] & 0 & \textbf{1,65\%} & 3,15\% & 5,05\% \\
Gorille & 1,65\% & 0 & 3,0\% & 4,9\% \\
Orang-outan & 3,15\% & 3,0\% & 0 & 4,5\% \\
Gibbon & 5,05\% & 4,9\% & 4,5\% & 0 \\ \hline
\end{tabular}
\end{center}

\textbf{Étape 3 :} Plus petite distance = \textbf{1,65\% ([Homme,Chimp] -- Gorille)}. 
$\rightarrow$ Regroupement \textbf{[[Homme,Chimp], Gorille]} (Nœud N2). Hauteur N2 = $1,65 / 2 = \mathbf{0,825\%}$.
Longueur branche Gorille $\rightarrow$ N2 = $0,825\%$. Longueur branche N1 $\rightarrow$ N2 = $0,825 - 0,600 = \mathbf{0,225\%}$.

\textbf{Étape 4 :} Distances du groupe [[Homme,Chimp], Gorille] aux autres :
\begin{itemize}
    \item vers Orang-outan : $(3,15 + 3,0) / 2 = \mathbf{3,075\%}$
    \item vers Gibbon : $(5,05 + 4,9) / 2 = \mathbf{4,975\%}$
\end{itemize}
Matrice :
\begin{center}
\begin{tabular}{|c|ccc|}\hline
 & [[H,C],G] & Orang-outan & Gibbon \\ \hline
[[H,C],G] & 0 & \textbf{3,075\%} & 4,975\% \\
Orang-outan & 3,075\% & 0 & 4,5\% \\
Gibbon & 4,975\% & 4,5\% & 0 \\ \hline
\end{tabular}
\end{center}

\textbf{Étape 5 :} Plus petite distance = \textbf{3,075\% ([[H,C],G] -- Orang-outan)}. 
$\rightarrow$ Regroupement \textbf{[[[H,C],G], Orang]} (Nœud N3). Hauteur N3 = $3,075 / 2 = \mathbf{1,5375\%}$.
Longueur branche Orang $\rightarrow$ N3 = $1,5375\%$. Longueur branche N2 $\rightarrow$ N3 = $1,5375 - 0,825 = \mathbf{0,7125\%}$.

\textbf{Étape 6 :} Distance finale vers Gibbon : $(4,975 + 4,5) / 2 = \mathbf{4,7375\%}$.
$\rightarrow$ Racine N4 (hauteur $4,7375 / 2 = \mathbf{2,36875\%}$).

\textbf{Arbre final (topologie et hauteurs relatives) :} \emph{((Homme, Chimpanzé), Gorille), Orang-outan), Gibbon}.
Cela correspond exactement à la classification actuelle : Hominini (Homme+Chimpanzé), Homininae (Hominini+Gorille), Hominidae (Homininae+Orang-outan), Hominoidea (Hominidae+Gibbon).
\end{exoresolu}

\courssub{Interprétation évolutive : la fusion chromosomique comme marqueur partagé}

La fusion du chromosome 2 est un \cle{caractère dérivé partagé} (synapomorphie) de \textbf{tous les Homininés} (lignée humaine après séparation d'avec le clade Pan, il y a \textasciitilde 6-7 Ma). Elle est absente chez \emph{Pan}, \emph{Gorilla}, \emph{Pongo}.

\begin{experience}[Démonstration par hybridation moléculaire (FISH -- Fluorescence In Situ Hybridization)]
\textbf{Principe :} On utilise des sondes ADN marquées par fluorescence hybridées sur des chromosomes en métaphase :
\begin{itemize}
    \item Sonde \textbf{télomérique} (répétitions \texttt{TTAGGG}n) : marque les extrémités chromosomiques.
    \item Sonde \textbf{centromérique} (ADN satellite $\alpha$ spécifique) : marque les centromères.
    \item Sonde \textbf{peinture chromosomique} (chromosome 2 humain entier) : marque l'homologie de séquence.
\end{itemize}
\textbf{Résultats observés au microscope à fluorescence :}
\begin{itemize}
    \item \textbf{Homme (46 chr) :} La sonde "peinture chr 2" marque \textbf{un seul} grand chromosome métacentrique. La sonde télomérique marque \textbf{les deux extrémités} + \textbf{un signal interstitiel brillant} en position 2q13 (vestige de la fusion). La sonde centromérique marque \textbf{un seul centromère} fonctionnel (2p11) ; le centromère ancestral 2q21 est inactivé (pas de signal ou signal faible).
    \item \textbf{Chimpanzé/Gorille (48 chr) :} La sonde "peinture chr 2 humain" marque \textbf{deux chromosomes distincts} (acrocentriques). La sonde télomérique marque seulement les quatre extrémités. Deux centromères fonctionnels.
\end{itemize}
\textbf{Conclusion :} L'hybridation \textbf{in situ} (FISH) visualise directement la structure fusionnée du chromosome 2 humain et confirme l'homologie de séquence bande par bande. C'est une preuve \textbf{expérimentale, directe et irréfutable} de l'origine commune des chromosomes et de l'événement de fusion spécifique à la lignée humaine.
\legende{Résultat de FISH : Sonde "peinture chromosome 2 humain" sur métaphase d'homme (1 signal métacentrique) et de chimpanzé (2 signaux acrocentriques). Sonde télomérique : signaux aux extrémités + signal interstitiel chez l'homme (flèche rouge, 2q13).}
\end{experience}

\begin{savaistu}[Le chromosome 2 et la spéciation ?]
La fusion chromosomique réduit la fertilité des hétérozygotes (problèmes de ségrégation en méiose I : formation de trivalents, risque de gamètes aneuploïdes déséquilibrés). Certains évolutionnistes (ex: J. Navarro, N. Barton) proposent que cette fusion a pu contribuer à l'\textbf{isolement reproductif} initial entre la lignée humaine naissante et la lignée chimpanzée, agissant comme une barrière post-zygotique partielle (spéciation chromosomique). C'est une hypothèse stimulante mais débattue (d'autres fusions sont polymorphes sans spéciation), qui illustre le lien entre génomique structurale et macroévolution.
\end{savaistu}

\begin{exercice}[Application : Analyse d'un arbre phylogénétique correct]
On considère l'arbre phylogénétique chronogramme ci-dessous (topologie validée) :
\emph{(Homme, (Chimpanzé, Bonobo)), Gorille), Orang-outan)}.
\begin{enumerate}
    \item Quel est le groupe frère de l'Homme ? Justifier par la topologie.
    \item Quel est le dernier ancêtre commun (DAC) de l'Homme et du Gorille ? Le situer sur l'arbre.
    \item L'affirmation "Le gorille est plus proche de l'Homme que du Chimpanzé" est-elle correcte ? Expliquer en utilisant la notion de DAC.
    \item Si on découvre un fossile daté de 7 Ma présentant des caractères dérivés partagés avec l'Homme et le Chimpanzé (ex: bipédie facultative, canines réduites, foramen magnum antérieur), où le placer sur l'arbre ? Quel nom de clade reçoit-il ?
\end{enumerate}
\end{exercice}

\begin{corrige}[Exercice n° Application arbre phylogénétique]
\begin{enumerate}
    \item \textbf{Le clade Pan (Chimpanzé + Bonobo).} Sur l'arbre, l'Homme et le clade Pan partagent un nœud exclusif (Nœud Hominini). Ils sont groupes frères. Le Bonobo et le Chimpanzé sont groupes frères l'un de l'autre (Nœud Panina), plus récents.
    \item \textbf{Le nœud unissant le clade Hominini (Homme+Pan) au Gorille (Nœud Homininae / Homininés africains).} C'est le nœud le plus récent ancestral à la fois à l'Homme et au Gorille. Il est plus profond (plus ancien, \textasciitilde 9-10 Ma) que le nœud Hominini (\textasciitilde 6-7 Ma).
    \item \textbf{Faux.} \textbf{Le gorille est équidistant de l'Homme et du Chimpanzé (et du Bonobo).} Le DAC(Gorille, Homme) = DAC(Gorille, Chimpanzé) = DAC(Gorille, Bonobo) = Nœud Homininae. La distance phylogénétique (temps de divergence ou nombre de nœuds) est strictement identique. La parenté ne se mesure pas à la "ressemblance" mais au DAC.
    \item \textbf{Sur la branche menant à l'Homme, après la divergence avec le clade Pan (donc $< 6-7$ Ma).} Il appartient au clade des \textbf{Homininés (sens strict / Hominina)} : lignée humaine exclusive postérieure à la séparation d'avec les Panins. C'est un "homininé basal" ou "premier homininé" (ex: \emph{Sahelanthropus}, \emph{Orrorin}).
\end{enumerate}
\end{corrige}

\begin{aretenir}[Synthèse : Preuves de la parenté Homme -- Grands Singes]
\begin{enumerate}
    \item \textbf{Anatomie comparée :} Synapomorphies des Hominidés (pas de queue, bassin orthograde, denture Y-5, encéphale volumineux, pouce opposable).
    \item \textbf{Embryologie :} Stades précoces quasi identiques, placenta hémochorial, phylétisme.
    \item \textbf{Cytogénétique :} Bandes G quasi identiques sur 23 paires ; \textbf{fusion unique du chromosome 2} chez l'Homme (46 chr) vs 48 chez les autres Hominidés. Confirmé par FISH (séquences télomériques interstitielles, centromère vestigial).
    \item \textbf{Moléculaire (ADN/Protéines) :} Moins de différences Homme/Chimpanzé-Bonobo ($\sim 1,2\%$ génome, $\sim 0,7\%$ $\beta$-globine) qu'avec tout autre primate. Horloge moléculaire $\rightarrow$ divergence Homme/Pan $\sim 6-7$ Ma.
    \item \textbf{Phylogénie :} Arbre robuste (chronogramme) : \emph{(Homme, (Chimpanzé, Bonobo)), Gorille), Orang-outan), Gibbon}. L'Homme est un \textbf{grand singe africain} (Hominidé), nested au sein des Homininés (tribu Hominini), groupe frère des Panins.
\end{enumerate}
\end{aretenir}

\coursec{V. Capacité d'adaptation de l'Homme à son environnement et développement durable}

Après avoir reconstitué la lignée humaine et établi notre parenté avec les grands singes anthropomorphes, nous comprenons que l'Homme est le produit d'une longue histoire évolutive. Cette histoire ne s'arrête pas à l'apparition de \emph{Homo sapiens} : notre espèce continue d'évoluer, mais elle a acquis une capacité unique — l'adaptation culturelle et technique — qui transforme radicalement sa relation à l'environnement. Au Cameroun comme partout dans le monde, cette relation est aujourd'hui au cœur des enjeux du développement durable.

\courssub{L'adaptation biologique : l'héritage de la sélection naturelle}

\begin{definition}[Adaptation]
Ensemble des caractères (morphologiques, physiologiques, comportementaux) héréditaires qui augmentent la valeur sélective d'un individu dans un milieu donné, c'est-à-dire sa capacité à survivre et à se reproduire pour transmettre ses gènes à la génération suivante.
\end{definition}

L'adaptation biologique repose sur la \cle{sélection naturelle} : les mutations aléatoires créent de la variabilité génétique ; les pressions du milieu (climat, prédateurs, ressources, pathogènes) favorisent la reproduction différentielle des phénotypes les plus aptes. Ce processus est lent, s'étalant sur des centaines à des milliers de générations.

\begin{experience}[Adaptation humaine à la haute altitude : le cas des populations andines et tibétaines]
\textbf{Observation :} Les populations natives des hauts plateaux andins (\SIrange{4000}{4500}{\meter}) et de l'Himalaya tolèrent l'hypoxie chronique (pression partielle d'\ce{O2} ~ \SI{60}{\%} de celle au niveau de la mer) sans maladie aiguë des montagnes.
\textbf{Hypothèses :} Des mutations avantageuses ont été sélectionnées dans les gènes régulant la respiration, l'hématopoïèse ou la vascularisation.
\textbf{Documents / Résultats :} 
\begin{itemize}
    \item Andes : taux d'hémoglobine élevé (jusqu'à \SIrange{19}{21}{\gram\per\deci\liter} vs \SIrange{13}{15}{\gram\per\deci\liter} au niveau de la mer), volume pulmonaire accru, sélection sur le gène \emph{EPAS1} (régule la réponse à l'hypoxie).
    \item Tibet : taux d'hémoglobine proche de la normale, mais ventilation accrue et meilleure perfusion capillaire ; allèle dérivé de \emph{EPAS1} (introgression dénisovienne) fixé en ~ \num{3000} ans.
\end{itemize}
\textbf{Interprétation :} Deux solutions génétiques différentes (\cle{évolution convergente}) pour un même défi environnemental. La sélection a été forte : l'allèle tibétain de \emph{EPAS1} confère un avantage reproductif majeur (moins de fausses couches, meilleurs poids de naissance).
\textbf{Conclusion :} L'adaptation biologique à un stress environnemental fort peut être rapide (quelques millénaires) mais reste limitée par le temps de génération et l'intensité de la sélection.
\end{experience}

Autres exemples classiques d'adaptation biologique récente chez \emph{Homo sapiens} :
\begin{itemize}
    \item \textbf{Couleur de peau et rayonnement UV :} Gradient latitudinal de la mélanisation (protection contre la dégradation du folate par les UV-B vs synthèse de vitamine D). Au Cameroun, les populations équatoriales présentent une peau fortement mélanisée (phototype V--VI), adaptation à un UV index annuel moyen > 11.
    \item \textbf{\cle{Persistance de la lactase} (LP) :} Maintien de l'expression de la lactase-phlorizine hydrolase à l'âge adulte. Allèles distincts sélectionnés indépendamment en Europe (\emph{-13910*T}), en Afrique de l'Est (\emph{-14010*C}, \emph{-13915*G}) et au Sahel (\emph{-13907*G}) il y a \\ umrange{3000}{7000} ans, corrélés à l'élevage bovin/camelin. Au Nord-Cameroun, les Peuls et Arabes Shuwa montrent une fréquence élevée de LP (~ \SIrange{60}{80}{\%}), contre < \SI{10}{\%} dans les populations forestières du Sud non éleveuses.
    \item \textbf{Résistance au paludisme :} Hétérozygotes pour l'allèle drépanocytaire ($Hb^S$) ou déficients en G6PD, ou porteurs d'hémoglobine C ($Hb^C$) ou d'ovale ($Hb^E$) en zone endémique. Au Cameroun, la fréquence de $Hb^S$ atteint \SIrange{10}{25}{\%} selon les régions, maintenant un \cle{polymorphisme équilibré} (avantage hétérozygote).
\end{itemize}

\begin{attention}[Erreur fréquente : « L'Homme n'évolue plus »]
\textbf{Faux.} L'évolution biologique se poursuit dès qu'il y a variabilité génétique, hérédité et sélection différentielle. La médecine moderne, l'hygiène et l'abondance alimentaire \emph{modifient} les pressions de sélection (relaxation sur la mortalité infantile, nouvelle sélection sur la fécondité, les maladies métaboliques, la résistance aux antibiotiques), mais ne la suppriment pas. Exemple : la fréquence de l'allèle $Hb^S$ diminue lentement dans les populations africaines émigrées vers des zones non paludéennes (désavantage des homozygotes $Hb^S Hb^S$ sans avantage hétérozygote). Au Bac, ne confondez pas « évolution biologique » et « progrès technique ».
\end{attention}

\courssub{La spécificité humaine : l'adaptation culturelle et technique}

\begin{propriete}[Adaptation culturelle vs adaptation biologique]
L'adaptation culturelle humaine est \textbf{cumulative} (\cle{effet de ratchet} : chaque génération hérite des innovations des précédentes et y ajoute les siennes), \textbf{transmise horizontalement et oblique} (apprentissage, enseignement, imitation) bien plus vite que la transmission verticale génétique, et \textbf{réversible} (on peut abandonner une technique). Elle permet une réponse adaptative en une seule génération là où la sélection naturelle exige des millénaires.
\end{propriete}

\begin{exemplebox}[Vêtements, habitat, agriculture, médecine : des boucliers culturels]
\begin{itemize}
    \item \textbf{Vêtements / habitat :} Permettent à l'Homme nu (perte de pilosité il y a ~ \num{1,2} Ma) de coloniser la Sibérie (\SI{-50}{\celsius}) ou le Sahara (\SI{50}{\celsius}) sans attendre de mutations pour la fourrure ou la sudation.
    \item \textbf{Agriculture / élevage :} Domestication (~ \num{10000} ans) = coévolution homme-plantes/animaux. Sélection artificielle sur les espèces cultivées/élevées + sélection naturelle sur l'Homme (ex. persistance de la lactase).
    \item \textbf{Médecine :} Vaccins, antibiotiques, chirurgie, moustiquaires imprégnées (lutte contre le paludisme au Cameroun) réduisent drastiquement la mortalité infectieuse, modifiant le paysage sélectif.
\end{itemize}
\end{exemplebox}

Cette capacité culturelle a permis une expansion démographique sans précédent : de ~ \num{1} à \num{10} millions d'humains il y a \num{10000} ans (paléolithique supérieur) à \textbf{\num{8,1} milliards en 2024}. Mais elle a un coût : l'Homme modifie massivement son propre environnement, créant de nouvelles pressions de sélection... sur lui-même et sur la biosphère entière.

\courssub{L'Homme, agent majeur de modification de l'environnement}

\begin{popfigure}
\centering
\begin{tikzpicture}
\begin{axis}[
    width=\linewidth,
    height=8cm,
    xlabel={Année},
    ylabel={Superficie (\si{\square\kilo\meter})},
    title={Évolution de la superficie du lac Tchad (1960--2020)},
    xmin=1960, xmax=2020,
    ymin=0, ymax=26000,
    xtick={1960,1970,1980,1990,2000,2010,2020},
    ytick={0,5000,10000,15000,20000,25000},
    grid=major,
    grid style={popline},
    legend pos=south west,
    legend style={fill=popcream,draw=popdark,font=\small},
    popcurve/.style={thick,popblue,mark=*,mark size=2}
]
\addplot[popcurve] coordinates {
    (1963,25000)
    (1970,22000)
    (1973,18000)
    (1975,15000)
    (1983,10000)
    (1985,3000)
    (1992,2500)
    (1994,2000)
    (2000,1500)
    (2005,2500)
    (2010,2000)
    (2015,1800)
    (2020,1500)
};
\addlegendentry{Superficie observée}
\addplot[poporange,dashed,thick] coordinates {(1960,25000) (2020,25000)};
\addlegendentry{Niveau de référence 1960}
\end{axis}
\end{tikzpicture}
\legende{Figure 1 : Régression dramatique du lac Tchad. Source : données NASA/UNEP, CBLT. La superficie est passée d'environ \SI{25000}{\square\kilo\meter} en 1963 à moins de \SI{2000}{\square\kilo\meter} dans les années 1980--1990, avec une légère reprise variable depuis. Causes : sécheresses sahéliennes (variabilité climatique), prélèvements d'eau pour l'irrigation (projet d'irrigation de la Komadougou Yobé, barrages en amont sur le Chari-Logone), évapotranspiration accrue. Conséquences : perte de biodiversité, conflits intercommunautaires (pêcheurs, éleveurs, agriculteurs), migration climatique.}
\end{popfigure}

\begin{popfigure}
\centering
\begin{tikzpicture}
\begin{axis}[
    width=\linewidth,
    height=8cm,
    xlabel={Année},
    ylabel={\ce{CO2} atmosphérique (ppm)},
    title={Courbe de Keeling : \ce{CO2} atmosphérique à Mauna Loa (1960--2023)},
    xmin=1960, xmax=2023,
    ymin=310, ymax=425,
    xtick={1960,1970,1980,1990,2000,2010,2020},
    ytick={310,330,350,370,390,410,420},
    grid=major,
    grid style={popline},
    legend pos=north west,
    legend style={fill=popcream,draw=popdark,font=\small},
    popcurve/.style={thick,popgreen,mark=none}
]
\addplot[popcurve] coordinates {
    (1960,316.9) (1965,320.0) (1970,325.7) (1975,331.1) (1980,338.7)
    (1985,346.0) (1990,354.2) (1995,360.8) (2000,369.5) (2005,379.8)
    (2010,389.9) (2015,400.8) (2020,414.2) (2023,421.1)
};
\addlegendentry{\ce{CO2} (moyenne annuelle)}
\addplot[popred,dashed,thick] coordinates {(1960,350) (2023,350)};
\addlegendentry{Seuil 350 ppm (sécurité climatique)}
\end{axis}
\end{tikzpicture}
\legende{Figure 2 : Augmentation continue du \ce{CO2} atmosphérique depuis 1960 (mesures de l'observatoire de Mauna Loa, NOAA/Scripps). Le seuil de 350 ppm, considéré comme limite supérieure de sécurité pour le climat, a été franchi vers 1988. En 2023, la moyenne annuelle dépasse 420 ppm (+ \SI{50}{\%} depuis l'ère préindustrielle). Cause principale : combustion d'énergies fossiles (charbon, pétrole, gaz) et changement d'affectation des sols (déforestation).}
\end{popfigure}

Au Cameroun, les impacts sont déjà visibles et documentés :
\begin{itemize}
    \item \textbf{Extrême-Nord :} Avancée du désert (progression du front sahélien de ~ \SIrange{5}{10}{\kilo\meter\per\an}), recul du lac Tchad (Figure 1), salinisation des sols, insécurité alimentaire chronique, conflits agro-pastoraux exacerbés par la raréfaction des ressources en eau et pâturages.
    \item \textbf{Bassin du Congo (Sud et Est) :} Déforestation pour l'agriculture itinérante sur brûlis, l'exploitation forestière industrielle (bois d'œuvre : ayous, sapelli, azobé), l'agro-industrie (palmier à huile, hévéa, cacao). Perte de ~ \SI{1,5}{\%} de la couverture forestière par an dans certaines zones. Le bassin du Congo stocke ~ \SI{60}{\giga\tonne} de carbone ; sa dégradation accélère le changement climatique global.
    \item \textbf{Urbanisation rapide :} Douala (~ \num{4,5} millions d'habitants) et Yaoundé (~ \num{4} millions) concentrent pollution de l'air (\ce{PM2,5} > \SI{35}{\micro\gram\per\cubic\meter}, seuil OMS = \SI{15}{\micro\gram\per\cubic\meter}), gestion insuffisante des déchets solides (décharges à ciel ouvert, pollution du Wouri et de la Mfoundi), inondations récurrentes par imperméabilisation des sols et occupation des zones inondables.
    \item \textbf{Littoral :} Érosion côtière (Kribi, Limbé), montée des eaux, salinisation des nappes phréatiques, menace sur les mangroves (zones de reproduction halieutique).
\end{itemize}

\begin{savaistu}[Le bassin du Congo : deuxième poumon de la planète]
Le bassin du Congo couvre ~ \num{3,7} millions de \si{\square\kilo\meter} (6 pays : Cameroun, RCA, RDC, Congo, Guinée équatoriale, Gabon). C'est la deuxième plus grande forêt tropicale humide après l'Amazonie. Elle abrite ~ \num{10000} espèces de plantes (\SI{30}{\%} endémiques), \num{400} mammifères (gorilles, chimpanzés, éléphants de forêt, okapis), \num{1000} oiseaux. Elle évapotranspire ~ \SI{1300}{\cubic\kilo\meter\per\an}, alimentant les pluies du Sahel et de l'Afrique australe (pompe biotique). Au Cameroun, les parcs nationaux de Lobéké, Boumba Bek, Nki et le sanctuaire de la Dja (patrimoine mondial UNESCO) en sont des joyaux. Sa préservation est une condition \emph{sine qua non} de l'atteinte des Objectifs de Développement Durable (ODD 13, 15).
\end{savaistu}

\courssub{Le développement durable : réponse collective et consciente}

\begin{definition}[Développement durable (Rapport Brundtland, 1987)]
Mode de développement qui répond aux besoins du présent sans compromettre la capacité des générations futures à répondre aux leurs. Il repose sur l'articulation de trois piliers indissociables : \textbf{environnemental}, \textbf{économique} et \textbf{social}.
\end{definition}

\begin{popfigure}
\centering
\begin{tikzpicture}[scale=0.9]
% Three intersecting circles
\def\r{3.2}
\def\d{2.0}
% Environnement (vert)
\fill[popgreen!30,draw=popgreen,thick] (0,0) circle (\r);
% Économie (bleu)
\fill[popblue!30,draw=popblue,thick] (\d,-\d*0.866) circle (\r);
% Social (orange)
\fill[poporange!30,draw=poporange,thick] (-\d,-\d*0.866) circle (\r);

% Labels
\node[popgreen,font=\bfseries\large] at (0,3.8) {Environnement};
\node[popblue,font=\bfseries\large] at (4.2,-3.5) {Économie};
\node[poporange,font=\bfseries\large] at (-4.2,-3.5) {Social};

% Intersections labels
\node[align=center,font=\small] at (0,0.5) {\textbf{Viable}\\(Environnement $\cap$ Économie)};
\node[align=center,font=\small] at (-1.8,-1.5) {\textbf{Équitable}\\(Social $\cap$ Environnement)};
\node[align=center,font=\small] at (1.8,-1.5) {\textbf{Supportable}\\(Économie $\cap$ Social)};
\node[align=center,font=\bfseries\small] at (0,-1.5) {\textbf{DURABLE}\\(Les 3 piliers)};

% Examples in each zone
\node[align=center,font=\footnotesize,text=popdark] at (0,2.5) {Protection de la biodiversité,\\ gestion des déchets,\\  énergies renouvelables};
\node[align=center,font=\footnotesize,text=popdark] at (3.5,-1.0) {Économie circulaire,\\ efficacité énergétique,\\ finance verte};
\node[align=center,font=\footnotesize,text=popdark] at (-3.5,-1.0) {Éducation environnementale,\\ justice climatique,\\ droits autochtones};
\node[align=center,font=\footnotesize,text=popdark] at (0,-3.0) {Agroforesterie,\\ éco-tourisme communautaire,\\ villes durables};
\end{tikzpicture}
\legende{Figure 3 : Les trois piliers du développement durable et leurs intersections. Au centre, la zone « Durable » où les trois dimensions sont satisfaites simultanément. Aucune politique sectorielle ne peut réussir durablement si elle ignore les deux autres piliers.}
\end{popfigure}

\begin{methode}[Analyser une situation locale à la lumière du développement durable]
\begin{enumerate}
    \item \textbf{Identifier les pressions :} Quelles activités humaines modifient le milieu ? (ex. coupe de bois de chauffe à Douala, irrigation au Lac Tchad, pesticides cotonnières dans le Bénoué).
    \item \textbf{Évaluer les impacts sur les 3 piliers :} 
    \begin{itemize}
        \item Environnemental : biodiversité, qualité eau/air/sol, climat, ressources.
        \item Économique : revenus, emplois, coûts de la dégradation, résilience.
        \item Social : santé, équité, sécurité alimentaire, conflits, patrimoine culturel.
    \end{itemize}
    \item \textbf{Proposer des leviers d'action :} Réglementation, incitations économiques, éducation, technologies propres, gouvernance participative.
    \item \textbf{Vérifier la cohérence intergénérationnelle :} La solution d'aujourd'hui ne doit pas créer le problème de demain.
\end{enumerate}
\end{methode}

Exemples d'actions concrètes au Cameroun :
\begin{itemize}
    \item \textbf{Reboisement / restauration :} Opération « Green Sahel » (Extrême-Nord), programme de restauration des paysages forestiers (AFR100 : engagement de restaurer \num{12} millions d'\si{\hectare} d'ici 2030), corridors écologiques entre aires protégées (TRIDOM, TNS).
    \item \textbf{Énergies renouvelables :} Centrales solaires de Maroua (\SI{20}{\mega\watt}), Guider (\SI{30}{\mega\watt}), projet Nachtigal (\SI{420}{\mega\watt} hydroélectrique, réduction ~ \SI{1}{\mega\tonne} \ce{CO2}/an vs thermique). Promotion des foyers améliorés (réduction bois de chauffe, pollution intérieure).
    \item \textbf{Gestion durable de l'eau :} Commission du Bassin du Lac Tchad (CBLT) : plan d'action pour le transfert d'eau du bassin de l'Oubangui (projet Transaqua, à l'étude), gestion concertée des barrages (Lagdo, Maga, Yagoua).
    \item \textbf{Éducation environnementale :} Intégration dans les programmes scolaires (SVT, Géographie, Éducation à la citoyenneté), clubs « Jeunes pour l'Environnement », journées nationales de l'arbre (14 juillet) et de l'environnement (5 juin).
\end{itemize}

\courssub{Lien entre évolution humaine et nécessité d'adaptation consciente}

\begin{aretenir}[Synthèse évolutive]
La capacité d'adaptation qui a permis à la lignée humaine de survivre à des changements climatiques majeurs (alternances glaciaires/interglaciaires, aridification africaine il y a ~ \num{2,8} Ma), de coloniser tous les biomes et de devenir l'espèce dominante, repose sur deux piliers :
\begin{enumerate}
    \item \textbf{Plasticité phénotypique et diversité génétique} (adaptation biologique lente, stochastique).
    \item \textbf{Capacité culturelle cumulative} (adaptation rapide, dirigée, transmissible).
\end{enumerate}
Aujourd'hui, la vitesse et l'ampleur des modifications anthropiques de l'environnement (changement climatique, érosion de la biodiversité, cycles biogéochimiques perturbés) dépassent les capacités d'adaptation \emph{biologique} de la plupart des espèces, y compris la nôtre. Seule une adaptation \emph{culturelle consciente, collective, anticipative et éthique} — c'est-à-dire le développement durable — peut préserver les conditions d'habitabilité de la Terre pour \emph{Homo sapiens} et le vivant dont il dépend.
\end{aretenir}

\begin{exoresolu}[Analyse de la régression du lac Tchad]
\textbf{Énoncé :} À partir de la Figure 1 (superficie du lac Tchad 1960--2020) et de vos connaissances, expliquez les causes de la régression du lac Tchad et proposez deux mesures d'adaptation durable pour les populations riveraines.

\textbf{Solution détaillée :}
\begin{enumerate}
    \item \textbf{Analyse de la courbe :} La superficie passe de ~ \SI{25000}{\square\kilo\meter} (1963) à ~ \SIrange{1500}{2500}{\square\kilo\meter} (années 1980--2020), soit une perte de ~ \SI{90}{\%}. Deux phases : chute brutale 1963--1985 (sécheresses sahéliennes majeures 1972--1973, 1983--1984), puis stabilisation à bas niveau avec variations interannuelles.
    \item \textbf{Causes combinées :}
    \begin{itemize}
        \item \textbf{Climatiques :} Déficit pluviométrique persistant sur le bassin versant (Sahel), augmentation de l'évapotranspiration (+ \SI{1,5}{\celsius} depuis 1960), variabilité accrue (ENSO).
        \item \textbf{Anthropiques :} Prélèvements d'eau pour l'irrigation en amont (barrages de Maga, Yagoua sur la Logone ; projet Komadougou Yobé), déforestation du bassin versant (érosion, colmatation), croissance démographique (besoins en eau, bois, terres).
    \end{itemize}
    \item \textbf{Mesures d'adaptation durable (3 piliers) :}
    \begin{itemize}
        \item \textbf{Environnemental + Social :} Restauration des zones humides (réserves de biosphère de Waza-Logone), reboisement du bassin versant (fixation des sols, infiltration), gestion communautaire des ressources halieutiques (pêche durable, périodes de repos biologique).
        \item \textbf{Économique + Social :} Développement de l'agroforesterie (acacia Sénégal/gommier, néreé) et de l'élevage sédentarisé amélioré (fourrage, santé animale) pour réduire la pression sur les pâturages lacustres ; diversification des revenus (artisanat, éco-tourisme ornithologique) pour réduire la vulnérabilité climatique.
    \end{itemize}
    \item \textbf{Conclusion :} L'adaptation ne peut être seulement technique (forage, pompage) ; elle doit être systémique, intégrant la gouvernance transfrontalière (CBLT) et l'équité intergénérationnelle.
\end{enumerate}
\end{exoresolu}

\begin{exercice}[Application : Bilan carbone d'une exploitation cacaoyère au Sud-Cameroun]
Une exploitation de \SI{50}{\hectare} de cacaoyers (variété hybride, \num{1100} plants/ha) produit \SI{800}{\kilo\gram\per\hectare\per\an} de fèves sèches. L'itinéraire technique comprend : débroussaillage manuel, traitement phytosanitaire (insecticide + fongicide, 3 passages/an), fertilisation NPK (\SI{200}{\kilo\gram\per\hectare\per\an}), récolte manuelle, fermentation (bois de chauffe), séchage au soleil. 
\begin{enumerate}
    \item Estimez les principales sources d'émissions de GES (\ce{CO2}, \ce{N2O}, \ce{CH4}) pour cette exploitation.
    \item Proposez 3 pratiques d'agroforesterie cacaoyère « climato-intelligente » qui réduisent le bilan carbone tout en maintenant la productivité.
    \item Expliquez en quoi ces pratiques relèvent d'une adaptation culturelle durable au sens de la leçon.
\end{enumerate}
\end{exercice}

\begin{corrige}[Exercice : Bilan carbone cacaoyère]
\begin{enumerate}
    \item \textbf{Sources d'émissions :} (a) Déforestation initiale (si conversion forêt) : ~ \SIrange{200}{300}{\tonne\ CO2\per\hectare} émis une fois (perte biomasse + sol). (b) Fertilisation NPK : fabrication (énergie fossile) + émission \ce{N2O} du sol (~ \SI{1}{\%} N appliqué) → \SI{200}{\kilo\gram\per\hectare\per\an} N × \SI{1}{\%} × 298 (PRG \ce{N2O}) ≈ \SI{0,6}{\tonne\ CO2eq\per\hectare\per\an}. (c) Phytosanitaires : synthèse + transport ≈ \SIrange{50}{100}{\kilo\gram\ CO2eq\per\hectare\per\an}. (d) Fermentation : bois de chauffe (non renouvelable si prélevé en forêt) ≈ \SIrange{0,5}{1}{\tonne\ bois\per\tonne\ feves} → ~ \SIrange{400}{800}{\kilo\gram\ CO2\per\hectare\per\an}. (e) Transport fèves (camion) : variable. Total hors changement d'affectation des sols : ~ \SIrange{1,5}{2,5}{\tonne\ CO2eq\per\hectare\per\an}.
    \item \textbf{Pratiques climato-intelligentes :} 
    \begin{itemize}
        \item \textbf{Associations arbres d'ombrage (espèces légumineuses : \emph{Gliricidia}, \emph{Leucaena}, \emph{Albizia}) :} Fixation azotée (réduction engrais N), séquestration carbone biomasse aérienne/racinaire (~ \SIrange{5}{10}{\tonne\ CO2\per\hectare\per\an}), régulation microclimatique (température, humidité), biodiversité auxiliaires (réduction phytos).
        \item \textbf{Couverture du sol (mucuna, pueraria, résidus de taille) :} Stockage C sol, réduction érosion, suppression herbicides, activité biologique.
        \item \textbf{Séchage solaire amélioré (séchoirs à effet de serre) :} Élimination bois de chauffe, qualité fèves supérieure (prime prix).
    \end{itemize}
    \item \textbf{Lien adaptation culturelle :} Ces pratiques sont des innovations techniques transmises par formation (vulgarisation, écoles-champs), cumulatives (amélioration variétale + itinéraire), réversibles (choix de l'agriculteur), et répondent aux 3 piliers : environnemental (C stocké, biodiversité), économique (coûts intrants ↓, rendement stable, prime qualité), social (sécurité alimentaire, santé paysanne, transmission intergénérationnelle du capital sol). Elles illustrent la capacité de l'Homme à adapter \emph{conscientement} son système de production face au changement climatique qu'il a lui-même contribué à créer.
\end{enumerate}
\end{corrige}

\courssub{Perspectives : l'évolution continue, à nous de la diriger}

L'histoire évolutive de l'Homme montre que notre lignée a survécu grâce à sa plasticité et à son ingéniosité. Le défi actuel est inédit : pour la première fois, une espèce \emph{connaît} les mécanismes de l'évolution (génétique, culturelle, écosystémique) et \emph{mesure} en temps réel l'impact de ses actions sur sa propre niche écologique. Le développement durable n'est pas une option morale, c'est une \textbf{nécessité adaptative} : il s'agit d'aligner notre évolution culturelle sur les limites biophysiques de la Terre. Au Cameroun, pays « Afrique en miniature » par sa diversité écologique et culturelle, cet alignement passe par la valorisation des savoirs locaux (agroforesterie traditionnelle, gestion communautaire des ressources), l'appropriation des sciences (suivi climatique, biotechnologies vertes) et une gouvernance inclusive. C'est le sens de votre formation en SVT : comprendre le vivant pour agir en citoyen éclairé, capable de contribuer à l'adaptation durable de notre espèce dans son environnement.

\newpage

\coursec{Activité d'intégration}

\coursec{Leçon 16 : Sensibilisation sur la nécessité d'adaptation de l'Homme avec son environnement}

\courssub{Objectifs de l'activité}
Cette activité d'intégration vise à mobiliser l'ensemble des savoir-faire de la leçon dans une démarche d'investigation scientifique complète. Elle permet à l'élève de :
\begin{itemize}
    \item Analyser des données morphométriques fossiles pour identifier les \cle{critères de l'hominisation} (bipédie, encéphalisation, réduction du prognathisme) et proposer une chronologie relative.
    \item Exploiter des données moléculaires (distances génétiques) pour \cle{construire un arbre phylogénétique} de la lignée humaine et dater les divergences majeures.
    \item Établir un lien argumenté entre la \cle{capacité d'adaptation biologique et culturelle} de l'Homme au cours de son évolution et la \cle{nécessité d'un développement durable} face aux changements environnementaux actuels (régression du lac Tchad, déforestation au Cameroun).
    \item Communiquer ses résultats par une restitution écrite structurée et une argumentation critique.
\end{itemize}

\courssub{Situation : De Toumaï au développement durable — Enquête d'un paléoanthropologue camerounais}
\textbf{Contexte.} Vous êtes stagiaire au Laboratoire de Paléoanthropologie et de Préhistoire de l'Université de Yaoundé I, en collaboration avec l'Institut de Recherche pour le Développement (IRD). Votre tuteur, le Dr \emph{Ngoufo}, vous confie un dossier pluridisciplinaire pour préparer une communication au colloque « \emph{Homme, Environnement et Développement Durable en Afrique Centrale} ». Ce dossier comporte trois volets indépendants mais convergents.

\textbf{Volet 1 : Analyse morphométrique de crânes fossiles (Données de terrain).}
Lors d'une mission de prospection dans le bassin du Djurab (Tchad, région du lac Tchad) et en Afrique de l'Est, quatre spécimens crâniens (références Crâne A, B, C, D) ont été découverts dans des couches géologiques distinctes. Les mesures suivantes ont été relevées par l'équipe :

\begin{center}
\begin{tabular}{|c|S[table-format=4.0]|S[table-format=1.2]|S[table-format=2.0]|c|S[table-format=2.1]|}\hline
\rowcolor{popblueL}
\textbf{Paramètre} & \textbf{Crâne A} & \textbf{Crâne B} & \textbf{Crâne C} & \textbf{Crâne D} \\ \hline
Volume crânien estimé (cm$^3$) & 350 & 450 & 650 & 1350 \\ \hline
Position du trou occipital (indice bipédie) & 0.48 & 0.45 & 0.49 & 0.52 \\ \hline
Angle de prognathisme facial (°) & 65 & 55 & 45 & 15 \\ \hline
Présence d'une crête sagittale & Oui (marquée) & Oui (faible) & Non & Non \\ \hline
Datation radiométrique ($^{40}$Ar/$^{39}$Ar, Ma) & 7.0 & 3.2 & 1.8 & 0.3 \\ \hline
\end{tabular}
\end{center}
\emph{Note :} L'indice de bipédie = distance (trou occipital -- bord antérieur du foramen) / longueur base crânienne. Valeur proche de 0,5 = bipédie ; $<$ 0,3 = quadrupédie.

\textbf{Volet 2 : Données moléculaires — Matrice des distances génétiques.}
Le laboratoire de génétique moléculaire a séquencé le gène mitochondrial \emph{COI} (cytochrome c oxydase sous-unité I) et une région non codante de l'ADN nucléaire pour quatre espèces actuelles. Le tableau ci-dessous donne le \textbf{pourcentage de divergence de séquences nucléotidiques} (distance $p$) calculé par alignement multiple (ClustalW) :

\begin{center}
\begin{tabular}{|c|cccc|}\hline
\rowcolor{popgreenL}
\textbf{Distances $p$ (\%)} & \textbf{Homme (\emph{Homo sapiens})} & \textbf{Chimpanzé (\emph{Pan troglodytes})} & \textbf{Gorille (\emph{Gorilla gorilla})} & \textbf{Orang-outan (\emph{Pongo pygmaeus})} \\ \hline
Homme & 0,0 & 1,2 & 1,6 & 3,1 \\ \hline
Chimpanzé & 1,2 & 0,0 & 1,5 & 3,0 \\ \hline
Gorille & 1,6 & 1,5 & 0,0 & 2,9 \\ \hline
Orang-outan & 3,1 & 3,0 & 2,9 & 0,0 \\ \hline
\end{tabular}
\end{center}
\emph{Hypothèse de travail :} L'horloge moléculaire pour ces séquences est calibrée à \textbf{0,15 \% de divergence par million d'années (Ma)} pour la lignée hominidée (taux moyen de substitution neutre).

\textbf{Volet 3 : Environnement actuel et enjeux durables au Cameroun.}
\begin{itemize}
    \item \textbf{Lac Tchad :} La superficie de l'eau libre a drastiquement diminué : \SI{25000}{km^2} en 1963 $\rightarrow$ \SI{2500}{km^2} en 1985 (sécheresses majeures) $\rightarrow$ \SI{14000}{km^2} en 2020 (légère reprise pluvieuse mais instable). La population riveraine est passée de 13 millions (1990) à plus de 30 millions (2020).
    \item \textbf{Déforestation au Cameroun :} Taux de perte de couverture forestière primaire : \SI{0,6}{\%} par an (période 2001-2021), soit environ \SI{200000}{ha/an}. Causes principales : agriculture itinérante sur brûlis (\SI{60}{\%}), exploitation forestière légale/illégale (\SI{25}{\%}), extension agro-industrielle (palmier à huile, hévéa) (\SI{15}{\%}).
    \item \textbf{Contexte local (Extrême-Nord) :} Avancée du désert, salinisation des sols, conflits agro-pasteurs, pression sur les ressources halieutiques du lac.
\end{itemize}

\courssub{Consignes}
\textbf{Volet 1 : Preuves fossiles et critères d'hominisation}
\begin{enumerate}
    \item En vous appuyant sur le tableau du Volet 1, identifiez les \textbf{trois critères majeurs de l'hominisation} illustrés par l'évolution des valeurs numériques entre le Crâne A et le Crâne D.
    \item Classez les quatre crânes dans l'ordre chronologique (du plus ancien au plus récent) en justifiant par \textbf{deux arguments morphologiques} et la \textbf{datation absolue}.
    \item Attribuez à chaque crâne le genre le plus probable (\emph{Sahelanthropus}, \emph{Australopithecus}, \emph{Homo habilis}/\emph{erectus}, \emph{Homo sapiens}) en explicitant le raisonnement.
\end{enumerate}

\textbf{Volet 2 : Preuves moléculaires et phylogénie}
\begin{enumerate}\setcounter{enumi}{3}
    \item À partir de la matrice de distances, construisez l'\textbf{arbre phylogénétique} (cladogramme) des quatre espèces en indiquant les nœuds de divergence.
    \item Calculez l'\textbf{âge approximatif de la divergence Homme--Chimpanzé} et de la divergence Gorille--(Homme+Chimpanzé) en utilisant l'horloge moléculaire fournie. Exprimez les résultats en millions d'années (Ma).
    \item Comparez la date de divergence Homme--Chimpanzé obtenue par l'horloge moléculaire avec l'âge du Crâne A (Volet 1). Commentez la cohérence entre données fossiles et données moléculaires.
\end{enumerate}

\textbf{Volet 3 : Adaptation, environnement et développement durable}
\begin{enumerate}\setcounter{enumi}{6}
    \item Calculez le \textbf{taux de variation moyen annuel} de la superficie du lac Tchad sur la période 1963--1985 (période de régression majeure) et sur la période 1985--2020 (période de reprise partielle). Exprimez en km$^2$/an.
    \item En vous fondant sur l'ensemble du dossier (évolution biologique/culturelle de l'Homme, données environnementales actuelles), rédigez un \textbf{texte argumenté de 12 à 15 lignes} expliquant en quoi la capacité d'adaptation de l'Homme, qui a permis sa survie et son expansion, impose aujourd'hui des choix de \textbf{développement durable}. Votre texte doit proposer \textbf{deux actions concrètes et localisées} (région de l'Extrême-Nord ou votre localité) pour concilier besoins humains et préservation des écosystèmes.
\end{enumerate}

\courssub{Ressources à mobiliser}
\begin{itemize}
    \item \textbf{Critères d'hominisation :} Bipédie (position trou occipital, bassin, fémur), Encéphalisation (volume crânien, complexification cérébrale), Réduction du prognathisme et de la mastication (dents, crêtes), Main préhensile (pouce opposant), Fabrication d'outils (culture).
    \item \textbf{Arbre phylogénétique :} Méthode UPGMA ou Neighbor-Joining simplifiée (recherche des plus faibles distances pour grouper les taxons). Racine = groupe externe (Orang-outan).
    \item \textbf{Horloge moléculaire :} $T_{\text{divergence}} = \frac{\text{Distance } p \text{ (\%)}}{2 \times \text{Taux de substitution (\%/Ma)}}$. Le facteur 2 tient compte de l'accumulation indépendante des mutations sur les deux lignées depuis l'ancêtre commun.
    \item \textbf{Développement durable :} Définition Brundtland (1987) : « Développement qui répond aux besoins du présent sans compromettre la capacité des générations futures de répondre aux leurs ». Trois piliers : Environnemental, Social, Économique.
    \item \textbf{Formules :} Taux de variation moyen $= \frac{V_{\text{final}} - V_{\text{initial}}}{\Delta t}$.
\end{itemize}

\courssub{Démarche et corrigé commenté}
\begin{exoresolu}[Corrigé détaillé de l'activité d'intégration]
\textbf{Énoncé de l'activité (rappel des questions)}
\begin{enumerate}
    \item Identifier les trois critères majeurs de l'hominisation (Volet 1).
    \item Classer les crânes chronologiquement et les attribuer (Volet 1).
    \item Construire l'arbre phylogénétique et dater les nœuds (Volet 2).
    \item Comparer datations fossiles et moléculaires (Volet 2).
    \item Calculer les taux de variation du Lac Tchad (Volet 3).
    \item Rédiger un texte argumenté adaptation / développement durable (Volet 3).
\end{enumerate}
\tcblower
\textbf{Solution détaillée}

\textbf{Volet 1 : Preuves fossiles et critères d'hominisation}

\textbf{1. Identification des critères.}
L'analyse du tableau montre trois tendances évolutives majeures de A à D :
\begin{itemize}
    \item \textbf{L'encéphalisation :} Volume crânien passe de \SI{350}{cm^3} (A) à \SI{1350}{cm^3} (D), soit une multiplication par ~3,9. C'est l'augmentation de la taille et de la complexité du cerveau.
    \item \textbf{La bipédie :} L'indice de position du trou occipital évolue de 0,48 (A) à 0,52 (D), se rapprochant de la valeur théorique 0,5 caractéristique de la bipédie stricte (trou occipital centré sous le crâne). Le Crâne A montre déjà une adaptation à la bipédie (indice > 0,45).
    \item \textbf{La réduction du prognathisme et de l'appareil masticatoire :} L'angle de prognathisme diminue fortement (65° $\rightarrow$ 15°) et la crête sagittale (insertion des muscles masticateurs puissants) disparaît entre B et C.
\end{itemize}
\emph{Note :} La fabrication d'outils (critère culturel) n'est pas directement visible sur ces crânes mais est corrélée à l'encéphalisation à partir de \emph{Homo}.

\textbf{2. Classement chronologique et attribution.}
L'ordre chronologique est donné directement par la datation radiométrique ($^{40}$Ar/$^{39}$Ar) : \textbf{A (7,0 Ma) $<$ B (3,2 Ma) $<$ C (1,8 Ma) $<$ D (0,3 Ma)}.
La morphologie confirme cet ordre :
\begin{itemize}
    \item \textbf{Crâne A (7,0 Ma) :} Volume chimpanzé (350 cm$^3$), mais trou occipital antérieur (bipédie), fort prognathisme, crête sagittale. $\rightarrow$ \textbf{\emph{Sahelanthropus tchadensis}} (Toumaï). C'est le plus ancien représentant connu de la lignée humaine post-divergence.
    \item \textbf{Crâne B (3,2 Ma) :} Volume légèrement supérieur (450 cm$^3$), bipédie confirmée, prognathisme encore fort, crête sagittale faible. $\rightarrow$ \textbf{\emph{Australopithecus afarensis}} (type « Lucy ») ou \emph{A. africanus}.
    \item \textbf{Crâne C (1,8 Ma) :} Saut encéphalique (650 cm$^3$), bipédie moderne, prognathisme réduit, pas de crête. $\rightarrow$ \textbf{\emph{Homo habilis}} (ou \emph{H. rudolfensis} / début \emph{H. erectus}). Apparition de l'Oldowayen (outils taillés).
    \item \textbf{Crâne D (0,3 Ma) :} Volume moderne (1350 cm$^3$), bipédie stricte, face orthognathe (15°). $\rightarrow$ \textbf{\emph{Homo sapiens}} (archaïque ou moderne précoce).
\end{itemize}

\textbf{Volet 2 : Preuves moléculaires et phylogénie}

\textbf{3. Construction de l'arbre phylogénétique (Méthode UPGMA simplifiée).}
On recherche les plus faibles distances pour former les groupes frères (clades) :
\begin{itemize}
    \item Distance min = 1,2 \% (Homme--Chimpanzé). $\rightarrow$ Groupe frère : \textbf{(Homme, Chimpanzé)}. Nœud 1.
    \item Distance suivante min vers ce groupe : Moyenne (Homme-Gorille ; Chimpanzé-Gorille) = (1,6 + 1,5)/2 = 1,55 \%. $\rightarrow$ Rattachement du Gorille. Nœud 2.
    \item Distance vers Orang-outan : Moyenne ~ 3,0 \%. $\rightarrow$ Rattachement de l'Orang-outan (groupe externe). Racine.
\end{itemize}

\begin{popfigure}
\begin{tikzpicture}[scale=0.9, every node/.style={font=\small}]
    % Coordonnées : x = temps avant présent (Ma) * 1cm/Ma. Présent à x=10. Racine à x=0.
    % Racine (Orang-outan split) : 10 Ma -> x=0
    % Nœud 2 (Gorille split) : 5.2 Ma -> x=4.8
    % Nœud 1 (Homme-Chimpanzé split) : 4.0 Ma -> x=6.0
    % Présent (Tips) : 0 Ma -> x=10
    
    \coordinate (root) at (0,0);
    \coordinate (nodeG) at (4.8,0);
    \coordinate (nodeHC) at (6.0,0);
    \coordinate (tip) at (10,0);
    
    % Branches principales
    \draw[popdark, thick] (root) -- (nodeG); % Tronc commun
    \draw[popdark, thick] (nodeG) -- (nodeHC); % Tronc Homme+Chimpanzé+Gorille
    \draw[popdark, thick] (nodeHC) -- (tip); % Tronc Homme+Chimpanzé
    
    % Espèces actuelles (Tips à x=10)
    \coordinate (H) at (10, 2.5);
    \coordinate (C) at (10, 0.8);
    \coordinate (G) at (10, -1.0);
    \coordinate (O) at (10, -2.8);
    
    % Branches terminales
    \draw[popdark, thick] (nodeHC) -- (H) node[right] {\textbf{Homme} (\emph{Homo sapiens})};
    \draw[popdark, thick] (nodeHC) -- (C) node[right] {\textbf{Chimpanzé} (\emph{Pan troglodytes})};
    \draw[popdark, thick] (nodeG) -- (G) node[right] {\textbf{Gorille} (\emph{Gorilla gorilla})};
    \draw[popdark, thick] (root) -- (O) node[right] {\textbf{Orang-outan} (\emph{Pongo pygmaeus})};
    
    % Nœuds
    \fill[poporange] (root) circle (3pt) node[left=2pt] {\textbf{Racine (10,0 Ma)}};
    \fill[popgreen] (nodeG) circle (3pt) node[above=2pt] {\textbf{Nœud 2 (5,2 Ma)}};
    \fill[poppink] (nodeHC) circle (3pt) node[above=2pt] {\textbf{Nœud 1 (4,0 Ma)}};
    
    % Échelle de temps
    \draw[<->, popblue, thick] (0,-3.5) -- (10,-3.5);
    \foreach \x/\label in {0/10, 4.8/5,2, 6.0/4,0, 10/0}
        \draw[popblue, thick] (\x,-3.5) -- (\x,-3.7) node[below] {\label Ma};
    \node[popblue, below] at (5,-3.9) {Temps avant présent (Ma)};
\end{tikzpicture}
\legende{Arbre phylogénétique des Hominidés construit par UPGMA à partir des distances $p$. L'Orang-outan est le groupe externe. Les nœuds sont datés par l'horloge moléculaire (taux 0,15 \%/Ma). Les longueurs de branches sont proportionnelles au temps.}
\end{popfigure}

\textbf{4. Datation par l'horloge moléculaire.}
Formule : $T = \frac{p}{2 \times \mu}$ avec $\mu = 0,15 \%\text{/Ma}$.
\begin{itemize}
    \item \textbf{Divergence Homme--Chimpanzé (Nœud 1) :} $p = 1,2 \%$.
    $T_{HC} = \frac{1,2}{2 \times 0,15} = \frac{1,2}{0,3} = \mathbf{4,0 \text{ Ma}}$.
    \item \textbf{Divergence Gorille--(Homme+Chimpanzé) (Nœud 2) :} Distance moyenne au groupe (H+C) = 1,55 \%.
    $T_{G(HC)} = \frac{1,55}{0,3} \approx \mathbf{5,2 \text{ Ma}}$.
    \item \textbf{Divergence Orang-outan--Autres (Racine) :} Distance moyenne ~ 3,0 \%.
    $T_{\text{Racine}} = \frac{3,0}{0,3} = \mathbf{10,0 \text{ Ma}}$.
\end{itemize}
\emph{Remarque :} Les valeurs réelles acceptées par la communauté scientifique sont plus anciennes (Homme-Chimpanzé $\sim$ 6-7 Ma, Gorille $\sim$ 8-9 Ma). Notre taux $\mu = 0,15$ est un peu rapide ou les distances $p$ sous-estimées (saturation). L'exercice montre la méthode.

\textbf{5. Comparaison Fossiles / Moléculaire.}
Le Crâne A (\emph{Sahelanthropus}) est daté de \textbf{7,0 Ma}. L'horloge moléculaire (avec ce taux) donne une divergence Homme-Chimpanzé à \textbf{4,0 Ma}.
\textbf{Incohérence apparente :} Le fossile (7 Ma) est plus ancien que la divergence moléculaire calculée (4 Ma). C'est impossible : un représentant de la lignée humaine post-divergence ne peut pas exister avant la divergence elle-même.
\textbf{Interprétation :} Cela indique que le taux de substitution $\mu = 0,15 \%\text{/Ma}$ est \textbf{surévalué} pour cette période/gène, ou que les distances $p$ ne sont pas corrigées des substitutions multiples (saturation). Un taux plus réaliste ($\mu \approx 0,08 - 0,10 \%\text{/Ma}$) donne $T_{HC} \approx 6-7,5 \text{ Ma}$, compatible avec Toumaï. Cela illustre la \textbf{nécessité de calibrer l'horloge moléculaire sur des fossiles datés} (points de calibration).

\textbf{Volet 3 : Adaptation, environnement et développement durable}

\textbf{6. Taux de variation de la superficie du Lac Tchad.}
\begin{itemize}
    \item \textbf{1963--1985 (Régression) :} $\Delta S = 2\,500 - 25\,000 = -22\,500 \text{ km}^2$. $\Delta t = 22 \text{ ans}$.
    Taux $= \frac{-22\,500}{22} \approx \mathbf{-1\,023 \text{ km}^2/\text{an}}$ (Perte massive, ~4\%/an).
    \item \textbf{1985--2020 (Reprise partielle) :} $\Delta S = 14\,000 - 2\,500 = +11\,500 \text{ km}^2$. $\Delta t = 35 \text{ ans}$.
    Taux $= \frac{+11\,500}{35} \approx \mathbf{+329 \text{ km}^2/\text{an}}$ (Reprise 3 fois plus lente que la perte).
\end{itemize}

\textbf{7. Texte argumenté (Proposition de rédaction).}
\begin{quote}
\emph{L'évolution de la lignée humaine, de Toumaï à l'Homme moderne, témoigne d'une capacité d'adaptation exceptionnelle : la bipédie a libéré les mains, l'encéphalisation a permis l'innovation technique et la culture, transformant l'Homme en un agent géologique majeur. Cependant, cette même capacité, amplifiée par la révolution industrielle et l'agriculture intensive, conduit aujourd'hui à une surexploitation des ressources. Au Cameroun, la régression du lac Tchad (-1023 km$^2$/an entre 1963-1985) et la déforestation (200 000 ha/an) menacent la sécurité alimentaire et hydrique de 30 millions de personnes dans le bassin. Le développement durable n'est plus un choix éthique mais une condition de survie, imposant de remplacer l'adaptation « par destruction » par une adaptation « par régulation ». }

\emph{Concrètement, dans l'Extrême-Nord, il faut : (1) Généraliser l'agroforesterie et les techniques de zaï (fosses de rétention d'eau et de matière organique) pour restaurer les sols dégradés et augmenter les rendements sans défricher ; (2) Mettre en place une gestion concertée et transfrontalière de l'eau du bassin du Logone-Chari (barrages de régulation, irrigation goutte-à-goutte) pour sécuriser l'agriculture et la pêche face à la variabilité climatique.}
\end{quote}
\end{exoresolu}

\courssub{Bilan de l'activité}
Cette activité a permis de mettre en évidence la \textbf{convergence des preuves} (paléontologiques, anatomiques, moléculaires, géochronologiques) pour reconstituer l'histoire évolutive de l'Homme. Elle a illustré la rigueur de la démarche scientifique : les données fossiles (Toumaï, 7 Ma) servent à calibrer l'horloge moléculaire, dont les prédictions doivent être cohérentes avec le registre stratigraphique. 

Sur le plan conceptuel, l'activité ancre la notion d'\textbf{hominisation} comme un processus mosaïque (bipédie précoce, encéphalisation tardive) et non linéaire. Enfin, elle ancre l'enseignement de l'évolution dans une \textbf{responsabilité citoyenne} : la compréhension des mécanismes d'adaptation passés (sélection naturelle, innovation culturelle) éclaire l'urgence d'une adaptation présente \emph{consciente et durable} face aux changements globaux (climat, biodiversité) dont l'Homme est désormais le principal acteur. La capacité d'anticipation, fruit de l'encéphalisation, est l'outil principal pour éviter l'effondrement des socio-écosystèmes comme celui du bassin du Lac Tchad.

\newpage

\coursec{Méthodes et exercices résolus}

\coursec{I. Les preuves de l'évolution de l'Homme}

\courssub{Établir les preuves de l'évolution humaine}

\begin{methode}[Identifier et classer les preuves de l'évolution humaine]
Pour établir les preuves de l'évolution de l'Homme à partir de documents (fossiles, données anatomiques, moléculaires, embryologiques), procéder en étapes :
\begin{enumerate}
\item \cle{Identifier} la nature du document : fossile (crâne, mandibule, squelette post-crânien), outil lithique, vestige de sépulture, séquence d'ADN ou de protéine, datation radiochronologique.
\item \cle{Classer} la preuve dans l'une des quatre grandes catégories : preuves \textbf{paléontologiques} (fossiles stratifiés), preuves \textbf{anatomiques comparées} (organes homologues, organes vestigiaux), preuves \textbf{embryologiques} (similitudes des stades précoces), preuves \textbf{moléculaires} (comparaison séquences ADN/protéines, caryotypes).
\item \cle{Dater} le vestige si possible : méthode au carbone 14 (\ce{^{14}C}) jusqu'à environ \num{50000} ans ; méthode potassium-argon (\ce{^{40}K}/\ce{^{40}Ar}) ou argon-argon pour les couches volcaniques anciennes (millions d'années) ; paléomagnétisme, biochronologie.
\item \cle{Comparer} les caractères du fossile avec ceux des formes actuelles (Homme moderne, grands singes anthropomorphes) : capacité crânienne (\si{cm^3}), indice d'enveloppement, prognathisme facial, position du trou occipital (foramen magnum), forme du bassin (largeur, hauteur), courbure de la colonne vertébrale (cyphose/lordose), angulation fémorale.
\item \cle{Conclure} : montrer que les formes fossiles anciennes diffèrent des formes actuelles et que les caractères dérivés « humains » apparaissent progressivement et de manière ordonnée au cours du temps géologique — c'est la preuve d'une \cle{évolution} buissonnante.
\end{enumerate}
\textbf{Réflexes de rédaction} : toujours citer la donnée chiffrée du document (capacité crânienne en \si{cm^3}, âge en millions d'années \textbf{Ma}) avant de conclure ; ne jamais écrire « le fossile prouve » sans dire \emph{comment} (quel caractère, quelle comparaison, quelle tendance temporelle).
\end{methode}

\begin{exoresolu}[Toumaï, Lucy et l'Homme moderne : preuves paléontologiques]
On dispose des données suivantes sur trois spécimens fossiles ou actuels :
\begin{center}
\begin{tabularx}{\linewidth}{|l|X|X|X|}
\hline
\rowcolor{popblueL} \textbf{Spécimen} & \textbf{Âge estimé} & \textbf{Capacité crânienne} & \textbf{Position du trou occipital (foramen magnum)} \\
\hline
Toumaï (\emph{Sahelanthropus tchadensis}, Toros-Menalla, Tchad) & environ 7 Ma & environ \SI{350}{cm^3} & position antérieure (sous la base du crâne) \\
\hline
Lucy (\emph{Australopithecus afarensis}, Hadar, Éthiopie) & environ 3,2 Ma & environ \SI{450}{cm^3} & position intermédiaire \\
\hline
\emph{Homo sapiens} actuel & actuel & environ \SI{1350}{cm^3} (moyenne \SIrange{1200}{1500}{cm^3}) & centré sous le crâne (base crânienne flexionnée) \\
\hline
\end{tabularx}
\end{center}
\begin{enumerate}
\item À quelle catégorie de preuves appartiennent ces documents ? Préciser la méthode de datation utilisée pour les plus anciens.
\item En exploitant les données du tableau, montrer que ces documents apportent des preuves directes de l'évolution de la lignée humaine.
\end{enumerate}
\tcblower
\textbf{1. Nature des documents et datation.} Il s'agit de restes fossiles (crânes, mandibules, squelettes post-crâniens) découverts dans des couches sédimentaires datées : ce sont des preuves \cle{paléontologiques} directes de l'évolution humaine. Pour Toumaï (7 Ma) et Lucy (3,2 Ma), les datations reposent sur la méthode \textbf{potassium-argon} (\ce{^{40}K}/\ce{^{40}Ar}) appliquée aux couches volcaniques (tufs, basaltes) encaissant les fossiles, car la méthode au \ce{^{14}C} est inutilisable au-delà de \num{50000} ans.

\textbf{2. Exploitation des données : tendance évolutive.} On raisonne caractère par caractère, en comparant les formes anciennes à la forme actuelle.
\begin{itemize}
\item \textbf{Capacité crânienne (volume endocrânien) :} elle augmente de manière nette au cours du temps : \SI{350}{cm^3} chez Toumaï (7 Ma, valeur proche du chimpanzé actuel), \SI{450}{cm^3} chez Lucy (3,2 Ma), puis \SI{1350}{cm^3} en moyenne chez l'Homme actuel. Soit une multiplication par environ 3,9 entre Toumaï et \emph{Homo sapiens} ($1350/350 \approx 3,9$). Le développement progressif du cerveau (encéphalisation) est un fait évolutif mesurable et orienté.
\item \textbf{Position du trou occipital (foramen magnum) :} chez Toumaï, il est déjà en position antérieure sur la base du crâne, indice fort d'une tête portée au-dessus du rachis (station verticale) et d'une \cle{bipédie} facultative ou habituelle ; chez Lucy, il est en position intermédiaire, confirmant une bipédie acquise mais une base crânienne moins flexionnée ; chez l'Homme actuel, il est centré sous le crâne, position caractéristique de la bipédie permanente et de l'équilibre statique sur les membres inférieurs.
\end{itemize}
\textbf{Conclusion.} Les formes fossiles anciennes possèdent des caractères intermédiaires (mosaïque de caractères primitifs et dérivés) entre ceux des grands singes actuels (faible capacité crânienne, trou occipital postérieur) et ceux de l'Homme moderne (forte capacité crânienne, trou occipital centré, bipédie). L'apparition progressive et ordonnée dans le temps (7 Ma $\rightarrow$ 3,2 Ma $\rightarrow$ actuel) des caractères diagnostiques de la lignée humaine démontre que cette lignée a \cle{évolué} au cours des derniers millions d'années : les fossiles constituent donc des preuves directes, datées et stratifiées de l'évolution de l'Homme.
\end{exoresolu}

\coursec{II. Les critères de l'hominisation}

\courssub{Identifier les critères morphologiques et culturels de l'hominisation}

\begin{definition}[Hominisation]
L'\cle{hominisation} est l'ensemble des transformations morphologiques, physiologiques et culturelles ayant conduit, à partir d'un ancêtre commun partagé avec les chimpanzés, à l'espèce \emph{Homo sapiens}. Elle se définit par l'acquisition successive de quatre critères majeurs, dans un ordre chronologique précis.
\end{definition}

\begin{propriete}[Les quatre critères de l'hominisation (ordre d'apparition)]
\begin{enumerate}
\item \textbf{La bipédie (critère le plus ancien, acquis dès \SI{\sim 7}{Ma}) :} station verticale et locomotion sur les membres inférieurs. Indices squelettiques : trou occipital antérieur, colonne vertébrale à double courbure (cyphose thoracique, lordose lombaire), bassin court, large et évasé (ilions orientés latéro-antérieurement), fémur oblique (angle col-diaphysaire \SI{\sim 125}{\degree}) avec condyles inégaux, pied à voûte plantaire et gros orteil aligné (non opposable).
\item \textbf{L'augmentation de la capacité crânienne et la réduction du prognathisme (à partir de \SI{\sim 2,5}{Ma}, genre \emph{Homo}) :} passage de \SI{\sim 400-500}{cm^3} (\emph{Australopithecus}) à \SI{\sim 1350}{cm^3} (\emph{H. sapiens}). Réduction concomitante de la face (orthognathie), diminution de la taille des canines et des prémolaires, disparition du diastème.
\item \textbf{La libération de la main et la préhension fine (corrélée à la bipédie, exploitée par \emph{Homo}) :} main préhensile à pouce long et pleinement opposable, musculature thenar/hypothénar développée, permettant la fabrication et l'utilisation d'outils complexes (préhension de précision « pouce-index »).
\item \textbf{L'activité culturelle et technique (marqueur du genre \emph{Homo}) :} taille de la pierre (galets aménagés \emph{Oldowayen}, bifaces \emph{Acheuléen}), maîtrise du feu (\SI{\sim 400000}{ans} preuves solides, \SI{\sim 1}{Ma} indices discutés), sépultures intentionnelles (\emph{H. neanderthalensis}, \emph{H. sapiens}), art pariétal et mobilier, langage symbolique (\emph{H. sapiens}).
\end{enumerate}
\textbf{Point clé Bac} : La bipédie précède l'encéphalisation. \emph{Australopithecus} est bipède avec un cerveau de singe. L'évolution n'est pas linéaire mais \cle{buissonnante} : coexistence de plusieurs espèces (ex. \emph{Paranthropus} robustes et \emph{Homo} graciles vers 2-1 Ma).
\end{propriete}

\begin{attention}[Pièges classiques sur les critères]
\begin{itemize}
\item Ne pas écrire « l'Homme descend du singe » : Homme et grands singes actuels descendent d'un \cle{ancêtre commun} (proche de \emph{Sahelanthropus} ou \emph{Orrorin}) ; ce sont des cousins évolutifs.
\item Ne pas inverser l'ordre : \textbf{Bipédie $\rightarrow$ Cerveau $\rightarrow$ Main libre $\rightarrow$ Culture}. La libération de la main est une \emph{conséquence} de la bipédie, pas une cause.
\item Ne pas présenter l'évolution comme une échelle : \emph{Australopithecus} $\nrightarrow$ \emph{H. habilis} $\nrightarrow$ \emph{H. erectus} $\nrightarrow$ \emph{H. sapiens} sans branches latérales. C'est un buisson : \emph{Paranthropus}, \emph{H. neanderthalensis}, \emph{H. floresiensis}, \emph{H. naledi} sont des impasses évolutives.
\end{itemize}
\end{attention}

\coursec{III. Les grandes étapes de la lignée humaine}

\courssub{Reconstituer la lignée humaine et les étapes de l'hominisation}

\begin{methode}[Reconstituer les étapes de l'hominisation à partir de fossiles datés]
\begin{enumerate}
\item \cle{Repérer} dans les documents les indices des quatre critères d'hominisation (bipédie, capacité crânienne, main/outils, culture).
\item \cle{Associer} chaque ensemble de caractères à un genre/spèce représentatif en utilisant les datations :
\begin{itemize}
\item \textbf{6-4 Ma :} \emph{Sahelanthropus} (Tchad), \emph{Orrorin} (Kenya), \emph{Ardipithecus} (Éthiopie) — bipédie naissante, cerveau \SI{\sim 350}{cm^3}.
\item \textbf{4-2 Ma :} \emph{Australopithecus} (\emph{afarensis}, \emph{africanus}, \emph{sediba}) — bipédie acquise, cerveau \SI{400-500}{cm^3}, pas d'outils taillés.
\item \textbf{2,4-1,4 Ma :} \emph{Homo habilis} / \emph{rudolfensis} — premiers outils (Oldowayen : galets aménagés, éclats), cerveau \SI{600-750}{cm^3}.
\item \textbf{1,8 Ma - 300 ka :} \emph{Homo erectus} / \emph{ergaster} — bifaces (Acheuléen), sortie d'Afrique (\emph{H. georgicus} Dmanissi 1,8 Ma ; Java, Chine), feu (indices \SI{\sim 1}{Ma}, preuves \SI{\sim 400}{ka}), cerveau \SI{900-1100}{cm^3}.
\item \textbf{400-40 ka :} \emph{Homo neanderthalensis} (Europe/Proche-Orient) — sépultures, outillage Moustérien, cerveau \SI{\sim 1400}{cm^3} (supérieur à \emph{H. sapiens} actuel).
\item \textbf{300 ka - actuel :} \emph{Homo sapiens} (Afrique : Jebel Irhoud 300 ka) — art, langage symbolique, outillage complexe (Aurignacien), colonisation mondiale.
\end{itemize}
\item \cle{Ordonner} chronologiquement en précisant les coexistences (buissonnant).
\item \cle{Conclure} par une phrase de synthèse reliant acquisition de la bipédie (libération mains), développement cérébral (capacités cognitives) et essor culturel (adaptation technique).
\end{enumerate}
\end{methode}

\begin{exoresolu}[De la bipédie à la culture : séquence Afar]
Au cours d'une fouille dans la région de l'Afar (Éthiopie), on découvre dans une couche datée de 3,2 Ma (\ce{^{40}K}/\ce{^{40}Ar}) un squelette partiel dont le bassin est court, large, aux ilions évasés antérieurement ; le fémur présente un angle col-diaphysaire de \SI{125}{\degree} et des condyles inégaux ; le trou occipital est en position intermédiaire. La capacité crânienne est estimée à \SI{450}{cm^3}. Dans une couche supérieure datée de 1,6 Ma, on découvre un crâne de capacité \SI{950}{cm^3} associé à des bifaces symétriques et à des sédiments rubéfiés (indices de foyers).
\begin{enumerate}
\item Identifier le genre de chacun des deux spécimens en justifiant par deux arguments anatomiques ou culturels chacun.
\item Décrire, à partir de cet exemple, les grandes étapes de l'hominisation mises en évidence.
\end{enumerate}
\tcblower
\textbf{1. Identification des spécimens.}
\begin{itemize}
\item \textbf{Spécimen de 3,2 Ma (couche inférieure) :} Le bassin court/large (ilions évasés), le fémur oblique (angle \SI{125}{\degree}, condyles inégaux) et le trou occipital avancé sont des indices squelettiques indiscutables de la \cle{bipédie} habituelle. Or la capacité crânienne reste faible (\SI{450}{cm^3}, proche du chimpanzé). L'association « bipédie acquise + encéphale réduit + absence d'outils » à 3,2 Ma caractérise le genre \emph{Australopithecus} (spécimen type : Lucy, \emph{A. afarensis}).
\item \textbf{Spécimen de 1,6 Ma (couche supérieure) :} La capacité crânienne de \SI{950}{cm^3} (nettement supérieure au seuil \SI{\sim 700}{cm^3} du genre \emph{Homo}), associée à des outils élaborés (bifaces acheuléens = symétrie bilatérale, standardisation) et à des indices de maîtrise du feu (sédiments rubéfiés), caractérise \emph{Homo erectus} (ou \emph{H. ergaster} en Afrique).
\end{itemize}
\textbf{2. Étapes de l'hominisation mises en évidence.}
\begin{itemize}
\item \textbf{Étape 1 (avant 3,2 Ma) : Acquisition de la bipédie.} C'est le premier critère d'hominisation à apparaître, alors que le volume cérébral reste ancestral (\SI{\sim 450}{cm^3}). La bipédie précède donc l'encéphalisation. Elle libère les mains de la locomotion.
\item \textbf{Étape 2 (entre 3,2 et 1,6 Ma) : Augmentation majeure de la capacité crânienne.} Elle passe de \SI{450}{cm^3} à \SI{950}{cm^3} (+ \SI{500}{cm^3}, soit +110 \%), traduisant le développement cérébral dans le genre \emph{Homo} (apparition de \emph{H. habilis} vers 2,4 Ma, puis \emph{H. erectus}).
\item \textbf{Étape 3 : Développement technique et culturel.} La bipédie ayant libéré les mains, celles-ci deviennent des organes de fabrication : apparition des outils taillés (Oldowayen \SI{\sim 2,6}{Ma}, puis Acheuléen \SI{\sim 1,7}{Ma}), maîtrise du feu (\emph{H. erectus}). La culture devient le mode d'adaptation principal.
\end{itemize}
\textbf{Conclusion.} L'hominisation se déroule selon une succession ordonnée et non linéaire : \textbf{bipédie d'abord} (libération mains, économie énergétique), \textbf{puis accroissement du volume cérébral} (capacités cognitives, planification), \textbf{enfin essor des activités techniques et culturelles} (outils, feu, langage) qui permettent une adaptation rapide aux milieux variés. La libération de la main par la bipédie a été la condition \emph{sine qua non} du développement technique, lui-même rétroactivement lié à l'augmentation des capacités cérébrales (boucle de rétroaction positive).
\end{exoresolu}

\coursec{IV. Parenté entre l'Homme et les grands singes anthropomorphes : l'arbre phylogénétique}

\courssub{Construire et lire un arbre phylogénétique}

\begin{methode}[Construire et interpréter un arbre phylogénétique (cladistique)]
\begin{enumerate}
\item \cle{Recueillir} les caractères partagés (anatomiques : squelette, dentition ; moléculaires : pourcentage d'identité des séquences d'ADN ou d'acides aminés de protéines conservatrices comme l'hémoglobine, le cytochrome c).
\item \cle{Regrouper} d'abord les espèces partageant le plus grand nombre de caractères dérivés (synapomorphies) ou le plus fort pourcentage d'identité moléculaire : elles forment un \cle{nœud} qui représente leur \cle{ancêtre commun} hypothétique le plus récent.
\item \cle{Joindre} progressivement les groupes en remontant vers les ancêtres communs plus anciens ; chaque nœud correspond à une \cle{divergence} (séparation de deux lignées sœurs).
\item \cle{Placer} une échelle de temps (en millions d'années \textbf{Ma}) si les datations moléculaires (horloge moléculaire) ou fossiles sont connues : divergence Homme--Chimpanzé \SI{\sim 6-7}{Ma} ; divergence Gorille \SI{\sim 8-10}{Ma} ; divergence Orang-outan \SI{\sim 12-14}{Ma}.
\item \cle{Lire} l'arbre : deux espèces sont d'autant plus proches parentes que leur nœud commun est récent (proche des extrémités). \textbf{Jamais} un arbre ne montre qu'une espèce actuelle descend d'une autre espèce actuelle (les nœuds sont des ancêtres éteints).
\end{enumerate}
\textbf{Réflexe} : vérifier la cohérence avec les données moléculaires de référence : Homme--Chimpanzé \SI{\sim 98,8}{\%} identité ADN ; Homme--Gorille \SI{\sim 98,3}{\%} ; Homme--Orang-outan \SI{\sim 97,0}{\%}.
\end{methode}

\begin{exoresolu}[Arbre phylogénétique Homme--Grands singes d'après l'ADN]
Le tableau ci-dessous indique le pourcentage d'identité des séquences d'ADN nucléaires entre l'Homme et trois grands singes anthropomorphes :
\begin{center}
\begin{tabular}{|l|c|c|c|}
\hline
\rowcolor{popblueL} \textbf{Espèce comparée à \emph{Homo sapiens}} & Chimpanzé (\emph{Pan troglodytes}) & Gorille (\emph{Gorilla gorilla}) & Orang-outan (\emph{Pongo pygmaeus}) \\
\hline
Identité des séquences d'ADN (\%) & 98,8 & 98,3 & 97,0 \\
\hline
\end{tabular}
\end{center}
\begin{enumerate}
\item Construire l'arbre phylogénétique de ces quatre espèces en indiquant les nœuds de divergence.
\item Indiquer le grand singe le plus proche parent de l'Homme et justifier par la donnée numérique.
\item Préciser l'âge estimé de la divergence Homme--Chimpanzé.
\end{enumerate}
\tcblower
\textbf{1. Construction de l'arbre.} Principe : plus le pourcentage d'identité est élevé, plus la divergence est récente (horloge moléculaire). On regroupe donc par ordre décroissant d'identité :
\begin{itemize}
\item Groupe 1 (plus proche) : Homme + Chimpanzé (98,8 \%) $\rightarrow$ Nœud 1 (ancêtre commun Homme-Chimpanzé).
\item Groupe 2 : (Homme+Chimpanzé) + Gorille (98,3 \%) $\rightarrow$ Nœud 2 (ancêtre commun Homininés-Gorillini).
\item Groupe 3 (racine) : Ensemble précédent + Orang-outan (97,0 \%) $\rightarrow$ Nœud 3 (ancêtre commun Hominidés).
\end{itemize}
\begin{popfigure}
\begin{center}
\begin{tikzpicture}[
  grow=right,
  level distance=2.8cm,
  level 1/.style={sibling distance=3.5cm},
  level 2/.style={sibling distance=1.8cm},
  level 3/.style={sibling distance=1.2cm},
  edge from parent/.style={draw=popdark,thick,->},
  every node/.style={font=\small, align=center},
  tip/.style={circle, draw=popblue, thick, fill=popblueL, inner sep=2pt},
  anc/.style={circle, draw=poporange, thick, fill=poporangeL, inner sep=2pt}
]
\node[anc] (root) {}
  child { node[tip, align=center]{Orang-outan\\ (\emph{Pongo})} edge from parent node[above, sloped, pos=0.6] {$\sim$ 13 Ma} }
  child { node[anc] (n2) {}
    child { node[tip, align=center]{Gorille\\ (\emph{Gorilla})} edge from parent node[above, sloped, pos=0.6] {$\sim$ 9 Ma} }
    child { node[anc] (n1) {}
      child { node[tip, fill=popgreenL, draw=popgreen, align=center]{\textbf{Homme}\\ (\emph{Homo sapiens})} }
      child { node[tip, fill=popgreenL, draw=popgreen, align=center]{Chimpanzé\\ (\emph{Pan})} }
      edge from parent node[above, sloped, pos=0.6] {$\sim$ 7 Ma}
    }
    edge from parent node[below, sloped, pos=0.6] {$\sim$ 9 Ma}
  };
% Labels des nœuds
\node[above=0.2cm of root, font=\tiny, text=poporange] {Nœud 3 : Ancêtre commun Hominidés};
\node[above=0.2cm of n2, font=\tiny, text=poporange] {Nœud 2 : Ancêtre commun Homininés/Gorillini};
\node[above=0.2cm of n1, font=\tiny, text=popgreen] {Nœud 1 : Ancêtre commun Homme/Chimpanzé};
\end{tikzpicture}
\end{center}
\legende{Arbre phylogénétique simplifié (cladogramme) de l'Homme et des grands singes anthropomorphes. Les nœuds (cercles orange/vert) représentent les ancêtres communs hypothétiques. Les flèches indiquent le sens du temps (passé $\rightarrow$ présent). L'échelle de temps (Ma) est basée sur l'horloge moléculaire calibrée par les fossiles.}
\end{popfigure}
\textbf{2. Plus proche parent.} Le \cle{chimpanzé} (\emph{Pan troglodytes} et \emph{Pan paniscus}) est le grand singe le plus proche de l'Homme, car c'est avec lui que l'Homme partage le plus fort pourcentage d'identité moléculaire (\textbf{98,8 \%}), donc l'ancêtre commun le plus récent (Nœud 1, \SI{\sim 6-7}{Ma}).
\textbf{3. Divergence.} La séparation des lignées humaine et chimpanzé est estimée entre \textbf{6 et 7 millions d'années} (fossiles : \emph{Sahelanthropus} 7 Ma, \emph{Orrorin} 6 Ma).
\textbf{Conclusion.} Les données moléculaires confirment les données anatomiques : l'Homme et le chimpanzé forment un \textbf{groupe frère} (clade Homininés), issu d'un ancêtre commun ayant vécu il y a environ 6 à 7 millions d'années ; l'Homme ne descend pas du chimpanzé actuel, mais partage avec lui une histoire évolutive commune très longue.
\end{exoresolu}

\courssub{Établir le lien de parenté entre l'Homme et les grands singes}

\begin{methode}[Argumenter rigoureusement la parenté Homme--Anthropomorphes]
Pour établir le lien de parenté entre l'Homme et les grands singes (chimpanzé, gorille, orang-outan, gibbon) :
\begin{enumerate}
\item \cle{Mobiliser} plusieurs ordres de preuves \textbf{convergentes} (indépendantes) :
\begin{itemize}
\item \textbf{Anatomiques :} même plan d'organisation du squelette (membres pentadactyles, clavicules, omoplates dorsales), absence de queue (apomorphie des Hominoïdes), dentition de même type (hétérodonte, diphyodonte, formule dentaire 2.1.2.3 / 2.1.2.3 = 32 dents chez Homme et singes africains), encéphalisation relative.
\item \textbf{Embryologiques :} grandes similitudes des stades précoces (vésicule optique, arcs branchiaux, queue embryonnaire résorbée, membres en bourgeons).
\item \textbf{Moléculaires :} forte identité des séquences d'ADN (gènes codants, ADN non codant, ERV -- rétrovirus endogènes partagés aux mêmes loci) et des protéines (hémoglobine $\alpha$ et $\beta$ quasi identiques Homme/Chimpanzé, cytochrome c).
\item \textbf{Caryotypiques :} caryotypes très proches (46 chromosomes chez l'Homme, 48 chez les grands singes) ; bandeamento G/R identique bande par bande, prouvant la fusion chromosomique humaine.
\end{itemize}
\item \cle{Interpréter} : des similitudes nombreuses, précises et portant sur des caractères indépendants (structurels, séquentiels, développementaux) ne s'expliquent pas par le hasard ni par convergence adaptative, mais par l'héritage d'un \cle{ancêtre commun}.
\item \cle{Nuancer} : les différences majeures (bipédie obligatoire, volume cérébral triplé, langage articulé, culture cumulative) résultent d'évolutions divergentes (sélection directionnelle) depuis la séparation des lignées.
\item \cle{Conclure} : l'Homme appartient au clade des \textbf{Hominidés} (grands singes + homme) et au sous-clade des \textbf{Homininés} (lignée humaine post-divergence chimpanzé) ; il est le \textbf{cousin} des grands singes actuels, non leur descendant.
\end{enumerate}
\end{methode}

\begin{exoresolu}[Le caryotype, témoin cytogénétique de la parenté]
L'Homme possède 46 chromosomes (23 paires) alors que le chimpanzé, le gorille et l'orang-outan en possèdent 48 (24 paires). L'étude des bandes chromosomiques (bandes G) montre que le chromosome 2 humain correspond, bout à bout, à deux chromosomes distincts du chimpanzé (chimpanzé chr 12 et 13, orthologues de l'ancêtre), avec au milieu des restes de séquences télomériques (répétitions \texttt{TTAGGG}) et deux centromères (un actif, un inactif vestigial).
\begin{enumerate}
\item Expliquer l'origine du chromosome 2 humain à partir du caryotype ancestral.
\item En quoi cette observation renforce-t-elle l'hypothèse de la parenté étroite entre l'Homme et les grands singes ?
\end{enumerate}
\tcblower
\textbf{1. Origine du chromosome 2 humain.} Dans la lignée humaine, \textbf{après} sa séparation d'avec la lignée du chimpanzé (donc il y a moins de 7 Ma), deux chromosomes acrocentriques de l'ancêtre commun (orthologues aux chromosomes 12 et 13 du chimpanzé actuel) ont subi une \cle{fusion centrique (de type Robertsonnienne)} bout à bout (télomère-télomère). La présence de séquences télomériques résiduelles \emph{interstitielles} au milieu du chromosome 2 humain — alors que les télomères se situent normalement aux extrémités — est la signature moléculaire irréfutable de cette fusion. La présence d'un centromère vestigial inactif sur le bras long confirme l'origine bicentromérique. Cela explique le passage de 48 à 46 chromosomes : $48 - 2 + 1 = 47$ paires potentielles $\rightarrow$ 23 paires (46 chromosomes) après fixation dans la population.
\textbf{2. Renforcement de l'hypothèse de parenté.} La correspondance quasi parfaite, \textbf{bande par bande} (bandes G), entre le chromosome 2 humain et la juxtaposition des deux chromosomes du chimpanzé montre que les deux espèces possèdent le \emph{même patrimoine génétique}, organisé de manière presque identique. Une telle similitude structurale (syntenie conservée) ne peut s'expliquer que par l'héritage d'un \cle{ancêtre commun} dont le caryotype comportait 48 chromosomes, conservé tel quel chez le chimpanzé (et gorille, orang-outan) et modifié par une unique fusion chez l'Homme.
\textbf{Conclusion.} Le caryotype constitue une preuve cytogénétique majeure de la parenté étroite entre l'Homme et les grands singes : les différences observées sont des réarrangements chromosomiques récents (fusion, inversions péricentriques mineures) survenus dans un fond chromosomique commun, ce qui confirme sans ambiguïté les données anatomiques, embryologiques et moléculaires de séquence.
\end{exoresolu}

\coursec{V. Capacité d'adaptation de l'Homme à son environnement et développement durable}

\courssub{Relier évolution humaine, adaptation culturelle et développement durable}

\begin{methode}[Construire un raisonnement structuré : Évolution $\rightarrow$ Adaptation $\rightarrow$ Développement Durable]
\begin{enumerate}
\item \cle{Rappeler} que l'évolution biologique de l'Homme (bipédie, libération mains, encéphalisation) est elle-même une réponse adaptative aux changements environnementaux passés (ex. : ouverture des milieux forestiers $\rightarrow$ extension des savanes africaines au Miocène/Pliocène favorisant la bipédie et la thermorégulation).
\item \cle{Montrer} que, contrairement aux autres espèces dont l'adaptation repose sur l'évolution biologique (sélection naturelle, temps long, générations), l'Homme s'adapte \textbf{principalement par la culture et la technique} (temps court, transmission horizontale/verticale) : vêtements, habitat, agriculture, élevage, médecine, énergie, transformation des écosystèmes. Il \emph{transforme activement son milieu} (niche écologique) au lieu de seulement le subir (construction de niche / \emph{niche construction}).
\item \cle{Identifier} les conséquences \textbf{négatives} (maladaptation à long terme) de cette capacité technique débridée : déforestation massive (bassin du Congo, Amazonie, Asie), érosion des sols, pollution (plastiques, métaux lourds, azote/phosphore), émissions de gaz à effet de serre (GES) $\rightarrow$ changement climatique, érosion de la biodiversité (6\textsuperscript{e} extinction de masse), épuisement des ressources halieutiques, forestières, minières. \textbf{Exemples camerounais :} déforestation pour cacao/palmier à huile/bois énergie (régions Sud, Est, Centre), urbanisation rapide non maîtrisée (Douala, Yaoundé) $\rightarrow$ inondations, glissements de terrain, perte de biodiversité ; érosion éolienne/hydrique au Nord (Sahel) ; pollution plastique (fleuves Wouri, Sanaga).
\item \cle{Dégager} la nécessité d'une adaptation \textbf{raisonnée et éthique} : le \cle{développement durable} (Brundtland 1987), c'est-à-dire un « développement qui répond aux besoins du présent sans compromettre la capacité des générations futures à répondre aux leurs ». Trois piliers indissociables : \textbf{Économique} (efficacité), \textbf{Social} (équité), \textbf{Environnemental} (viabilité).
\item \cle{Proposer} des solutions concrètes, contextualisées (Cameroun) et mesurables :
\begin{itemize}
\item \textbf{Forêt/Biodiversité :} reboisement (opération « Green Sahel »), agroforesterie (cacaoyers sous couvert forestier, systèmes agroforestiers à \emph{Faidherbia albida} au Nord), aires protégées gérées participativement, lutte contre le braconnage.
\item \textbf{Énergie/Climat :} développement hydroélectrique (barrages Lom Pangar, Nachtigal, Memve'ele), solaire photovoltaïque (rural), foyers améliorés (réduction bois/charbon), biogaz (déchets agricoles).
\item \textbf{Sols/Agriculture :} agriculture de conservation (semis direct sous couverture végétale), rotation cultures-légumineuses, compostage, lutte anti-érosive (cordons pierreux, haies vives).
\item \textbf{Déchets/Urbanisme :} tri sélectif, valorisation déchets organiques (compost), plastiques (recyclage), assainissement durable, planification urbaine (zones non constructibles, corridors écologiques).
\item \textbf{Gouvernance/Éducation :} éducation à l'environnement (programmes scolaires, ONG), application stricte loi forestière (1994) et loi sur l'environnement (1996), évaluation d'impact environnemental (EIE) obligatoire.
\end{itemize}
\end{enumerate}
\textbf{Réflexe de rédaction Bac} : structurer la réponse en \textbf{Constat (données chiffrées)} $\rightarrow$ \textbf{Mécanisme (lien évolution/adaptation)} $\rightarrow$ \textbf{Conséquence (maladaptation)} $\rightarrow$ \textbf{Solution (DD : mesure concrète + pilier concerné)}.
\end{methode}

\begin{exoresolu}[Déforestation, érosion et adaptation durable au Cameroun]
Entre 1990 et 2020, le Cameroun a perdu environ \SI{15}{\%} de son couvert forestier dense humide (passant de \SI{\sim 22}{millions d'hectares} à \SI{\sim 18,5}{millions d'hectares}), sous l'effet conjugué de l'agriculture itinérante sur brûlis (principal moteur), de l'exploitation forestière industrielle (bois d'œuvre), de la demande en bois énergie (\SI{> 80}{\%} des ménages) et de l'extension des plantations industrielles (palmier à huile, hévéa, cacao). Parallèlement, les épisodes de sécheresse pluviométrique et d'érosion hydrique/éolienne des sols se sont aggravés dans les régions septentrionales (Adamaoua, Nord, Extrême-Nord), et les inondations urbaines sont récurrentes à Douala et Yaoundé.
\begin{enumerate}
\item Expliquer en quoi ces activités humaines traduisent une capacité d'adaptation spécifique de l'Homme à son environnement, issue de son histoire évolutive.
\item Montrer que cette adaptation peut devenir \emph{maladaptive} (non durable) et proposer \textbf{trois} mesures concrètes de développement durable adaptées au contexte camerounais, en les rattachant aux trois piliers (Économique, Social, Environnemental).
\end{enumerate}
\tcblower
\textbf{1. Capacité d'adaptation culturelle issue de l'hominisation.} Contrairement aux autres espèces dont l'adaptation repose uniquement sur l'évolution biologique (sélection naturelle agissant sur le génome sur des milliers de générations), l'Homme s'adapte par la \cle{technique} et la \cle{culture} (transmission extra-génétique, cumulatif). Il défriche la forêt pour cultiver (agriculture vivrière : manioc, maïs, plantain ; cultures de rente : cacao, café, coton), exploite le bois pour se chauffer, construire et exporter, aménage des plantations pour produire des revenus monétaires. Il \emph{transforme activement son milieu} (construction de niche) pour satisfaire ses besoins alimentaires, énergétiques et économiques immédiats : c'est l'aboutissement direct de l'hominisation (main libérée par la bipédie $\rightarrow$ fabrication d'outils $\rightarrow$ cerveau développé $\rightarrow$ langage/planification $\rightarrow$ culture cumulative). Cette adaptation est puissante et rapide.

\textbf{2. Adaptation devenue maladaptive et solutions de développement durable.} Cette transformation devient \textbf{maladaptive} (non viable) lorsqu'elle détruit les conditions mêmes de la survie future (capacité porteuse de l'écosystème) :
\begin{itemize}
\item \textbf{Environnemental :} déforestation $\rightarrow$ perte d'habitat/biodiversité, érosion sols (ruissellement nu), baisse évapotranspiration $\rightarrow$ pluviométrie locale réduite, libération carbone stocké $\rightarrow$ aggravation effet de serre.
\item \textbf{Social :} perte services écosystémiques (eau, bois, plantes médicinales, protéines brousse) pour populations rurales/autochtones (Baka, Bagyeli) ; conflits fonciers ; insécurité alimentaire à long terme (sol lessivé).
\item \textbf{Économique :} perte capital naturel (bois d'œuvre, services écosystémiques non valorisés) ; coûts dégâts inondations/érosion ; dépendance exportations brutes (volatilité cours).
\end{itemize}
Trois mesures de développement durable (DD) contextualisées :
\begin{enumerate}
\item \textbf{Agroforesterie cacaoyère et systèmes à \emph{Faidherbia albida} (Pilier Environnemental + Économique).} Associer arbres d'ombrage (espèces forestières locales, \emph{Faidherbia} au Nord) et cultures (cacao, café, coton, mil). Maintient fertilité sols (fixation azote, litière), séquestre carbone, diversifie revenus (bois, fruits, cacao), réduit intrants chimiques. Exemple : projets « Cacao durable » (certification Rainforest Alliance/UTZ) dans le Sud/Centre.
\item \textbf{Transition énergétique : foyers améliorés + solaire rural + hydroélectricité (Pilier Environnemental + Social + Économique).} Diffusion massive de foyers à haut rendement (réduction 40-50\% consommation bois/charbon $\rightarrow$ déforestation, pollution air intérieur santé femmes/enfants). Électrification rurale solaire (kits PV, mini-réseaux) pour pompage eau, transformation produits, éclairage (éducation). Achèvement barrages (Nachtigal 420 MW, Memve'ele 200 MW) pour substituer thermique (fioul) et alimenter industrie/urbain.
\item \textbf{Gestion intégrée des déchets urbains et économie circulaire (Pilier Social + Économique + Environnemental).} Mise en place tri à la source (ménages, marchés) dans les grandes villes (Douala, Yaoundé, Bafoussam). Unités de compostage déchets organiques (engrais agriculture périurbaine). Filières recyclage plastiques (pavés, granulés), métaux. Valorisation biogaz (déchets abattoirs, marchés). Création emplois verts (jeunes, femmes). Réduction pollution fleuves (Wouri, Mfoundi) et océans.
\end{enumerate}
\textbf{Conclusion.} L'histoire évolutive a doté l'Homme d'une capacité d'adaptation culturelle unique (technique, langage, coopération), lui permettant de coloniser tous les biomes. Mais cette puissance technique, si elle n'est pas encadrée par une éthique de responsabilité intergénérationnelle, détruit le capital naturel dont elle dépend. Seule une exploitation \textbf{raisonnée, équitable et régénérative} des ressources, conforme aux principes du \cle{développement durable}, garantira l'adaptation \textbf{perpétuelle} de l'espèce humaine à son environnement planétaire.
\end{exoresolu}

\begin{attention}[Les 5 erreurs qui coûtent des points au Bac (SVT Terminale D)]
\begin{enumerate}
\item \textbf{Écrire que « l'Homme descend du singe » ou « l'Homme descend du chimpanzé ».} \textbf{Faux et sanctionné (0 pt sur la question).} L'Homme et les grands singes actuels descendent d'un \cle{ancêtre commun} (proche de \emph{Sahelanthropus}) ; ce sont des \textbf{cousins évolutifs}. Une espèce actuelle ne descend \textbf{jamais} d'une autre espèce actuelle.
\item \textbf{Présenter l'évolution humaine comme une échelle linéaire} (\emph{Australopithecus} $\rightarrow$ \emph{Homo habilis} $\rightarrow$ \emph{Homo erectus} $\rightarrow$ \emph{Homo sapiens} sans nuance). L'évolution est \cle{buissonnante} (arbre) : plusieurs espèces d'hominines ont coexisté (ex. \emph{Paranthropus boisei} et \emph{Homo habilis} à 1,8 Ma ; \emph{Homo neanderthalensis} et \emph{Homo sapiens} à 40 ka ; \emph{Homo floresiensis} jusqu'à 50 ka).
\item \textbf{Inverser l'ordre chronologique des critères d'hominisation.} La \cle{bipédie} apparaît \textbf{avant} l'augmentation de la capacité crânienne (Lucy, 3,2 Ma, \SI{450}{cm^3}, était déjà bipède ; Toumaï, 7 Ma, bipède probable). Ne jamais écrire que « le gros cerveau a permis la bipédie » ou « les outils ont libéré la main ».
\item \textbf{Citer une donnée chiffrée sans l'exploiter / sans la comparer.} Écrire « la capacité crânienne est de \SI{900}{cm^3} » ne suffit pas : il faut la \emph{comparer} explicitement (à \SI{450}{cm^3} chez \emph{Australopithecus}, à \SI{1350}{cm^3} chez \emph{H. sapiens}) et en tirer la conclusion évolutive (augmentation, encéphalisation). Toute preuve doit être \textbf{argumentée} (donnée + comparaison + interprétation).
\item \textbf{Confondre adaptation biologique et adaptation culturelle, et oublier de définir le développement durable.} L'adaptation humaine moderne est \textbf{majoritairement technique et culturelle} (rapide, non génétique) ; le développement durable doit être \textbf{défini} (citation Brundtland : besoins présent / générations futures) et \textbf{illustré par des mesures concrètes, chiffrées et localisées} (pas de généralités vagues type « protéger la nature »).
\end{enumerate}
\end{attention}

\begin{savaistu}[Le Cameroun, berceau de l'humanité ?]
Le Tchad (Toros-Menalla), limitrophe du Cameroun au Nord, a livré \emph{Sahelanthropus tchadensis} (Toumaï, \SI{\sim 7}{Ma}), plus ancien représentant connu de la lignée humaine. Au Cameroun même, les sites de \textbf{Shum Laka} (Nord-Ouest, \SI{\sim 30000}{ans}) et \textbf{Mbi Crater} (Ouest) ont livré des sépultures et de l'ADN ancien révélant une structure génétique complexe des populations africaines du Pléistocène supérieur. La forêt du bassin du Congo, dont le Sud-Cameroun est une partie intacte, est un refuge climatique majeur ayant permis la survie de lignées anciennes (gorilles, chimpanzés, humains). Étudier l'évolution de l'Homme, c'est aussi étudier l'histoire géologique et climatique de l'Afrique centrale.
\end{savaistu}

\begin{exercice}[Entraînement Bac : Preuves et Phylogénie]
\begin{enumerate}
\item On compare les séquences d'une protéine conservée (cytochrome c) chez quatre espèces. Le nombre de différences d'acides aminés par rapport à l'Homme est : Chimpanzé 0 ; Gorille 1 ; Orang-outan 3 ; Macaque 9.
\begin{enumerate}
\item Construire l'arbre phylogénétique correspondant.
\item Quel est le groupe frère de l'Homme ? Justifier.
\end{enumerate}
\item Un fossile d'hominine découvert au Kenya est daté de 6 Ma par la méthode \ce{^{40}K}/\ce{^{40}Ar}. Son fémur présente un angle col-diaphysaire de \SI{120}{\degree} et un trou occipital antérieur. Sa capacité crânienne est de \SI{360}{cm^3}.
\begin{enumerate}
\item Quel critère d'hominisation ce fossile illustre-t-il principalement ?
\item Proposer un genre possible pour ce fossile.
\item Expliquer pourquoi la méthode au \ce{^{14}C} n'a pas pu être utilisée.
\end{enumerate}
\end{enumerate}
\end{exercice}

\begin{corrige}[Exercice Entraînement Bac]
\textbf{1. Arbre phylogénétique du cytochrome c.}
\begin{itemize}
\item[a)] Différences : Homme-Chimpanzé (0) < Homme-Gorille (1) < Homme-Orang-outan (3) < Homme-Macaque (9).
\item Regroupement : (Homme, Chimpanzé) $\rightarrow$ Nœud 1 ; + Gorille $\rightarrow$ Nœud 2 ; + Orang-outan $\rightarrow$ Nœud 3 ; + Macaque (groupe externe) $\rightarrow$ Racine.
\item Arbre : ((Homme, Chimpanzé), Gorille), Orang-outan), Macaque.
\item[b)] Le \textbf{Chimpanzé} est le groupe frère de l'Homme (0 différence = identité protéique totale pour cette protéine = divergence la plus récente).
\end{itemize}
\textbf{2. Fossile kényan 6 Ma.}
\begin{itemize}
\item[a)] Critère principal : la \textbf{bipédie} (angle col-diaphysaire \SI{120}{\degree} proche valeur humaine \SI{125}{\degree} + trou occipital antérieur). C'est le critère le plus ancien.
\item[b)] Genre possible : \emph{Orrorin tugenensis} (découvert au Kenya, \SI{\sim 6}{Ma}, fémur bipède, dents petites) ou \emph{Sahelanthropus tchadensis} (proche géographiquement, \SI{\sim 7}{Ma}, trou occipital antérieur). \emph{Ardipithecus} (\SI{5,5-4,4}{Ma}) est plus récent.
\item[c)] La méthode au \ce{^{14}C} ne remonte qu'à \SI{\sim 50000}{ans} (demi-vie \SI{5730}{ans}). À 6 Ma (\SI{6 000 000}{ans}), tout le \ce{^{14}C} initial a disparu depuis longtemps (plus de 1000 demi-vies). Il faut une méthode à longue demi-vie : \ce{^{40}K}/\ce{^{40}Ar} (demi-vie \SI{1,25}{milliard d'années}) sur les couches volcaniques encaissantes.
\end{itemize}
\end{corrige}

\newpage

\coursec{Exercices}

\courssub{Je vérifie mes connaissances}

\begin{exercice}[QCM : Preuves et critères de l'hominisation]
Parmi les affirmations suivantes, cochez la ou les bonne(s) réponse(s). Justifiez brièvement votre choix.
\begin{enumerate}
    \item Les principales preuves de l'évolution de l'Homme reposent sur :\\
    a) Les seuls fossiles découverts en Europe.\\
    b) La paléontologie (fossiles), l'anatomie comparée, la biochimie (protéines, ADN) et la datation absolue.\\
    c) Les récits mythologiques et les textes anciens.\\
    d) L'observation exclusive des grands singes actuels.
    \item La bipédie stricte (obligatoire) se distingue de la bipédie facultative par :\\
    a) La position du trou occipital (foramen magnum) centré sous le crâne.\\
    b) L'absence de bassin.\\
    c) Un volume crânien supérieur à \SI{1400}{cm^3}.\\
    d) La présence d'un torse supra-orbitaire très développé.
    \item Le genre \emph{Australopithecus} (ex. \emph{A. afarensis}, « Lucy ») présente :\\
    a) Un volume crânien moyen de \SI{1350}{cm^3} et une bipédie parfaite.\\
    b) Un volume crânien proche de \SI{450}{cm^3}, une bipédie obligatoire mais une morphologie crânienne proche des singes (prognathisme, crête sagittale chez les mâles).\\
    c) Une taille corporelle supérieure à \SI{1,80}{m} et des outils en pierre taillée.\\
    d) Une absence totale de poils et une parole articulée.
    \item L'hominisation est un processus :\\
    a) Linéaire, sans branchement, menant directement à \emph{Homo sapiens}.\\
    b) Buissonnant, avec coexistence de plusieurs espèces d'homininés (ex. \emph{Paranthropus}, \emph{Homo} ancien).\\
    c) Qui s'est déroulé exclusivement en Asie il y a \SI{2}{Ma}.\\
    d) Où l'augmentation du volume crânien a précédé l'acquisition de la bipédie.
    \item La fusion chromosomique à l'origine du chromosome 2 humain (2n=46) alors que les grands singes en ont 48 (2n=48) est :\\
    a) Une preuve de spéciation par polyploïdie.\\
    b) Un argument majeur de parenté commune (synapomorphie partagée issue d'un ancêtre commun).\\
    c) Une erreur de méiose sans conséquence évolutive.\\
    d) Un caractère dérivé propre aux gorilles.
\end{enumerate}
\end{exercice}

\begin{exercice}[Vrai ou Faux justifié : Lignée humaine et parentés]
Pour chaque affirmation, indiquez VRAI ou FAUX et justifiez par une phrase précise en mobilisant vos connaissances.
\begin{enumerate}
    \item \emph{Sahelanthropus tchadensis} (Toumaï), découvert au Tchad, daté d'environ \SI{7}{Ma}, est considéré comme le plus ancien représentant connu de la lignée humaine après la divergence avec les chimpanzés.
    \item Le volume crânien a augmenté de manière régulière et continue de \emph{Australopithecus} à \emph{Homo sapiens} sans aucune phase de stagnation.
    \item L'Homme moderne (\emph{Homo sapiens}) et l'Homme de Néandertal (\emph{Homo neanderthalensis}) ont coexisté en Europe et au Proche-Orient il y a environ \SI{40000}{ans} et se sont hybridés.
    \item L'arbre phylogénétique basé sur l'ADN mitochondrial place l'Homme comme groupe frère du gorille, le chimpanzé étant plus éloigné.
\end{enumerate}
\end{exercice}

\begin{exercice}[Définitions et mise en relation]
Associez chaque concept à sa définition exacte, puis reliez les concepts entre eux par une flèche logique (ex. \emph{Critère} $\rightarrow$ \emph{Preuve}).
\begin{center}
\begin{tabularx}{\linewidth}{|l|X|}\hline
\textbf{Concept} & \textbf{Définition} \\ \hline
Hominisation & Ensemble des caractères morphologiques et comportementaux qui distinguent la lignée humaine de celle des grands singes (bipédie, volume crânien, face orthognathe, main préhensile, culture). \\ \hline
Bipédie obligatoire & Mode de locomotion exclusif sur les deux membres inférieurs, libérant les mains, lié à la réorganisation du bassin, du fémur et de la position du trou occipital. \\ \hline
Arbre phylogénétique & Représentation graphique des relations de parenté entre espèces, basée sur des caractères partagés dérivés (synapomorphies), matérialisant l'évolution divergente. \\ \hline
Développement durable & Mode de développement qui répond aux besoins du présent sans compromettre la capacité des générations futures à répondre aux leurs (environnement, social, économie). \\ \hline
Fusion de Robertson & Réarrangement chromosomique par fusion centrique de deux chromosomes acrocentriques en un seul chromosome métacentrique, réduisant le nombre chromosomique (48 $\rightarrow$ 46). \\ \hline
\end{tabularx}
\end{center}
\end{exercice}

\begin{exercice}[Ordre chronologique et datations absolues]
Le tableau ci-dessous présente cinq fossiles clés de la lignée humaine avec leur datation approximative (en millions d'années, Ma). Classez-les par ordre chronologique croissant (du plus ancien au plus récent) et indiquez le genre et l'espèce correspondants.
\begin{center}
\begin{tabularx}{\linewidth}{|X|c|X|}\hline
\textbf{Fossile (Surnom / Site)} & \textbf{Datation (Ma)} & \textbf{Genre / Espèce (à compléter)} \\ \hline
Lucy (Hadar, Éthiopie) & 3,2 & \\ \hline
Toumaï (Toros-Menalla, Tchad) & $\sim$ 7 & \\ \hline
Homme de Tautavel (Arago, France) & $\sim$ 0,45 & \\ \hline
Turkana Boy (Nariokotome, Kenya) & 1,6 & \\ \hline
Cro-Magnon (Abri de Cro-Magnon, France) & $\sim$ 0,03 & \\ \hline
\end{tabularx}
\end{center}
\emph{Question bonus :} Quelle méthode de datation absolue a permis de dater le tuf volcanique recouvrant le site de Lucy (potassium-argon ou carbone 14) ? Justifiez.
\end{exercice}

\courssub{Je m'entraîne}

\begin{exercice}[Analyse de documents : Critères de l'hominisation sur crânes]
On vous présente les mesures relevées sur trois crânes fossiles (Doc 1, 2, 3) et un crâne de chimpanzé actuel (Doc 0) pour comparaison.

\begin{popfigure}
\begin{tikzpicture}[scale=0.7, every node/.style={font=\small}]
% Chimpanzee skull schematic (left view)
\draw[popdark, thick] (0,0) ellipse (1.5 and 2); % braincase
\draw[popdark, thick] (-1.5,0) -- (-3.5, -0.5) -- (-3.5, 1.5) -- (-1.5, 1.5) -- cycle; % face/prognathism
\draw[popdark, thick] (-1.5, -1.5) arc (180:360:1.5 and 0.5); % jaw
\node[popdark] at (-2.5, 2.5) {Doc 0 : \emph{Pan troglodytes}};
% Foramen magnum position
\fill[popred] (-0.8, -1.8) circle (2pt) node[right] {FM postérieur};
% Australopithecus afarensis
\begin{scope}[xshift=7cm]
\draw[popblue, thick] (0,0) ellipse (1.8 and 2.2);
\draw[popblue, thick] (-1.8,0) -- (-3.2, -0.2) -- (-3.2, 1.2) -- (-1.8, 1.2) -- cycle;
\draw[popblue, thick] (-1.8, -1.5) arc (180:360:1.8 and 0.5);
\node[popblue] at (-1.5, 2.8) {Doc 1 : \emph{A. afarensis}};
\fill[popred] (0, -2.0) circle (2pt) node[below] {FM central};
\end{scope}
% Homo erectus
\begin{scope}[xshift=14cm]
\draw[popgreen, thick] (0,0) ellipse (2.2 and 2.5);
\draw[popgreen, thick] (-2.2,0) -- (-3.0, 0.2) -- (-3.0, 1.0) -- (-2.2, 1.0) -- cycle;
\draw[popgreen, thick] (-2.2, -1.5) arc (180:360:2.2 and 0.5);
\node[popgreen] at (-1.5, 3.1) {Doc 2 : \emph{H. erectus}};
\fill[popred] (0.2, -2.3) circle (2pt) node[below] {FM central/antérieur};
\end{scope}
% Homo sapiens
\begin{scope}[xshift=21cm]
\draw[poppurple, thick] (0,0) ellipse (2.5 and 2.8);
\draw[poppurple, thick] (-2.5,0) -- (-2.7, 0.5) -- (-2.7, 0.8) -- (-2.5, 0.8) -- cycle; % flat face
\draw[poppurple, thick] (-2.5, -1.5) arc (180:360:2.5 and 0.5);
\node[poppurple] at (-1.5, 3.4) {Doc 3 : \emph{H. sapiens}};
\fill[popred] (0.5, -2.6) circle (2pt) node[below] {FM central/antérieur};
\end{scope}
\end{tikzpicture}
\legende{Schémas comparatifs de crânes en vue latérale (non à l'échelle). FM = Foramen Magnum (trou occipital). Point rouge : position du FM.}
\end{popfigure}

\begin{center}
\begin{tabularx}{\linewidth}{|l|>{\centering\arraybackslash}X|>{\centering\arraybackslash}X|>{\centering\arraybackslash}X|>{\centering\arraybackslash}X|}\hline
\textbf{Critère mesuré} & \textbf{Doc 0 : Chimpanzé} & \textbf{Doc 1 : \emph{A. afarensis}} & \textbf{Doc 2 : \emph{H. erectus}} & \textbf{Doc 3 : \emph{H. sapiens}} \\ \hline
Volume crânien (\si{cm^3}) & 400 & 450 & 1000 & 1350 \\ \hline
Indice d'orthognathie (face plate) & 30 (prognathe) & 45 & 65 & 95 \\ \hline
Position trou occipital (distance du bord antérieur en \%) & 85\% (postérieur) & 50\% (central) & 40\% (central-antérieur) & 35\% (central-antérieur) \\ \hline
Torse supra-orbitaire & Très développé & Développé & Très développé (barre continue) & Faible / absent \\ \hline
Crête sagittale & Présente & Parfois (mâles) & Absente & Absente \\ \hline
\end{tabularx}
\end{center}

\begin{enumerate}
    \item À partir du tableau, identifiez \textbf{trois critères morphologiques} de l'hominisation qui évoluent de manière progressive du Doc 0 au Doc 3.
    \item Pour chaque critère, précisez le sens de l'évolution (augmentation, diminution, déplacement) et son avantage sélectif probable.
    \item La bipédie est-elle acquise dès \emph{Australopithecus} ? Justifiez par la position du trou occipital.
    \item L'augmentation du volume crânien est-elle synchronisée avec l'acquisition de la bipédie ? Concluez sur le caractère « mosaïque » de l'hominisation.
\end{enumerate}
\end{exercice}

\begin{exercice}[Construction d'un arbre phylogénétique : Données moléculaires (Hémoglobine $\beta$)]
Le tableau ci-dessous donne le nombre de différences d'acides aminés sur la chaîne $\beta$ de l'hémoglobine (146 acides aminés au total) entre cinq espèces de primates. \textbf{Note :} Les valeurs ont été ajustées pour refléter la proximité Homme-Chimpanzé.

\begin{center}
\begin{tabular}{|c|ccccc|}\hline
\textbf{Espèces comparées} & \textbf{Homme} & \textbf{Chimpanzé} & \textbf{Gorille} & \textbf{Orang-outan} & \textbf{Gibbon} \\ \hline
Homme & 0 & 1 & 2 & 16 & 25 \\
Chimpanzé & 1 & 0 & 2 & 17 & 26 \\
Gorille & 2 & 2 & 0 & 17 & 26 \\
Orang-outan & 16 & 17 & 17 & 0 & 15 \\
Gibbon & 25 & 26 & 26 & 15 & 0 \\ \hline
\end{tabular}
\end{center}

\begin{enumerate}
    \item Calculez le pourcentage d'identité de séquence entre l'Homme et le Chimpanzé, et entre l'Homme et l'Orang-outan.
    \item En vous basant sur le principe de l'horloge moléculaire (le nombre de différences est proportionnel au temps de divergence), construisez l'arbre phylogénétique de ces cinq espèces. Représentez-le sous forme de cladogramme (branches non proportionnelles au temps) en indiquant les nœuds de divergence.
    \item Quelle espèce est le groupe frère de l'Homme ? Justifiez par les données du tableau.
    \item L'Orang-outan et le Gibbon sont-ils plus proches parents entre eux qu'avec l'Homme ? Justifiez.
    \item \emph{Question de cours :} Pourquoi l'utilisation de l'ADN mitochondrial (ADNmt) ou de l'ADN nucléaire non codant est-elle souvent privilégiée pour dater les divergences récentes ($<$ \SI{5}{Ma}) plutôt que les séquences protéiques codantes ?
\end{enumerate}
\end{exercice}

\begin{exercice}[Preuve caryotypique : La fusion du chromosome 2 humain]
L'Homme possède un caryotype à \textbf{46 chromosomes} (2n=46), tandis que le Chimpanzé, le Gorille et l'Orang-outan en possèdent \textbf{48} (2n=48). L'analyse par bandage G (bandes G) révèle une correspondance précise entre le chromosome 2 humain et deux chromosomes acrocentriques distincts chez le chimpanzé (souvent notés chr 2a et 2b).

\begin{popfigure}
\begin{tikzpicture}[scale=0.9, every node/.style={font=\scriptsize}]
% Human Chr 2
\draw[popblue, thick, fill=popblueL] (0,0) rectangle (1, 5);
\foreach \y/\c in {0.5/popdark, 1.5/popmuted, 2.5/popdark, 3.5/popmuted, 4.5/popdark}
  \draw[\c, thick] (0,\y) -- (1,\y);
\node[popblue] at (0.5, 5.4) {Homme : Chr 2 (Métacentrique)};
\node[popblue] at (0.5, -0.4) {Centromère unique};

% Chimp Chr 2a & 2b
\begin{scope}[xshift=3cm]
\draw[poporange, thick, fill=poporangeL] (0,0) rectangle (0.6, 4.5); % 2a
\foreach \y/\c in {0.5/popdark, 1.5/popmuted, 2.5/popdark, 3.5/popmuted, 4.0/popdark}
  \draw[\c, thick] (0,\y) -- (0.6,\y);
\node[poporange] at (0.3, 4.9) {Chimpanzé : Chr 2a (Acrocentrique)};
\node[poporange] at (0.3, -0.4) {Centromère terminal};

\draw[poporange, thick, fill=poporangeL] (1.2,0) rectangle (1.8, 3.5); % 2b
\foreach \y/\c in {0.5/popdark, 1.5/popmuted, 2.5/popdark, 3.0/popdark}
  \draw[\c, thick] (1.2,\y) -- (1.8,\y);
\node[poporange] at (1.5, 3.9) {Chr 2b (Acrocentrique)};
\node[poporange] at (1.5, -0.4) {Centromère terminal};
\end{scope}

% Arrows showing fusion
\draw[->, thick, popgreen] (1.2, 2.5) -- (2.8, 2.5) node[midway, above] {Fusion centrique};
\draw[->, thick, popgreen] (1.2, 1.5) -- (2.8, 1.5);
\end{tikzpicture}
\legende{Schéma simplifié de la fusion de Robertson (centrique) à l'origine du chromosome 2 humain. Les bandes de bandage G (alternance clair/foncé) sont conservées, prouvant l'homologie de séquence.}
\end{popfigure}

\begin{enumerate}
    \item Expliquez le mécanisme cytogénétique (type de mutation chromosomique) qui a réduit le nombre de chromosomes de 48 à 46 chez l'ancêtre de la lignée humaine.
    \item En quoi la conservation de l'ordre des bandes (syntenie) entre le chromosome 2 humain et les deux chromosomes chimpanzés constitue-t-elle une preuve forte de parenté commune ?
    \item Cette fusion est-elle une synapomorphie (caractère dérivé partagé) de l'espèce \emph{Homo sapiens} seulement, ou de la lignée humaine (homininés) dans son ensemble ? Justifiez en citant le caryotype de l'Homme de Néandertal ou de Denisova (connus par paléogénomique).
    \item Pourquoi cette fusion n'a-t-elle pas entraîné de stérilité immédiate chez les hétérozygotes (porteurs d'un chr 2 fusionné et de deux chr ancestraux) ? (Indice : méiose et ségrégation).
\end{enumerate}
\end{exercice}

\begin{exercice}[Exploitation de courbe : Évolution du volume crânien au cours du temps]
La courbe ci-dessous représente l'évolution de la capacité crânienne moyenne (en \si{cm^3}) en fonction du temps (en millions d'années, Ma) pour les principaux représentants de la lignée humaine.

\begin{popfigure}
\begin{tikzpicture}
\begin{axis}[
    width=\linewidth, height=8cm,
    xlabel={Temps (Ma avant présent)}, ylabel={Volume crânien moyen (\si{cm^3})},
    xmin=7, xmax=0, % Past (left) to Present (right)
    ymin=300, ymax=1500,
    xtick={0,1,2,3,4,5,6,7}, ytick={400,600,800,1000,1200,1400},
    grid=major, grid style={popline},
    axis lines=left,
    enlargelimits=0.05,
    title style={font=\small}, title={Évolution du volume crânien chez les Homininés},
    x tick label style={/pgf/number format/fixed, /pgf/number format/precision=0}
]
% Data points (Time Ma, Volume cc)
% Sahelanthropus ~7 Ma, 350cc
\addplot[only marks, mark=*, mark size=3, poppurple] coordinates {(7, 360)} node[above left] {Toumaï};
% Orrorin ~6 Ma, ~350
\addplot[only marks, mark=*, mark size=3, poppurple] coordinates {(6, 350)};
% Ardipithecus ~4.4 Ma, 350
\addplot[only marks, mark=*, mark size=3, popblue] coordinates {(4.4, 350)} node[above] {Ardi};
% Australopithecus afarensis ~3.2 Ma, 450
\addplot[only marks, mark=square*, mark size=3, popblue] coordinates {(3.2, 450)} node[above] {Lucy};
% Australopithecus africanus ~2.5 Ma, 480
\addplot[only marks, mark=square*, mark size=3, popblue] coordinates {(2.5, 480)};
% Paranthropus boisei ~1.8 Ma, 520 (robust)
\addplot[only marks, mark=triangle*, mark size=3, poporange] coordinates {(1.8, 520)} node[above right] {Paranthropus};
% Homo habilis ~2.0-1.6 Ma, 650
\addplot[only marks, mark=*, mark size=3, popgreen] coordinates {(1.8, 650)} node[above] {H. habilis};
% Homo erectus/ergaster ~1.6-0.8 Ma, 900-1100
\addplot[only marks, mark=*, mark size=3, popgreen] coordinates {(1.6, 900) (1.0, 1050) (0.8, 1100)} node[above] {H. erectus};
% Homo heidelbergensis ~0.6-0.3 Ma, 1200
\addplot[only marks, mark=*, mark size=3, popdark] coordinates {(0.5, 1200)} node[above] {H. heidelbergensis};
% Homo neanderthalensis ~0.4-0.04 Ma, 1400-1600
\addplot[only marks, mark=diamond*, mark size=3, poppink] coordinates {(0.1, 1500)} node[above] {Néandertal};
% Homo sapiens ~0.3-0 Ma, 1350
\addplot[only marks, mark=*, mark size=3, popred] coordinates {(0.05, 1350)} node[above right] {H. sapiens};

% Trend lines (schematic)
\addplot[popblue, dashed, domain=4:2, samples=2] {350 + 50*(4-x)}; % Australopiths slight rise
\addplot[popgreen, dashed, domain=2:0.5, samples=10] {500 + 400*(2-x)/1.5}; % Homo rapid rise
\addplot[popdark, dashed, domain=0.5:0, samples=2] {1200 + 300*(0.5-x)/0.5}; % Late Homo
\end{axis}
\end{tikzpicture}
\legende{Évolution du volume crânien moyen (estimations paléoanthropologiques). Carrés : Australopithèques ; Triangles : Paranthropes ; Cercles : Homo ; Losange : Néandertal. Les lignes pointillées sont des tendances schématiques. L'axe des temps va du passé (7 Ma, gauche) au présent (0 Ma, droite).}
\end{popfigure}

\begin{enumerate}
    \item Relevez les volumes crâniens moyens approximatifs pour : \emph{Australopithecus afarensis}, \emph{Homo habilis}, \emph{Homo erectus}, \emph{Homo neanderthalensis} et \emph{Homo sapiens}.
    \item La courbe montre-t-elle une augmentation linéaire et continue du volume crânien depuis \SI{7}{Ma} ? Décrivez les phases (stagnation, rupture, accélération).
    \item L'espèce \emph{Paranthropus boisei} (branche latérale robuste) a-t-elle suivi la même tendance d'encéphalisation que le genre \emph{Homo} ? Quelle hypothèse fonctionnelle explique cette différence (régime alimentaire) ?
    \item \emph{Homo neanderthalensis} présente un volume crânien moyen légèrement supérieur à celui de \emph{Homo sapiens} actuel. Cela signifie-t-il qu'il était « plus intelligent » ? Quels autres critères anatomiques (organisation cérébrale, culture) distinguent les deux espèces ?
    \item Reliez l'accélération de l'encéphalisation chez \emph{Homo} (à partir de \SI{2}{Ma}) aux innovations culturelles majeures (outils Oldowayen, Acheuléen, maîtrise du feu, langage).
\end{enumerate}
\end{exercice}

\courssub{Je me prépare au Bac}

\begin{exercice}[Sujet type Bac : L'hominisation, un processus buissonnant et mosaïque (Documents à exploiter)]
\textbf{Document 1 :} Photographies de crânes en vue de face et latérale (non reproduites ici, se reporter aux schémas de l'exercice 5) : Crâne A (\emph{Sahelanthropus tchadensis}, \SI{7}{Ma}, Tchad), Crâne B (\emph{Australopithecus afarensis}, \SI{3,2}{Ma}, Éthiopie), Crâne C (\emph{Homo erectus}, \SI{1,6}{Ma}, Kenya), Crâne D (\emph{Homo sapiens}, \SI{0,03}{Ma}, France).

\textbf{Document 2 :} Tableau récapitulatif de caractères anatomiques.
\begin{center}
\begin{tabular}{|l|c|c|c|c|}\hline
\textbf{Caractère} & \textbf{Crâne A} & \textbf{Crâne B} & \textbf{Crâne C} & \textbf{Crâne D} \\ \hline
Volume crânien (\si{cm^3}) & 360 & 450 & 950 & 1350 \\ \hline
Position trou occipital & Postérieur & Central & Antérieur & Antérieur \\ \hline
Prognathisme facial & Très fort & Fort & Modéré & Faible (orthognathe) \\ \hline
Torse supra-orbitaire & Massif & Développé & Barre continue & Absent \\ \hline
Dentition (canines, diastème) & Canines moyennes, diastème & Canines réduites, pas de diastème & Canines petites & Canines très petites \\ \hline
Datation absolue & \SI{7}{Ma} & \SI{3,2}{Ma} & \SI{1,6}{Ma} & \SI{0,03}{Ma} \\ \hline
\end{tabular}
\end{center}

\textbf{Document 3 :} Extrait de carte paléogéographique. Les sites de découverte des crânes A, B, C sont localisés en Afrique (Tchad, Éthiopie, Kenya). Le crâne D est découvert en Europe, mais les plus anciens \emph{Homo sapiens} (\SI{300000}{ans}) sont africains (Jebel Irhoud, Maroc ; Omo Kibish, Éthiopie).

\begin{enumerate}
    \item \textbf{Classement chronologique :} Rangez les quatre crânes par ordre chronologique croissant en vous basant sur le Document 2.
    \item \textbf{Critères de l'hominisation :} En exploitant le Document 2, dégagez \textbf{trois critères morphologiques} qui évoluent au cours de l'hominisation. Pour chacun, précisez le sens de l'évolution (ex : augmentation, diminution, déplacement antérieur).
    \item \textbf{Processus mosaïque :} Montrez, à l'aide d'exemples précis pris dans le tableau, que l'acquisition des caractères de l'hominisation n'est pas synchronisée (ex : bipédie avant encéphalisation).
    \item \textbf{Origine africaine et dispersion :} En vous appuyant sur le Document 3 et vos connaissances, expliquez le scénario « Out of Africa » (Sortie d'Afrique) pour le peuplement de la planète par \emph{Homo sapiens}. Citez deux vagues de dispersion majeures.
    \item \textbf{Coexistence :} Le Document 2 montre-t-il une coexistence d'espèces ? En utilisant vos connaissances, citez un exemple de coexistence entre deux espèces d'homininés distinctes (genre \emph{Homo} ou \emph{Paranthropus}) à la même époque et sur le même continent.
\end{enumerate}
\end{exercice}

\begin{exercice}[Sujet type Bac : Preuves moléculaires et chromosomiques de la parenté Homme / Grands Singes]
\textbf{Document 1 :} Distances génétiques (nombre de substitutions nucléotidiques pour 100 sites) pour un gène mitochondrial (COI) et un gène nucléaire (intron) entre cinq primates.
\begin{center}
\begin{tabular}{|l|ccccc|}\hline
& \textbf{Homme} & \textbf{Chimpanzé} & \textbf{Gorille} & \textbf{Orang-outan} & \textbf{Gibbon} \\ \hline
Homme & 0 & 1,2 & 1,6 & 3,1 & 5,8 \\
Chimpanzé & 1,2 & 0 & 1,7 & 3,2 & 5,9 \\
Gorille & 1,6 & 1,7 & 0 & 3,0 & 5,7 \\
Orang-outan & 3,1 & 3,2 & 3,0 & 0 & 4,5 \\
Gibbon & 5,8 & 5,9 & 5,7 & 4,5 & 0 \\ \hline
\end{tabular}
\end{center}

\textbf{Document 2 :} Schéma de bandage G (représentation simplifiée) du chromosome 2 humain et des chromosomes homologues 2a et 2b du chimpanzé. (Voir figure de l'exercice 7). On observe une correspondance parfaite des bandes entre le bras long du chr 2 humain et le chr 2a chimpanzé, et entre le bras court du chr 2 humain et le chr 2b chimpanzé. Le centromère du chr 2 humain correspond à l'un des centromères ancestraux (celui du 2a), l'autre (celui du 2b) étant inactivé (vestige centromérique détectable par séquençage : séquences alpha-satellites).

\begin{enumerate}
    \item \textbf{Arbre phylogénétique moléculaire :} À partir du Document 1, construisez l'arbre phylogénétique (cladogramme) des cinq espèces. Indiquez les nœuds et les groupes frères. Justifiez la topologie de l'arbre par les valeurs du tableau.
    \item \textbf{Datation relative :} En admettant une horloge moléculaire constante pour ce gène, quel couple d'espèces a divergé le plus récemment ? Le plus anciennement ?
    \item \textbf{Preuve caryotypique (Doc 2) :} Expliquez comment l'observation du bandage G et la présence de séquences centromériques vestigiales (alpha-satellites) au milieu du bras long du chromosome 2 humain permettent de confirmer l'hypothèse d'une fusion centrique (Robertsonienne) chez l'ancêtre commun de la lignée humaine.
    \item \textbf{Synthèse :} En combinant les arguments moléculaires (Doc 1) et chromosomiques (Doc 2), rédigez une conclusion argumentée (10-15 lignes) sur la proximité phylogénétique entre l'Homme et le Chimpanzé, et sur le moment de la divergence par rapport au Gorille.
\end{enumerate}
\end{exercice}

\begin{exercice}[Sujet type Bac : Adaptation de l'Homme, responsabilité environnementale et Développement Durable]
\textbf{Document 1 :} L'Homme (\emph{Homo sapiens}) est le seul homininé actuel. Son évolution a été marquée par l'acquisition d'une \cle{bipédie} libérant les mains, d'une \cle{main préhensile} avec pouce opposant long, d'une \cle{encéphalisation} majeure (volume crânien $\approx$ \SI{1350}{cm^3}), et d'une \cle{culture cumulative} (langage symbolique, techniques complexes, transmission intergénérationnelle). Ces adaptations biologiques et culturelles lui ont permis de coloniser tous les biomes de la planète (du désert du Sahara aux forêts du bassin du Congo, aux hautes altitudes de l'Himalaya).

\textbf{Document 2 :} Au Cameroun, la pression anthropique sur les écosystèmes s'intensifie : déforestation pour l'agriculture extensive (cacao, huile de palme) et l'exploitation forestière dans le Sud (forêt dense humide du bassin du Congo) ; surexploitation des ressources halieutiques sur la côte ; érosion des sols et désertification dans le Nord (région de l'Extrême-Nord, pourtour du Lac Tchad) ; pollution urbaine (Yaoundé, Douala) et gestion des déchets plastiques ; risque volcanique (Mont Cameroun) et limnique (Lac Nyos, Lac Monoun) exacerbés par l'installation humaine à proximité.

\textbf{Document 3 :} Définition du Développement Durable (Rapport Brundtland, 1987) : « Un développement qui répond aux besoins du présent sans compromettre la capacité des générations futures à répondre aux leurs. » Il repose sur trois piliers indissociables : \textbf{Environnement} (préservation de la biodiversité, climat, ressources), \textbf{Social} (équité, santé, éducation), \textbf{Économique} (efficacité, lutte contre la pauvreté).

\begin{enumerate}
    \item \textbf{Capacité d'adaptation biologique vs culturelle :} En vous appuyant sur le Document 1, distinguez les adaptations \emph{biologiques} (génétiques, lentes, sélection naturelle) des adaptations \emph{culturelles} (acquises, rapides, transmission horizontale/verticale) qui ont permis à \emph{Homo sapiens} de s'adapter à des environnements extrêmes sans modification génétique majeure (ex : vêtement, habitat, agriculture, médecine). Donnez un exemple camerounais d'adaptation culturelle au climat (Nord vs Sud).
    \item \textbf{Rupture du rapport à l'environnement :} Expliquez en quoi la révolution néolithique (domestication plantes/animaux, sédentarisation) puis la révolution industrielle ont modifié le rapport de l'Homme à son environnement, le faisant passer de « prédateur-cueilleur » intégré aux écosystèmes à « ingénieur des écosystèmes » capable de les transformer massivement (niche écologique élargie).
    \item \textbf{Responsabilité éthique et Développement Durable :} En exploitant les Documents 2 et 3, montrez que la capacité d'adaptation exceptionnelle de l'Homme (Document 1) lui confère aujourd'hui une \textbf{responsabilité unique} vis-à-vis de la biosphère. Illustrer votre propos par \textbf{deux exemples concrets au Cameroun} où l'absence de gestion durable (piliers Environnement/Économique/Social non respectés) menace la survie même des populations humaines à long terme (ex : Lac Nyos, déforestation, érosion côtière).
    \item \textbf{Synthèse argumentée (Rédaction) :} Rédigez un texte argumenté de \textbf{20 à 25 lignes} répondant à la question suivante :
    \emph{« En quoi la connaissance de l'évolution de l'Homme (preuves, mécanismes, histoire) éclaire-t-elle notre responsabilité actuelle dans la mise en œuvre du développement durable ? »}
    Votre texte devra structurer le raisonnement en : 1) Rappel des preuves d'une histoire commune avec le vivant (parenté), 2) Spécificité de l'adaptation humaine (culture/technique), 3) Conséquences sur la biosphère (Anthropocène), 4) Impératif éthique et pratique du développement durable (exemple local).
\end{enumerate}
\end{exercice}

\newpage

\coursec{Corrigés des exercices}

\begin{corrige}[Exercice 1 — QCM : Preuves et critères de l'hominisation]
\begin{enumerate}
    \item \textbf{Bonne réponse : b).} Les preuves de l'évolution humaine sont pluridisciplinaires : fossiles (paléontologie), comparaison anatomique des squelettes, comparaison des molécules (séquences d'ADN, protéines) et datations absolues (radiochronologie). Les réponses a), c) et d) sont fausses : les plus anciens fossiles d'homininés sont africains, les mythes ne sont pas des preuves scientifiques, et les singes actuels ne sont pas nos ancêtres mais nos cousins évolutifs.
    \item \textbf{Bonne réponse : a).} Chez un bipède obligatoire, le trou occipital (foramen magnum) est en position centrale sous la boîte crânienne : la tête est équilibrée au sommet d'une colonne vertébrale verticale. Chez le quadrupède, il est en position postérieure. Le volume crânien (c) n'est pas un critère de bipédie ; le torse supra-orbitaire (d) est un caractère archaïque lié à la puissance masticatrice.
    \item \textbf{Bonne réponse : b).} \emph{A. afarensis} (Lucy) : volume crânien $\approx$ \SI{450}{cm^3} (proche de celui du chimpanzé), bassin et fémur de bipède obligatoire, mais face prognathe et crâne de type simien. C'est l'exemple type de l'évolution en \cle{mosaïque} : la bipédie est acquise bien avant l'encéphalisation.
    \item \textbf{Bonne réponse : b).} L'hominisation est un processus \cle{buissonnant} : plusieurs espèces d'homininés ont coexisté (ex. \emph{Paranthropus boisei} et \emph{Homo habilis} vers \SI{1,8}{Ma} en Afrique de l'Est). La bipédie a \emph{précédé} l'augmentation du volume crânien (d est faux) et le berceau de l'hominisation est l'Afrique (c est faux).
    \item \textbf{Bonne réponse : b).} La fusion centrique (translocation robertsonienne) de deux chromosomes acrocentriques ancestraux a produit le chromosome 2 humain ; la conservation des bandes de coloration et les vestiges de centromère et de télomères internes prouvent la descendance commune. Ce n'est ni une polyploïdie (a) ni une erreur sans conséquence évolutive (c).
\end{enumerate}
\end{corrige}

\begin{attention}[Points de vigilance]
Ne pas confondre \cle{bipédie facultative} (chimpanzé, occasionnelle) et \cle{bipédie obligatoire} (lignée humaine, exclusive). Ne jamais écrire que « l'Homme descend du singe » : l'Homme et le chimpanzé partagent un \emph{ancêtre commun} disparu.
\end{attention}

\begin{corrige}[Exercice 2 — Vrai ou Faux justifié]
\begin{enumerate}
    \item \textbf{VRAI.} \emph{Sahelanthropus tchadensis} (Toumaï, Toros-Menalla, Tchad, $\approx$ \SI{7}{Ma}) présente un trou occipital en position avancée suggérant une bipédie, ainsi que des canines réduites ; son âge est proche de la date de divergence estimée entre les lignées humaine et chimpanzé (6 à \SI{8}{Ma}).
    \item \textbf{FAUX.} L'augmentation du volume crânien n'est ni régulière ni continue : stagnation autour de 350 à \SI{500}{cm^3} pendant près de \SI{5}{Ma} (de Toumaï aux australopithèques), puis accélération rapide chez le genre \emph{Homo} à partir d'environ \SI{2}{Ma}.
    \item \textbf{VRAI.} \emph{H. sapiens} et \emph{H. neanderthalensis} ont cohabité en Europe et au Proche-Orient entre environ \num{45000} et \num{40000} ans ; la paléogénomique montre des hybridations : les non-Africains actuels possèdent environ 1 à \SI{4}{\%} d'ADN d'origine néandertalienne.
    \item \textbf{FAUX.} Les données moléculaires (ADN mitochondrial et ADN nucléaire) placent le \emph{chimpanzé} comme groupe frère de l'Homme (divergence $\approx$ 6 à \SI{8}{Ma}) ; le gorille a divergé plus tôt ($\approx$ 8 à \SI{10}{Ma}).
\end{enumerate}
\end{corrige}

\begin{corrige}[Exercice 3 — Définitions et mise en relation]
Les définitions données dans le tableau sont exactes ; il fallait vérifier l'adéquation concept/définition :
\begin{itemize}
    \item \textbf{Hominisation} : ensemble des transformations morphologiques (bipédie, encéphalisation, face orthognathe, main préhensile) et comportementales (culture) ayant conduit à la lignée humaine.
    \item \textbf{Bipédie obligatoire} : locomotion exclusive sur les membres inférieurs, corrélée à la réorganisation du bassin (court et large), du fémur (angle bicondylien marqué) et à la position centrale du trou occipital.
    \item \textbf{Arbre phylogénétique} : graphe ramifié représentant les relations de parenté entre espèces, construit à partir des caractères dérivés partagés (synapomorphies).
    \item \textbf{Développement durable} : développement répondant aux besoins du présent sans compromettre la capacité des générations futures à répondre aux leurs (trois piliers : environnemental, social, économique).
    \item \textbf{Fusion de Robertson} : fusion centrique de deux chromosomes acrocentriques en un chromosome métacentrique (passage de 48 à 46 chromosomes dans la lignée humaine).
\end{itemize}
\textbf{Flèches logiques attendues :}
\emph{Fusion de Robertson} $\rightarrow$ \emph{Preuve de parenté} $\rightarrow$ \emph{Arbre phylogénétique} ; \emph{Bipédie obligatoire} $\rightarrow$ \emph{Critère de l'hominisation} ; \emph{Hominisation} $\rightarrow$ (capacité d'adaptation) $\rightarrow$ \emph{Développement durable}.
\end{corrige}

\begin{attention}[Point de vigilance]
Une définition doit être générale et précise : éviter les définitions par l'exemple (« la bipédie, c'est marcher comme Lucy »). Chaque définition doit pouvoir s'appliquer à tous les cas et seulement à eux.
\end{attention}

\begin{corrige}[Exercice 4 — Ordre chronologique et datations absolues]
\textbf{Classement du plus ancien au plus récent :}
\begin{center}
\begin{tabularx}{\linewidth}{|c|X|c|X|}\hline
\textbf{Rang} & \textbf{Fossile} & \textbf{Datation} & \textbf{Genre / Espèce} \\ \hline
1 & Toumaï (Tchad) & $\sim$ \SI{7}{Ma} & \emph{Sahelanthropus tchadensis} \\ \hline
2 & Lucy (Éthiopie) & \SI{3,2}{Ma} & \emph{Australopithecus afarensis} \\ \hline
3 & Turkana Boy (Kenya) & \SI{1,6}{Ma} & \emph{Homo ergaster} (ou \emph{H. erectus}) \\ \hline
4 & Homme de Tautavel (France) & $\sim$ \SI{0,45}{Ma} & \emph{Homo heidelbergensis} \\ \hline
5 & Cro-Magnon (France) & $\sim$ \SI{0,03}{Ma} & \emph{Homo sapiens} \\ \hline
\end{tabularx}
\end{center}
\textbf{Question bonus :} c'est la méthode \cle{potassium-argon} (\ce{^{40}K}/\ce{^{40}Ar}) qui a permis de dater les tufs volcaniques encaissant Lucy. Justification : la période radioactive du \ce{^{14}C} est de \SI{5730}{ans}, ce qui limite son usage à environ \num{50000} ans ; à \SI{3,2}{Ma}, la totalité du \ce{^{14}C} initial est désintégrée. La méthode K/Ar (période du \ce{^{40}K} : \SI{1,25}{Ga}) convient aux roches volcaniques anciennes de plusieurs millions d'années.
\end{corrige}

\begin{attention}[Point de vigilance]
On ne date pas le fossile lui-même mais les couches volcaniques qui l'encadrent (datation indirecte, par encadrement). Citer toujours la période radioactive pour justifier le choix du chronomètre.
\end{attention}

\begin{corrige}[Exercice 5 — Analyse de documents : critères de l'hominisation sur crânes]
\begin{enumerate}
    \item \textbf{Trois critères progressifs} (Doc 0 $\rightarrow$ Doc 3) : le \cle{volume crânien} (400 $\rightarrow$ \SI{1350}{cm^3}), l'\cle{indice d'orthognathie} (30 $\rightarrow$ 95, la face s'aplatit sous le front), la \cle{position du trou occipital} (85 \% postérieur $\rightarrow$ 35 \% central-antérieur). On peut ajouter la réduction du torse supra-orbitaire et la disparition de la crête sagittale.
    \item \textbf{Sens de l'évolution et avantages sélectifs :}
    \begin{itemize}
        \item Volume crânien : \emph{augmentation} ($\times 3,4$) $\rightarrow$ développement du cerveau (cognition, langage, fabrication d'outils).
        \item Orthognathie : \emph{augmentation de l'indice} = réduction du prognathisme $\rightarrow$ réduction de l'appareil masticateur (régime alimentaire varié, puis cuisson des aliments).
        \item Trou occipital : \emph{déplacement antérieur} $\rightarrow$ équilibre de la tête au sommet d'une colonne vertébrale verticale, condition de la bipédie permanente (libération des mains, économie d'énergie locomotrice, vision au-dessus des herbes hautes).
    \end{itemize}
    \item \textbf{Oui}, la bipédie est acquise dès \emph{Australopithecus} : le trou occipital du Doc 1 est déjà en position \emph{centrale} (50 \%), contre 85 \% (postérieur) chez le chimpanzé quadrupède. La tête est donc déjà équilibrée au-dessus du rachis.
    \item \textbf{Non}, les deux acquisitions ne sont pas synchronisées : au Doc 1, la bipédie est acquise (trou occipital central) alors que le volume crânien reste proche de celui du chimpanzé (\SI{450}{cm^3} contre 400). L'encéphalisation majeure n'apparaît qu'avec le genre \emph{Homo} (Doc 2 : \SI{1000}{cm^3}). \textbf{Conclusion :} l'hominisation est un processus en \cle{mosaïque} : les caractères évoluent à des rythmes différents, la bipédie précédant l'encéphalisation.
\end{enumerate}
\end{corrige}

\begin{corrige}[Exercice 6 — Arbre phylogénétique : hémoglobine $\beta$]
\begin{enumerate}
    \item \textbf{Calculs d'identité} (chaîne de 146 acides aminés) :
    \begin{align*}
    \text{Homme--Chimpanzé : } & \frac{146-1}{146}\times 100 = \frac{145}{146}\times 100 \approx 99,3\,\% \\
    \text{Homme--Orang-outan : } & \frac{146-16}{146}\times 100 = \frac{130}{146}\times 100 \approx 89,0\,\%
    \end{align*}
    \item \textbf{Construction du cladogramme} (méthode d'agglomération par distances croissantes) : le couple le moins distant est Homme--Chimpanzé (1 différence) $\rightarrow$ premier nœud N1. Le Gorille les rejoint (2 différences avec l'Homme) $\rightarrow$ nœud N2. Puis l'Orang-outan (16--17 différences) $\rightarrow$ nœud N3, enfin le Gibbon (25--26 différences), le plus éloigné.
\begin{popfigure}
\begin{tikzpicture}[scale=0.85, every node/.style={font=\small}, grow=right, level distance=2.6cm,
  level 1/.style={sibling distance=3.2cm}, level 2/.style={sibling distance=1.6cm},
  level 3/.style={sibling distance=0.9cm}, edge from parent/.style={draw=popdark, thick}]
\node[popdark] {Ancêtre commun}
  child { node {Gibbon} }
  child { node[popdark] {N3}
    child { node {Orang-outan} }
    child { node[popdark] {N2}
      child { node {Gorille} }
      child { node[popdark] {N1}
        child { node {Chimpanzé} }
        child { node {Homme} }
      }
    }
  };
\end{tikzpicture}
\legende{Cladogramme des cinq primates d'après les distances sur l'hémoglobine $\beta$. N1 : divergence Homme/Chimpanzé ; N2 : divergence avec la lignée du Gorille ; N3 : divergence avec l'Orang-outan.}
\end{popfigure}
    \item Le \textbf{Chimpanzé} est le groupe frère de l'Homme : c'est l'espèce avec laquelle il présente le \emph{plus petit nombre de différences} (1 sur 146, soit 99,3 \% d'identité), donc celle avec laquelle il partage l'ancêtre commun le plus récent.
    \item \textbf{Oui} : l'Orang-outan et le Gibbon ne diffèrent que de 15 acides aminés entre eux, valeur inférieure à leurs distances respectives avec l'Homme (16--17 et 25--26). Ils partagent donc entre eux un ancêtre commun plus récent que celui qu'ils partagent avec l'Homme. Toutefois, ces distances restent élevées et proches les unes des autres : une seule protéine ne suffit pas à trancher finement l'ordre des divergences anciennes ; seules des données génomiques complètes confirment que le Gibbon a divergé avant l'Orang-outan.
    \item \textbf{ADN mitochondrial et ADN non codant :} les séquences \emph{non codantes} et l'ADN mitochondrial accumulent des mutations quasi neutres (non soumises à sélection) à un rythme rapide et régulier : elles offrent une \cle{horloge moléculaire} fine pour les divergences récentes. Les protéines codantes évoluent lentement (sélection purificatrice) et manquent de résolution pour des divergences de moins de \SI{5}{Ma}.
\end{enumerate}
\end{corrige}

\begin{corrige}[Exercice 7 — Preuve caryotypique : la fusion du chromosome 2]
\begin{enumerate}
    \item Il s'agit d'une \textbf{mutation chromosomique structurale} : une \cle{fusion centrique} (translocation robertsonienne) entre deux chromosomes acrocentriques ancestraux. Les deux chromosomes ont fusionné au niveau de leurs centromères, formant un unique chromosome métacentrique (le chromosome 2 humain), avec perte des petits bras et inactivation de l'un des deux centromères. Bilan : $48 \rightarrow 46$ chromosomes.
    \item La \textbf{synténie} (conservation de l'ordre des bandes G) montre que le contenu génétique du chromosome 2 humain correspond \emph{exactement}, bande par bande, à la juxtaposition des chromosomes 2a et 2b du chimpanzé. Une telle correspondance ne peut résulter du hasard : elle traduit une \textbf{homologie} héritée d'un ancêtre commun, la fusion étant survenue dans la seule lignée humaine après la divergence.
    \item C'est une \textbf{synapomorphie de la lignée humaine (homininés)} et non de la seule espèce \emph{H. sapiens} : la paléogénomique montre que l'Homme de Néandertal et l'Homme de Denisova possédaient déjà 46 chromosomes avec un chromosome 2 fusionné. La fusion est donc antérieure à leur ancêtre commun (au moins \SI{600000}{ans}).
    \item Chez un hétérozygote, la méiose reste possible grâce à la formation d'un \textbf{trivalent} (le chromosome fusionné s'apparie avec les deux chromosomes ancestraux). Une partie des ségrégations est équilibrée et produit des gamètes viables (porteurs soit de la fusion, soit des deux chromosomes ancestraux). La fertilité est réduite mais non nulle, ce qui a permis à la fusion de se fixer progressivement dans de petites populations d'homininés (dérive génétique).
\end{enumerate}
\end{corrige}

\begin{attention}[Point de vigilance]
Ne pas écrire que la fusion a « créé » l'espèce humaine : c'est un marqueur de la lignée, fixé après la divergence d'avec le chimpanzé. Distinguer la fusion robertsonienne (entre acrocentriques) d'une simple translocation réciproque.
\end{attention}

\begin{corrige}[Exercice 8 — Exploitation de courbe : évolution du volume crânien]
\begin{enumerate}
    \item \textbf{Lecture graphique} (valeurs approchées) : \emph{A. afarensis} $\approx$ \SI{450}{cm^3} ; \emph{H. habilis} $\approx$ \SI{650}{cm^3} ; \emph{H. erectus} $\approx$ 900 à \SI{1100}{cm^3} ; \emph{H. neanderthalensis} $\approx$ \SI{1500}{cm^3} ; \emph{H. sapiens} $\approx$ \SI{1350}{cm^3}.
    \item \textbf{Non}, l'augmentation n'est ni linéaire ni continue. Trois phases : \textbf{stagnation} de 7 à $\approx$ \SI{3}{Ma} (350--\SI{450}{cm^3}, de Toumaï aux australopithèques) ; \textbf{rupture et accélération} à partir de \SI{2}{Ma} avec le genre \emph{Homo} (650 $\rightarrow$ \SI{1100}{cm^3} en un million d'années) ; \textbf{poursuite de la hausse} puis stabilisation après \SI{0,5}{Ma} (1200--\SI{1500}{cm^3}).
    \item \textbf{Non} : \emph{Paranthropus boisei} reste à $\approx$ \SI{520}{cm^3}, sans encéphalisation comparable à celle de \emph{Homo}. Hypothèse fonctionnelle : son régime végétal dur et abrasif (graines, racines) a entraîné le développement d'un appareil masticateur puissant (molaires géantes, crête sagittale, mandibule massive) au détriment de l'expansion cérébrale — une spécialisation alimentaire qui l'a conduit à l'extinction.
    \item \textbf{Non}, un volume crânien supérieur ne signifie pas une intelligence supérieure : ce qui compte est l'\emph{organisation} du cerveau (développement du cortex préfrontal, aires du langage) et non la seule taille. \emph{H. sapiens} se distingue par un front vertical, un menton saillant, une culture symbolique développée (art pariétal, parures, sépultures ritualisées) et des techniques plus élaborées.
    \item L'accélération de l'encéphalisation à partir de \SI{2}{Ma} coïncide avec les grandes innovations culturelles : outils \textbf{oldowayens} (\SI{2,5}{Ma}, galets taillés), puis \textbf{acheuléens} (\SI{1,7}{Ma}, bifaces symétriques témoignant d'une pensée abstraite), \textbf{maîtrise du feu} ($\approx$ \SI{0,8}{Ma} ; la cuisson rend les aliments plus digestes et libère de l'énergie pour le cerveau, organe très coûteux), et émergence du \textbf{langage articulé}. Culture et cerveau évoluent en \emph{boucle de rétroaction positive} (co-évolution gènes-culture).
\end{enumerate}
\end{corrige}

\begin{corrige}[Exercice 9 — Sujet type Bac : hominisation buissonnante et mosaïque]
\textbf{Barème indicatif (sur 20) :} Q1 : 2 pts ; Q2 : 6 pts ; Q3 : 4 pts ; Q4 : 5 pts ; Q5 : 3 pts.
\begin{enumerate}
    \item \textbf{Classement chronologique croissant :} Crâne A (\emph{S. tchadensis}, \SI{7}{Ma}) $\rightarrow$ Crâne B (\emph{A. afarensis}, \SI{3,2}{Ma}) $\rightarrow$ Crâne C (\emph{H. erectus}, \SI{1,6}{Ma}) $\rightarrow$ Crâne D (\emph{H. sapiens}, \SI{0,03}{Ma}).
    \item \textbf{Trois critères évolutifs} (Doc 2) :
    \begin{itemize}
        \item Volume crânien : \emph{augmentation} (360 $\rightarrow$ 450 $\rightarrow$ 950 $\rightarrow$ \SI{1350}{cm^3}) ;
        \item Position du trou occipital : \emph{déplacement antérieur} (postérieur $\rightarrow$ central $\rightarrow$ antérieur) ;
        \item Prognathisme : \emph{diminution} (très fort $\rightarrow$ orthognathe) ; on peut ajouter la réduction du torse supra-orbitaire et des canines.
    \end{itemize}
    \item \textbf{Processus en mosaïque :} au crâne B (\emph{A. afarensis}), le trou occipital est déjà en position \emph{centrale} (bipédie acquise) alors que le volume crânien reste très faible (\SI{450}{cm^3}, proche du chimpanzé) et le prognathisme fort. La bipédie précède donc l'encéphalisation de plusieurs millions d'années : les caractères de l'hominisation ne sont pas acquis simultanément.
    \item \textbf{Scénario « Out of Africa » :} la lignée humaine naît en Afrique (crânes A, B, C tous africains). Première vague : \emph{Homo erectus/ergaster} quitte l'Afrique vers \SI{1,8}{Ma} et colonise l'Eurasie (Dmanissi, Java, Chine). Deuxième vague : \emph{Homo sapiens}, apparu en Afrique il y a $\approx$ \SI{300000}{ans} (Jebel Irhoud, Maroc ; Omo Kibish, Éthiopie), en sort vers \num{60000}--\num{70000} ans, atteint l'Asie, l'Australie ($\approx$ \num{50000} ans), l'Europe ($\approx$ \num{45000} ans, où il côtoie Néandertal) puis les Amériques.
    \item \textbf{Coexistence :} le Doc 2 seul ne montre pas de coexistence (dates distinctes), mais les connaissances en fournissent : \emph{Paranthropus boisei} et \emph{Homo habilis} coexistent vers \SI{1,8}{Ma} en Afrique de l'Est (Olduvaï, Turkana) ; \emph{H. sapiens} et \emph{H. neanderthalensis} coexistent en Europe entre \num{45000} et \num{40000} ans. L'hominisation est donc \cle{buissonnante} et non linéaire.
\end{enumerate}
\end{corrige}

\begin{attention}[Points de vigilance — Bac]
Exigences du correcteur : citer les \emph{valeurs chiffrées} du document pour justifier chaque critère ; employer « ancêtre commun » et non « descend du singe » ; distinguer clairement ce qui relève du document et ce qui relève des connaissances (Q5).
\end{attention}

\begin{corrige}[Exercice 10 — Sujet type Bac : preuves moléculaires et chromosomiques]
\textbf{Barème indicatif (sur 20) :} Q1 : 6 pts ; Q2 : 3 pts ; Q3 : 5 pts ; Q4 : 6 pts.
\begin{enumerate}
    \item \textbf{Arbre phylogénétique :} la plus petite distance est Homme--Chimpanzé (1,2) $\rightarrow$ nœud N1 (groupes frères). Le Gorille les rejoint (1,6--1,7) $\rightarrow$ nœud N2. Puis l'Orang-outan ($\approx$ 3,0--3,2) $\rightarrow$ nœud N3. Enfin le Gibbon (4,5--5,9), le plus éloigné. La topologie est justifiée par l'ordre croissant des distances : plus la distance génétique entre deux espèces est faible, plus leur ancêtre commun est récent. (Cladogramme de même topologie que celui de l'exercice 6.)
    \item \textbf{Datation relative :} divergence la plus \emph{récente} : couple Homme--Chimpanzé (distance 1,2, la plus faible). Divergence la plus \emph{ancienne} : celle séparant le Gibbon de l'ensemble (Homme + grands singes) (distances 5,7--5,9, les plus fortes).
    \item \textbf{Preuve caryotypique :} le bandage G montre que le bras long du chromosome 2 humain correspond bande par bande au chromosome 2a du chimpanzé, et son bras court au chromosome 2b. De plus, on détecte au milieu du bras long des séquences \emph{alpha-satellites} vestigiales (restes d'un centromère inactivé) et des séquences télomériques internes au point de jonction. Or un centromère et des télomères ne peuvent se retrouver au milieu d'un chromosome que par \emph{fusion centrique} de deux chromosomes acrocentriques suivie de l'inactivation de l'un des deux centromères. L'hypothèse d'une translocation robertsonienne chez un ancêtre de la lignée humaine (48 $\rightarrow$ 46) est donc confirmée par deux indices indépendants et concordants.
    \item \textbf{Synthèse (10--15 lignes) :} Les données moléculaires (Doc 1) placent l'Homme et le Chimpanzé comme groupes frères : leur distance génétique (1,2 substitution pour 100 sites) est la plus faible du tableau, signe d'un ancêtre commun récent ($\approx$ 6--\SI{8}{Ma}). Le Gorille, légèrement plus distant (1,6--1,7), a divergé avant la séparation Homme--Chimpanzé ($\approx$ 8--\SI{10}{Ma}). Les données chromosomiques (Doc 2) confirment cette parenté : la quasi-totalité du caryotype humain est homologue de celui du chimpanzé, et la seule différence majeure — le chromosome 2 — s'explique par une fusion robertsonienne survenue dans la lignée humaine \emph{après} la divergence, comme le prouvent la synténie des bandes G et les vestiges centromériques et télomériques internes. La concordance de deux ordres de preuves indépendants (séquences d'ADN et structure des chromosomes) établit solidement que l'Homme et le Chimpanzé partagent un ancêtre commun exclusif, plus récent que celui partagé avec le Gorille. L'Homme ne descend donc pas du chimpanzé actuel : les deux espèces sont des cousins évolutifs issus d'une même souche africaine.
\end{enumerate}
\end{corrige}

\begin{attention}[Points de vigilance — Bac]
Dans un tableau de distances, toujours raisonner sur la \emph{plus petite valeur} pour identifier les groupes frères ; ne pas conclure à une parenté directe (« descend de ») mais à un \emph{ancêtre commun} ; en synthèse, articuler explicitement les deux documents (convergence des preuves).
\end{attention}

\begin{corrige}[Exercice 11 — Sujet type Bac : adaptation, responsabilité et développement durable]
\textbf{Barème indicatif (sur 20) :} Q1 : 4 pts ; Q2 : 4 pts ; Q3 : 5 pts ; Q4 : 7 pts (dont 2 pts pour la structure et l'expression).
\begin{enumerate}
    \item \textbf{Adaptations biologiques} : modifications génétiques transmises par la reproduction, lentes (sélection naturelle sur des millénaires) : bipédie, main préhensile, encéphalisation. \textbf{Adaptations culturelles} : acquises, non génétiques, transmises par apprentissage (horizontalement et verticalement), rapides et réversibles : vêtements, habitats, outils, agriculture, médecine. Exemples camerounais : dans le \emph{Nord} (climat soudano-sahélien), cases en banco à toit de paille restant fraîches, greniers surélevés, cultures de mil résistant à la sécheresse ; dans le \emph{Sud} (climat équatorial humide), cases sur pilotis, toits pentus évacuant les pluies, culture du cacao et du manioc.
    \item \textbf{Rupture du rapport à l'environnement :} le prédateur-cueilleur prélevait dans les écosystèmes sans les transformer durablement (faible densité, nomadisme). La \emph{révolution néolithique} ($\approx$ \num{10000} ans) : domestication des plantes et des animaux, sédentarisation, défrichement $\rightarrow$ l'Homme devient producteur et transforme les paysages. La \emph{révolution industrielle} ($\approx$ \num{250} ans) : énergies fossiles, mécanisation, urbanisation $\rightarrow$ transformation massive et globale des écosystèmes (déforestation, pollution, changement climatique) : l'Homme devient « ingénieur de la biosphère », sa niche écologique s'étend à la planète entière.
    \item \textbf{Responsabilité unique :} étant la seule espèce capable de modifier la biosphère à grande échelle \emph{et} d'en prévoir les conséquences (grâce à la culture cumulative et à la science), l'Homme porte une responsabilité éthique inédite. Exemples camerounais où le non-respect des piliers menace les populations :
    \begin{itemize}
        \item \textbf{Déforestation du bassin du Congo} (Sud-Cameroun) : l'exploitation forestière et les plantations extensives détruisent la biodiversité, perturbent le régime des pluies (pilier Environnement) et appauvrissent à terme les communautés locales (piliers Social et Économique) ;
        \item \textbf{Désertification de l'Extrême-Nord} (bassin du Lac Tchad, qui a perdu $\approx$ 90 \% de sa superficie depuis les années 1960) : surpâturage, déboisement et sécheresses entraînent érosion, insécurité alimentaire et migrations forcées ;
        \item On peut aussi citer l'installation humaine autour des \textbf{lacs Nyos et Monoun} (risque limnique : l'éruption de 1986 au lac Nyos a tué $\approx$ \num{1700} personnes) — le dégazage artificiel actuel illustre au contraire une adaptation technique responsable.
    \end{itemize}
    \item \textbf{Synthèse argumentée (corrigé-type, 20--25 lignes) :}\\
    La paléontologie, l'anatomie comparée et la biologie moléculaire établissent que l'Homme partage une histoire commune avec l'ensemble du vivant : cousin du chimpanzé (environ 99 \% d'identité génétique, ancêtre commun il y a 6--\SI{8}{Ma}), il est un produit de l'évolution, inséré dans la biosphère et non au-dessus d'elle. Cette parenté fonde une première leçon d'humilité : détruire le vivant, c'est dégrader sa propre généalogie.\\
    Mais la lignée humaine s'est singularisée par une adaptation d'un type nouveau : à la lente adaptation génétique (bipédie, encéphalisation) s'est substituée l'adaptation \emph{culturelle}, cumulative et rapide — outils, feu, langage, agriculture, industrie. Cette plasticité a permis à \emph{Homo sapiens} de coloniser tous les biomes, de la savane africaine aux milieux polaires, sans modification génétique majeure.\\
    L'envers de cette réussite est l'Anthropocène : devenu une force géologique, l'Homme épuise les ressources, détruit les écosystèmes et dérègle le climat. Au Cameroun, la déforestation du bassin du Congo, la désertification du pourtour du Lac Tchad et la pollution urbaine menacent directement la santé et la survie des populations — preuve que l'espèce la plus adaptable peut, par excès de puissance, détruire les conditions de sa propre adaptation.\\
    La connaissance de notre évolution impose donc un impératif éthique : notre capacité unique de prévision et de coopération doit servir le \cle{développement durable}, conciliant les piliers environnemental, social et économique — reforestation communautaire, agriculture durable, gestion des déchets, énergies renouvelables. Protéger la biosphère, c'est appliquer à l'échelle planétaire la leçon de sept millions d'années d'évolution : une espèce ne survit durablement qu'en équilibre avec son environnement.
\end{enumerate}
\end{corrige}

\begin{attention}[Points de vigilance — Bac]
Pour la rédaction : annoncer un plan, articuler les quatre étapes demandées, mobiliser au moins un exemple local chiffré, et éviter le hors-sujet moralisateur sans appui scientifique. Distinguer rigoureusement adaptation \emph{biologique} (génétique, héréditaire) et adaptation \emph{culturelle} (acquise, transmise par apprentissage).
\end{attention}

\newpage

\coursec{Fiche bilan}

\courssub{Fiche Bilan : L'évolution de l'Homme et son adaptation}

\begin{popfigure}
\begin{tikzpicture}[
    mindmap,
    grow cyclic,
    every node/.style={concept, execute at begin node=\hskip0pt},
    concept color=popblue!80,
    text=white,
    level 1/.append style={level distance=4.2cm, sibling angle=72, font=\small\bfseries},
    level 2/.append style={level distance=3.0cm, sibling angle=45, font=\tiny, text width=2.2cm, align=center}
]
\node[concept, scale=1.1, align=center]{Évolution\\ de l'Homme} [clockwise from=90]
    child[concept color=poporange] { node[concept] {Preuves} [clockwise from=60]
        child { node[concept] {Fossiles \& Datation} }
        child { node[concept] {Anatomie comparée} }
        child { node[concept] {Biochimie \& Génétique} }
        child { node[concept] {Embryologie} }
    }
    child[concept color=popgreen] { node[concept, align=center]{Critères\\ d'hominisation} [clockwise from=10]
        child { node[concept] {Bipédie obligatoire} }
        child { node[concept] {Libération main/outil} }
        child { node[concept] {Encéphalisation} }
        child { node[concept] {Face orthognathe} }
    }
    child[concept color=poppurple] { node[concept, align=center]{Lignée\\ humaine} [clockwise from=-50]
        child { node[concept] {Toumaï (\SI{7}{Ma})} }
        child { node[concept] {Australopithèques} }
        child { node[concept] {Homo habilis} }
        child { node[concept] {Homo erectus} }
        child { node[concept] {Homo sapiens} }
    }
    child[concept color=poppink] { node[concept] {Phylogénie} [clockwise from=-130]
        child { node[concept] {Divergence Pan/Homo} }
        child { node[concept] {Buisson évolutif} }
        child { node[concept] {Proximité ADN (\SI{98}{\%})} }
    }
    child[concept color=popgold] { node[concept] {Adaptation \& DD} [clockwise from=-170]
        child { node[concept] {Adaptation culturelle} }
        child { node[concept] {Pression sélection} }
        child { node[concept] {Gestion ressources} }
        child { node[concept] {Éthique responsabilité} }
    };
\end{tikzpicture}
\legende{Carte mentale récapitulative de la leçon 16 : preuves, critères, lignée, phylogénie et enjeux d'adaptation durable.}
\end{popfigure}

\begin{bilan}

\textbf{1. Preuves de l'évolution de l'Homme}
\begin{itemize}
    \item \textbf{Paléontologie :} Fossiles datés (méthodes relatives : stratigraphie ; absolues : \ce{^{14}C} < \SI{50000}{ans}, \ce{^{40}K}/\ce{^{40}Ar} > \SI{100000}{ans}). Séquence stratigraphique continue.
    \item \textbf{Anatomie comparée :} Homologies structurelles (squelette, membres, crâne) entre Homme et grands singes ; caractères dérivés partagés (synapomorphies).
    \item \textbf{Biochimie \& Génétique :} Proximité séquentielle ADN Homme/Chimpanzé \approx \SI{98,7}{\%} ; horloge moléculaire date divergence \approx \SI{7}{Ma} ; cytochrome c, hémoglobines identiques.
    \item \textbf{Embryologie :} Stades embryonnaires communs (fentes branchiales, queue, préhension palmaire) attestant ancêtre commun.
\end{itemize}

\textbf{2. Critères de l'hominisation (synapomorphies des Homininés)}
\begin{center}
\begin{tabularx}{\linewidth}{|l|X|}\hline
\textbf{Critère} & \textbf{Caractères dérivés (Homininés)} \\ \hline
Bipédie obligatoire & Bassin court/large, fémur oblique (angle col-diaphysaire), pied arqué (grand orteil aligné), foramen magnum antérieur/central. \\ \hline
Libération de la main & Pouce long/opposable, phalanges courtes, préhension de précision (pince pouce-index) $\rightarrow$ fabrication outils. \\ \hline
Encéphalisation & Volume crânien croissant (\SI{400}{cm^3} $\rightarrow$ \SI{1350}{cm^3}), bossage frontal, réduction face, os pariétaux bombés. \\ \hline
Face orthognathe & Réduction prognathisme, arcade dentaire parabolique, canines réduites (non émoulues), absence diastème, molaires à cuspides basses (bunyodontes). \\ \hline
Comportement & Outils taillés (Oldowayen, Achuléen), maîtrise feu (\approx \SI{400000}{ans}), langage symbolique, art, rites funéraires. \\ \hline
\end{tabularx}
\end{center}

\textbf{3. Grandes étapes de la lignée humaine (Buisson évolutif)}
\begin{center}
\begin{tabularx}{\linewidth}{|c|c|X|X|}\hline
\textbf{Groupe / Espèce} & \textbf{Âge (Ma)} & \textbf{Caractères clés} & \textbf{Lieux emblématiques (Afrique)} \\ \hline
\emph{Sahelanthropus tchadensis} (Toumaï) & \approx 7 & Mosaïque : foramen magnum antérieur (bipédie ?), boîte crânienne petite (\SI{350}{cm^3}), face orthognathe, canines réduites. & Toros-Menalla (Tchad) \\ \hline
\emph{Orrorin tugenensis} & \approx 6 & Fémur suggérant bipédie, dents postérieures épaisses. & Tugen Hills (Kenya) \\ \hline
\emph{Ardipithecus ramidus} & \approx 4,4 & Bipédie facultative (pied opposant), canines petites, habitat forestier. & Aramis (Éthiopie) \\ \hline
\emph{Australopithecus afarensis} (Lucy) & 3,9--2,9 & Bipédie assurée (pied, bassin, Laetoli), encéphale \SI{450}{cm^3}, prognathisme fort, dimorphisme sexuel marqué. & Hadar, Laetoli (Éthiopie, Tanzanie) \\ \hline
\emph{Australopithecus africanus} & 3--2 & Proche \emph{afarensis}, face moins prognathe, dents postérieures larges. & Sterkfontein (Afrique du Sud) \\ \hline
\emph{Paranthropus} (\emph{robustus, boisei}) & 2,5--1 & Hyper-robustesse masticatrice (crête sagittale, molaires géantes), bipédie, encéphale \SI{500}{cm^3}. Voie évolutive éteinte. & Afrique Est/Sud \\ \hline
\emph{Homo habilis} & 2,4--1,6 & Premier \emph{Homo} : encéphale \SI{650}{cm^3}, face réduite, main habile $\rightarrow$ Oldowayen (galets taillés). & Olduvai (Tanzanie) \\ \hline
\emph{Homo erectus} / \emph{ergaster} & 1,8--0,1 & Encéphale \SI{900}{--}\SI{1100}{cm^3}, stature moderne, Achuléen (bifaces), maîtrise feu, sortie d'Afrique (Dmanissi, Java, Zhoukoudian). & Afrique, Eurasie \\ \hline
\emph{Homo neanderthalensis} & 0,4--0,04 & Encéphale \SI{1400}{cm^3} > sapiens, robustesse, Moustérien, rites funéraires, adaptation froid. Eurasie. & Europe, Proche-Orient \\ \hline
\emph{Homo sapiens} & 0,3--actuel & Encéphale \SI{1350}{cm^3}, front vertical, menton, technologie complexe (Aurignacien), art pariétal, langage syntaxique. Origine Afrique (\approx \SI{300}{ka}). & Monde entier \\ \hline
\end{tabularx}
\end{center}

\textbf{4. Parenté Homme -- Grands singes : Arbre phylogénétique}
\begin{itemize}
    \item \textbf{Groupe frère :} \emph{Pan} (Chimpanzé \emph{P. troglodytes}, Bonobo \emph{P. paniscus}). Divergence \approx \SI{7}{Ma} (horloge moléculaire calibrée sur fossiles).
    \item \textbf{Homininés (Hominidés au sens strict) :} Clade incluant \emph{Homo} et tous les taxons plus proches de l'Homme que de \emph{Pan} (Toumaï, Australopithèques, \emph{Paranthropus}, \emph{Homo}).
    \item \textbf{Structure :} Buisson (radiation adaptative) et non échelle linéaire. Coexistence fréquente (ex: \emph{H. erectus}, \emph{H. habilis}, \emph{Paranthropus} à \approx \SI{1,8}{Ma}).
    \item \textbf{Preuve génétique :} Caryotype \emph{Homo} 2n=46 vs \emph{Pan/Gorilla} 2n=48 $\rightarrow$ fusion chromosomique ancestrale (chromosome 2 humain = fusion 2a+2b ancestraux). Séquences ERV partagées, pseudogènes identiques (ex: \emph{GULO} vitamine C).
\end{itemize}

\textbf{5. Capacité d'adaptation, Environnement et Développement Durable}
\begin{itemize}
    \item \textbf{Double adaptation :} Biologique (lente, sélection naturelle, ex: pigmentation peau, tolérance lactose, résistance paludisme \emph{HbS}) + Culturelle (rapide, transmission horizontale/verticale : outils, feu, habitat, médecine, agriculture, industrie).
    \item \textbf{Retour d'action :} L'adaptation culturelle modifie l'environnement $\rightarrow$ nouvelles pressions de sélection (ex: densité population $\rightarrow$ maladies infectieuses ; pollution $\rightarrow$ cancers ; antibiotiques $\rightarrow$ résistances).
    \item \textbf{Développement Durable (DD) :} « Développement qui répond aux besoins du présent sans compromettre la capacité des générations futures » (Brundtland 1987). \textbf{3 piliers indissociables :} Écologique (biodiversité, climat), Social (équité, santé, éducation), Économique (efficacité, lutte pauvreté).
    \item \textbf{Contexte Cameroun :} Déforestation bassin du Congo (perte biodiversité, carbone), érosion côtière (Kribi, Limbé), éruption limnique Lac Nyos (1986, gestion risque), agriculture cacaoyère durable (agroforesterie), gestion déchets urbains (Yaoundé, Douala), transition énergétique.
    \item \textbf{Responsabilité éthique :} Espèce unique capable d'anticiper (prospective) et de choisir (éthique). Préservation services écosystémiques = condition survie \emph{H. sapiens}.
\end{itemize}

\textbf{Méthode express : Construire un argumentaire évolutif (Bac)}
\begin{enumerate}
    \item \textbf{Observation :} Donnée factuelle (fossile, séquence ADN, anatomie, datation).
    \item \textbf{Question / Problème :} Lien de parenté ? Chronologie ? Mode de vie ?
    \item \textbf{Hypothèse :} Relation phylogénétique, caractère ancestral/dérivé, pression de sélection.
    \item \textbf{Analyse documents / Expérience :} Comparaison caractères (tableau), arbre phylogénétique, courbe encéphalisation, datation.
    \item \textbf{Interprétation :} Identification synapomorphies, positionnement nœud, rejet convergence.
    \item \textbf{Conclusion :} Énoncé validé (ex: « Toumaï est un homininé basal car il partage le foramen magnum antérieur avec les homininés ultérieurs »).
\end{enumerate}

\end{bilan}

\begin{aretenir}[Checklist avant le Bac]
\begin{itemize}
    \item $\square$ Citer \textbf{4 preuves} majeures de l'évolution humaine (Paléontologie, Anatomie, Biochimie, Embryologie) avec un exemple précis chacune.
    \item $\square$ Définir \textbf{l'hominisation} et lister les \textbf{5 critères} officiels (Bipédie, Main, Encéphale, Face/Dents, Comportement) avec le caractère dérivé associé.
    \item $\square$ Classer chronologiquement les genres : \emph{Sahelanthropus} $\rightarrow$ \emph{Orrorin/Ardipithecus} $\rightarrow$ \emph{Australopithecus} $\rightarrow$ \emph{Paranthropus} / \emph{Homo}.
    \item $\square$ Caractériser \textbf{Toumaï} (âge, mosaicité, statut basal), \textbf{Lucy} (bipédie sûre, encéphale chimpanzé), \textbf{H. habilis} (premier \emph{Homo}, Oldowayen), \textbf{H. erectus} (sortie Afrique, feu, Achuléen), \textbf{H. sapiens} (Afrique \approx 300 ka, modernité anatomique/comportementale).
    \item $\square$ Dessiner et interpréter l'\textbf{arbre phylogénétique} Homme/Chimpanzé : divergence \approx 7 Ma, groupe frère \emph{Pan}, clade Homininés, structure buissonnante (coexistence).
    \item $\square$ Exploiter la \textbf{génétique} : % ADN (98,7\%), caryotype (fusion chr 2), horloge moléculaire, ERV/pseudogènes partagés.
    \item $\square$ Différencier \textbf{adaptation biologique} (génétique, lente, générations) et \textbf{adaptation culturelle} (apprise, rapide, cumulable) ; donner un exemple de chaque.
    \item $\square$ Expliquer le \textbf{retour d'action} : culture modifie environnement $\rightarrow$ nouvelles pressions de sélection (ex: agriculture $\rightarrow$ amylase, lactase ; urbanisation $\rightarrow$ maladies).
    \item $\square$ Définir le \textbf{Développement Durable} et ses \textbf{3 piliers} (Écologique, Social, Économique) ; citer une action concrète au Cameroun par pilier.
    \item $\square$ Éviter les \textbf{pièges} : « L'Homme descend du singe » (faux : cousin commun), « Évolution = progrès linéaire » (faux : buisson, adaptation locale), confusion datation relative/absolue, inversion cause/effet (bipédie $\rightarrow$ savane ou savane $\rightarrow$ bipédie ?).
    \item $\square$ Maîtriser la \textbf{rédaction scientifique} : Observations $\rightarrow$ Hypothèses $\rightarrow$ Arguments (données chiffrées, unités) $\rightarrow$ Conclusion. Vocabulaire précis (orthognathe, prognathe, synapomorphie, clade, taxons).
\end{itemize}
\end{aretenir}

\findecours

\end{document}
